\chapter{Metrized bundles and curvature currents}
\index{metrized vector bundle}\index{curvature current}
\label{chap:cld-metrized-bundles}

The preceding chapters constructed differential forms and currents from
tropical charts.  We now attach that calculus to line bundles.  A metric is
described locally by a weight, curvature is obtained by applying
\(\dprimeddprime\), and a change of frame disappears by Poincar\'e--Lelong.  Formal
metrics then bring the differential theory back to intersection theory on a
special fibre.

\section{Orientation}
\label{sec:ch17-orientation}

\subsection{What has to be built, and why}
\label{sec:ch17-why}

Up to this point \Cref{part:cld} has been a theory of \emph{functions}.
A superform lives on a polyhedron, a current is a functional on test forms,
and a psh function is a section of a sheaf.  Nothing yet connects that
apparatus to a line bundle, and nothing yet produces a number attached to a
projective variety.

The bridge is the same one as in complex geometry, and it is elementary.  A
metric \(\|\cdot\|\) on a line bundle \(L\) is not a function on the base;
but if \(e\) is a local frame, then \(\phi_e=-\log\|e\|\) \emph{is} a
function, and the ambiguity in \(\phi_e\) as \(e\) varies is exactly a
logarithm \(\log|g|\) of an invertible analytic function.  The operator
\(\dprimeddprime\) kills such logarithms
(\Cref{thm:cld-poincare-lelong} on the locus where \(g\) is invertible), so
the local currents \(\dprimeddprime[\phi_e]\) glue.  The result is a global
current attached to \(\overline L=(L,\|\cdot\|)\) alone.  That current is
the first Chern current.

Three things then have to be proved, and they organize the chapter.

\begin{enumerate}
\item\label{itm:ch17-plan:classes}
 \emph{The metrics have to be classified by how singular their curvature
 is.}  Smooth metrics give smooth forms; but---and this is the first
 surprise of the non-Archimedean theory---smooth metrics do not exist at
 all on bundles of rank \(\ge2\)
 (\Cref{prop:ch17-no-smooth-higher-rank}), and the metrics that come from
 formal models are \emph{not} smooth even in rank one.  The correct class is
 the parlinear one, and its curvature is a singular, atomic object.
\item\label{itm:ch17-plan:products}
 \emph{The curvatures have to be multiplied.}  This is
 \Cref{chap:cld-bedford-taylor} applied to metric weights, with one
 addition: parlinear metrics are not psh in general, only \emph{dpsh}, so
 the multilinear extension \Cref{thm:ch16-dpsh-product} is what is actually
 used, and the resulting mixed product is a signed measure in top degree.
\item\label{itm:ch17-plan:numbers}
 \emph{The top product has to be computed.}  For a metric coming from a
 formal model, the answer is an intersection number on the special fibre.
 This is the theorem that closes the circle with
 \Cref{chap:metrics}, and it is the reason the chapter exists.
\end{enumerate}

\subsection{The payoff, stated in advance}
\label{sec:ch17-payoff}

In \Cref{sec:metrics-cl-measure} the Chambert--Loir measure of a family of
model metrics was \emph{defined} by an intersection-theoretic formula on the
special fibre, and \Cref{thm:cl-model-properties} verified that the
definition is consistent and has the right total mass.  Nothing in
\Cref{chap:metrics} explained why that formula deserves the notation
\(c_1(\overline L_1)\wedge\cdots\wedge c_1(\overline L_n)\); the symbol was
a promissory note.

\Cref{thm:cld-recovers-cl} redeems it.  The wedge product on the left is the
genuine Bedford--Taylor product of genuine curvature currents, formed by the
local analytic calculus of
\Crefrange{chap:cld-forms}{chap:cld-bedford-taylor} with no reference to any
model; the measure on the right is Chambert--Loir's.  They are equal.  In
particular
\[
 \int_{X^{\An}}c_1(\overline L_1)\wedge\cdots\wedge c_1(\overline L_n)
 =(L_1\cdots L_n)_X,
\]
which we verify numerically against the computations of
\Cref{chap:metrics} in \Cref{sec:ch17-cl-checks}.

Two remarks about the logical shape of that theorem.  First, it is not a
tautology dressed as a theorem: the two sides are computed by genuinely
different mechanisms---one by residual intersection theory on a graded
reduction, the other by Chow-theoretic degrees on a special fibre---and the
comparison passes through three correction factors
(\Cref{def:ch17-defect}) that are invisible in the discretely valued
normal case and are \emph{not} invisible in general.  Second, the theorem
runs in the direction the book needs: it lets a global model-theoretic
measure be treated as a local differential object, so that locality,
restriction, and support statements proved in
\Cref{chap:cld-bedford-taylor} become available to
\Cref{part:bfj,part:bj}.

\subsection{Standing hypotheses}
\label{sec:ch17-hypotheses}

Throughout, \(K\) is a complete non-Archimedean field with absolute value
\(|\cdot|\), residue field \(\widetilde K\), and valuation ring
\(K^\circ\).  Unless the contrary is said:

\begin{itemize}
\item \(X\) is a topologically separated \(K\)-analytic space of pure
 dimension \(n\) \emph{without boundary}, so that the calculus of
 \Cref{chap:cld-forms,chap:cld-currents} applies and Green's formula
 (\Cref{cor:analytic-stokes-green}) carries no flux term.  Boundarylessness
 also guarantees that \(X\) is good, hence that vector bundles for the
 \(G\)-topology are vector bundles for the usual topology.
\item Sections that speak of \emph{strict} analytic spaces, of
 \(G_{\mathrm{str}}\)-coverings, or of strictly integral parlinear metrics
 assume in addition that \(|K^\times|\ne\{1\}\).  This restriction is
 flagged where it occurs; it is never used silently.
\item Currents, bidegrees, positivity cones, supports, and the weak topology
 are those of \Cref{chap:cld-currents}; psh means CLD-psh in the sense of
 \Cref{chap:cld-poincare-lelong}; products of curvature currents are the
 Bedford--Taylor products of \Cref{chap:cld-bedford-taylor}, extended to
 differences of psh potentials by \Cref{thm:ch16-dpsh-product}.
\item Signs and lattice normalizations are those of
 \Cref{chap:tropical-integration} and \Cref{app:normalizations}; in
 particular \eqref{eq:appJ-cld-c1} fixes
 \(c_1(L,\|\cdot\|)=\dprimeddprime\lambda_e\) with
 \(\lambda_e=-\log\|e\|\), and \eqref{eq:appJ-cld-pl} fixes
 \(\dprimeddprime[\log|f|]=\delta_{\divis(f)}\).
\item \emph{Positive} means nonnegative.  \emph{Parlinear} and \emph{PL} are
 synonyms, as in \Cref{chap:cld-bedford-taylor}; we prefer the former when
 the sheaf-theoretic definition of \Cref{chap:tropical-spaces} is in play
 and the latter when the reader is likely to think of \(\mathbf Q\)-model
 metrics.
\end{itemize}

One convention deserves a separate sentence because it is a frequent source
of confusion.  A \emph{semi-metric} allows the fibre norms to be
degenerate, and its weight takes values in \(\mathbf R\cup\{-\infty\}\).  We
allow semi-metrics only when they are locally integrable, and we then say so.
Every unqualified ``metric'' below is a genuine metric.

\subsection{What is proved here and what is imported}
\label{sec:ch17-imported}

Everything from \Cref{sec:ch17-metrics} to
\Cref{sec:ch17-positivity-detection} is proved, with the following
exceptions, each quoted with its hypotheses at the point of use.

\begin{enumerate}
\item The structure theory of the graded reduction of a normed vector
 space over \(\cH(x)\): existence of an orthogonal basis over a spherically
 complete algebraically closed field, and the resulting description of the
 graded residual space.  This is Temkin's theory, recalled in
 \Cref{sec:appG-graded-reduction} and used in
 \Cref{sec:ch17-residual}; the source is \cite{Tem16,Temkin2004}.
\item Raynaud's theorem on admissible formal models and their blowups,
 in the form of \cite{BoschLutkebohmert1993}, used in
 \Cref{thm:cld-formal-pl-comparison} and in \Cref{sec:ch17-formal}; the
 statement used is exactly the one already quoted in
 \Cref{thm:gm-metric-comparison}.
\item Intersection theory of line bundles on proper schemes over a graded
 field, including Kleiman's theorem that a product of nef classes is nef and
 the Nakai--Moishezon criterion.  These are used only through the graded
 dictionary of \Cref{sec:appG-graded-reduction}; the ungraded statements are
 standard.
\item Two graded geometric lemmas of \cite[\S18.9]{CLD25} that produce a
 fibration with positive relative degree from a nef and numerically
 nontrivial class.  They are stated as \Cref{lem:ch17-detect-fibration} and
 \Cref{lem:ch17-detect-model} and are used only in
 \Cref{thm:ch17-detect-subvariety}.
\item Semicontinuity of the graded transcendence degree along an affinoid
 reduction map, \cite[Proposition~11.10.4]{CLD25}, used once, in
 \Cref{prop:ch17-dimdef-discrete}.
\end{enumerate}

No local smoothing theorem for psh metrics is used anywhere.  In particular
the passage from a formal metric to a limit of smooth psh metrics, which is
how \Cref{sec:ch17-psh} produces psh-ness, goes through
\(\mathbf P^N\) and Demailly's regularized maximum, not through any
mollification on \(X\).

\subsection{Plan}
\label{sec:ch17-plan}

\Cref{sec:ch17-metrics} sets up metrics on vector bundles, the weight
cocycle, smooth metrics and their nonexistence in higher rank, and
orthogonal bases.  \Cref{sec:ch17-parlinear} defines parlinear, integral
parlinear, and formal metrics and proves their comparison.
\Cref{sec:ch17-residual} constructs the residual bundle and
\Cref{sec:ch17-dimdef} the dimension of definition, which is the invariant
that decides where a top curvature product can sit.
\Cref{sec:ch17-chern} constructs first Chern forms and currents, proves
Poincar\'e--Lelong for metrized bundles, and computes the total mass on a
proper curve.  \Cref{sec:ch17-psh} treats psh metrics and the equivalence
between positivity of the curvature current and residual nefness.
\Cref{sec:ch17-products} constructs the mixed products, proves that the
total mass is a metric-independent intersection number, and treats the
relative Monge--Amp\`ere operator.  \Cref{sec:ch17-formal} computes the
top product of formal metrics in terms of special-fibre intersection
numbers.  \Cref{sec:ch17-cl} is the recovery of Chambert--Loir measures and
\Cref{sec:ch17-cl-checks} its numerical verification.
\Cref{sec:ch17-bm-levine} proves the Bochner--Martinelli and Levine
identities.  \Cref{sec:ch17-functoriality} treats projection and base
change, and \Cref{sec:ch17-positivity-detection} the detection of positivity
on a subvariety.  \Cref{sec:ch17-dictionary} is the chapter dictionary.

\section{Metrics on line bundles and vector bundles}
\label{sec:ch17-metrics}

\subsection{Metrics, and why fibrewise norms are not enough}
\label{sec:ch17-metric-definition}

Let \(E\) be a vector bundle of rank \(r\) on a \(K\)-analytic space \(X\).

\begin{definition}[Continuous metric]
A \emph{continuous metric} on \(E\) is a family of ultrametric norms
\(\|\cdot\|_x\) on the \(\mathcal H(x)\)-vector spaces \(E(x)\), compatible
with scalar multiplication, such that \(x\mapsto\|s(x)\|_x\) is continuous
for every local analytic section \(s\).  For a line bundle this condition is
equivalent to continuity for one, and hence every, local frame.
\end{definition}

This is the definition we shall use, and for line bundles it is
unimpeachable.  For higher rank it hides a subtlety that must be brought
into the open before anything is done with it, because it is the reason the
whole chapter is organized around \emph{orthogonal} frames.

\begin{warning}{A fibrewise norm need not stay a norm after base change}
A ``family of ultrametric norms on the fibres, continuous in \(x\)'' is a
weaker datum than one might expect: it need not induce a norm after
extension of scalars.  Take \(K\) trivially valued, \(X=\mathcal M(K)\) a
point, and \(E=\cO_X\) the trivial line bundle, so that the total space is
\(\mathbf A^{1,\An}_K\).  Any continuous function
\(\mathbf A^{1,\An}_K\to\mathbf R_{\ge0}\) sending every nonzero \(K\)-point
to \(1\) and the origin to \(0\) restricts to a norm on \(E(K)=K\); this is
satisfied both by \(|T|\) and by \(|T|^2\).  But \(|T|^2\) does not induce a
norm on \(E(K)\otimes L=L\) as soon as \(L\) is a nontrivially valued
extension of \(K\), since it is not multiplicative there.  The remedy, and
the definition actually used in \cite[Definition~17.1.2]{CLD25}, is to
demand that \(\|\cdot\|\) be a continuous function on the total space
\(V(E)\) whose restriction to \(E(x)\otimes_{\cH(x)}L\) is an ultrametric
norm \emph{for every} complete valued extension \(L\) of \(\cH(x)\).  Our
formulation above is equivalent to that one once it is read, as we always
read it, as a condition on the total space; the point of the warning is that
``fibrewise ultrametric norm'' alone is not a base-change-stable condition,
and no argument below may appeal to it in that weaker form.
\end{warning}

Concretely, the datum of a continuous metric is equivalent to the following
bookkeeping, which is what we use in practice.  Let \((U_i)\) be a
\(G\)-covering trivializing \(L\), let \(s_i\) be a frame of \(L|_{U_i}\),
and set \(f_{ij}=s_i/s_j\in\cO^\times(U_i\cap U_j)\).  Giving a continuous
metric on \(L\) is the same as giving continuous positive functions
\(u_i=\|s_i\|\) on \(U_i\) satisfying \(u_i=|f_{ij}|u_j\); this is
\Cref{lem:metrics-weight-cocycle}, transported from the algebraic setting of
\Cref{chap:metrics} to the analytic one word for word.  We use the
multiplicative and the additive formulations interchangeably.

If \(e\) is a frame of a metrized line bundle \(L\), its additive weight is
\[
                         \phi_e=-\log\|e\|.
\]
For another frame \(e'=ge\),
\begin{equation}
                    \phi_{e'}=\phi_e-\log|g|.
\label{eq:cld-weight-change-frame}
\end{equation}
Thus weights are potentials modulo logarithms of invertible analytic
functions.  Tensor product, duality, and pullback are described by
\[
 \phi_{L\otimes M} =\phi_L+\phi_M,\qquad
 \phi_{L^\vee}=-\phi_L,\qquad
 \phi_{f^*L}=\phi_L\circ f.
\]

We record the formal content of these three formulas, because every
functoriality statement in the chapter is obtained from them by applying
\(\dprimeddprime\).

\begin{proposition}
\label{prop:ch17-metric-picard}
Let \(X\) be a \(K\)-analytic space.
\begin{enumerate}
\item\label{itm:ch17-pic:group}
 Tensor product makes the set \(\widehat{\Pic}(X)\) of isometry classes of
 continuously metrized line bundles on \(X\) an abelian group, and the
 forgetful map \(\widehat{\Pic}(X)\to\Pic(X)\) is a group homomorphism.
\item\label{itm:ch17-pic:kernel}
 Its kernel is canonically
 \(C^0(X,\mathbf R)/\log|\cO^\times(X)|\), where a continuous function
 \(u\) corresponds to the metric \(\|1\|=\rme^{-u}\) on \(\cO_X\).
\item\label{itm:ch17-pic:pullback}
 A morphism \(f:Y\to X\) of \(K\)-analytic spaces induces a group
 homomorphism \(f^*:\widehat{\Pic}(X)\to\widehat{\Pic}(Y)\) compatible with
 \(f^*:\Pic(X)\to\Pic(Y)\), and \(\phi_{f^*e}=f^*\phi_e\) for every frame
 \(e\).
\item\label{itm:ch17-pic:sub}
 The same statements hold verbatim with ``continuous'' replaced throughout
 by ``smooth'', by ``parlinear'', or by ``integral parlinear'', giving
 subgroups
 \(\widehat{\Pic}(X,\text{sm})\subseteq\widehat{\Pic}(X,\PL)
   \subseteq\widehat{\Pic}(X,C^0)\)
 stable under pullback.
\end{enumerate}
\end{proposition}

\begin{proof}
For \cref{itm:ch17-pic:group}, the tensor product of two metrics is the
unique metric with \(\|s\otimes s'\|=\|s\|\cdot\|s'\|\) for local sections;
uniqueness is clear because \(s\otimes s'\) generates \(L\otimes M\) locally,
and existence follows from the cocycle description, the transition functions
of \(L\otimes M\) being the products \(f_{ij}f'_{ij}\) and the candidate
norms \(u_iu'_i\) satisfying
\(u_iu'_i=|f_{ij}f'_{ij}|u_ju'_j\).  Associativity, commutativity and the
neutrality of the trivially metrized \(\cO_X\) are then immediate, and the
inverse of \((L,\|\cdot\|)\) is \((L^\vee,\|\cdot\|^\vee)\) with
\(\|s^{-1}\|=\|s\|^{-1}\) for invertible \(s\).  The forgetful map is a
homomorphism by construction.

For \cref{itm:ch17-pic:kernel}, an element of the kernel is a metric on
\(\cO_X\), that is, a positive continuous function \(u=\|1\|\); two such
give isometric metrized bundles exactly when they differ by \(|g|\) for
some \(g\in\cO^\times(X)\), since an isomorphism \(\cO_X\to\cO_X\) is
multiplication by such a \(g\).  Passing to \(-\log\) gives the stated
quotient.

For \cref{itm:ch17-pic:pullback}, define \(f^*\|\cdot\|\) as the composite of
the canonical map \(V(f^*L)\to V(L)\) with \(\|\cdot\|\); this is a
continuous metric because the map is continuous and is an isomorphism on
fibres over each point of \(Y\), and the displayed compatibility with frames
is the definition.  Compatibility with tensor products is checked on frames.

For \cref{itm:ch17-pic:sub}, all four operations---tensor, dual, pullback,
and the identification of the kernel---are expressed by sums, differences,
and pullbacks of local weights, and each of the three classes of functions
(smooth, parlinear, integral parlinear) is stable under those operations:
for smooth functions by \Cref{prop:cld-intrinsic-calculus} and
\Cref{prop:cld-form-pullback}, for parlinear functions by
\Cref{prop:ch15-pl-pullback}.  Integrality is preserved because the local
model of an integral parlinear weight is \(\log|g|\) with
\(g\in\cO^\times\), and \(f^*\log|g|=\log|f^*g|\).
\end{proof}

\subsection{Smooth metrics exist in rank one and only in rank one}
\label{sec:ch17-smooth}

\begin{definition}[Smooth metric]
\label{def:ch17-smooth-metric}
\index{metric!smooth}
A metric on a vector bundle \(E\) is \emph{smooth} if the induced function
\(V(E)\setminus\{0\}\to\mathbf R_{>0}\) on the complement of the zero
section is smooth in the sense of \Cref{chap:cld-forms}.  This is a local
condition on \(X\).
\end{definition}

\begin{proposition}
\label{prop:ch17-smooth-line}
Let \(L\) be a line bundle on \(X\), locally trivial for the usual topology.
A continuous metric on \(L\) is smooth if and only if \(\log\|s\|\) is a
smooth function on \(U\) for every invertible section \(s\) of \(L\) over an
open \(U\subseteq X\).
\end{proposition}

\begin{proof}
Necessity is immediate: \(\|s\|\) is the composite of \(\|\cdot\|\) with the
section \(U\to V(L)\setminus\{0\}\) determined by \(s\), and composition of
a smooth function with a morphism is smooth by
\Cref{prop:cld-form-pullback}.  For sufficiency, work locally and choose an
invertible \(s\) over \(U\).  It identifies \(V(L)|_U\) with
\(\mathbf A^{1,\An}\times U\), and under this identification the metric is
\((t,x)\mapsto|t|\cdot\|s(x)\|\).  Its restriction to the complement of the
zero section is \(\mathbf G_{\mathrm m}^{\An}\times U\to\mathbf R_{>0}\),
\((t,x)\mapsto|t|\,\|s(x)\|\), whose logarithm is
\(\log|t|+\log\|s(x)\|\): a sum of a pullback of a coordinate moment
function and a pullback of a smooth function, hence smooth.
\end{proof}

\begin{proposition}
\label{prop:ch17-smooth-existence}
Let \(X\) be paracompact and locally holomorphically separated, and let
\(L\) be a line bundle on \(X\) that is locally trivial for the usual
topology---for instance any line bundle on a good \(X\), and in particular
on a boundaryless \(X\).  Then \(L\) admits a smooth metric, and every
continuous metric on \(L\) is a uniform limit, on compact subsets, of smooth
metrics.
\end{proposition}

\begin{proof}
Choose an open covering \((U_i)\) of \(X\) with invertible sections \(s_i\)
of \(L|_{U_i}\), and a smooth partition of unity \((\lambda_i)\)
subordinate to it, which exists by \Cref{thm:cld-smooth-partition} because
\(X\) is paracompact and locally holomorphically separated.  There is a
unique metric on \(L\) for which
\[
 \log\|s\|(x)=\sum_i\lambda_i(x)\log|s/s_i|(x),
\]
the sum being over the indices with \(x\in U_i\): indeed the right-hand
side satisfies the cocycle rule, because replacing \(s\) by \(gs\) adds
\(\sum_i\lambda_i\log|g|=\log|g|\).  Each \(\log|s/s_i|\) is smooth where
\(s\) is invertible, so \(\log\|s\|\) is smooth there, and
\Cref{prop:ch17-smooth-line} applies.

For the density statement, fix one smooth metric \(\|\cdot\|_0\) and let
\(\|\cdot\|\) be continuous.  There is a unique continuous
\(u:X\to\mathbf R\) with \(\|\cdot\|=\rme^{u}\|\cdot\|_0\).  Given a compact
\(S\subseteq X\) and \(\varepsilon>0\), the smooth partition of unity
supplied by \Cref{thm:cld-smooth-partition} produces a smooth \(w\) with
\(|w-u|\le\varepsilon\) on \(S\): cover \(S\) by finitely many opens on
which \(u\) oscillates by less than \(\varepsilon\), and take the
corresponding convex combination of constants.  Then
\(\rme^{w}\|\cdot\|_0\) is smooth and agrees with \(\|\cdot\|\) on \(S\) up
to a factor in \([\rme^{-\varepsilon},\rme^{\varepsilon}]\).
\end{proof}

The next statement has no complex analogue and is the reason why this
chapter constructs Chern currents only for line bundles.

\begin{proposition}
\label{prop:ch17-no-smooth-higher-rank}
\index{metric!on a vector bundle}\index{metric!no smooth metric in higher rank}
Let \(r\ge2\).  A vector bundle of rank \(r\) on a nonempty
\(K\)-analytic space carries no smooth metric.  This already fails at a
point: the trivial bundle \(\cO^{r}\) on \(\mathcal M(K)\) has no smooth
metric.
\end{proposition}

\begin{proof}
Since smoothness is local and restriction of a smooth metric to a point is
smooth, it suffices to treat \(X=\mathcal M(K)\) and \(E=\cO^r\); and since
smoothness is insensitive to a complete base extension---by
\Cref{prop:cld-form-pullback} applied to the base-change morphism---we may
assume that \(K\) is algebraically closed and spherically complete.

Let \(\|\cdot\|\) be a metric on \(K^r\).  Spherical completeness makes
every finite-dimensional ultrametric normed \(K\)-vector space
orthogonalizable \cite[\S2.3]{Tem16}: there is a basis \(e_1,\dots,e_r\)
with \(\|\sum x_ie_i\|=\max_i|x_i|\,\|e_i\|\).  Writing \(f_i\) for the dual
basis and \(a_i=\|e_i\|>0\), this reads
\begin{equation}
\label{eq:ch17-orthogonal-normal-form}
 \|x\|=\max_{1\le i\le r}a_i|f_i(x)|,\qquad x\in K^r.
\end{equation}
The set \(\mathbf A^{r}(K)\) of \(K\)-points is dense in
\(\mathbf A^{r,\An}_K\), and both sides of
\eqref{eq:ch17-orthogonal-normal-form} are continuous, so the identity holds
on all of \(\mathbf A^{r,\An}_K\).

Suppose \(\|\cdot\|\) were smooth.  The function \(\log\|\cdot\|\) on
\(\mathbf A^{r,\An}_K\setminus\{0\}\) satisfies
\(\log\|tx\|=\log|t|+\log\|x\|\), so the \((1,1)\)-form
\(\dprimeddprime\log\|\cdot\|\) is invariant under the scaling action and
descends to a smooth \((1,1)\)-form \(\omega\) on
\(\mathbf P^{r-1,\An}_K\).  Its \((r-1)\)st power \(\omega^{r-1}\) would
then be a smooth \((r-1,r-1)\)-form on the \((r-1)\)-dimensional space
\(\mathbf P^{r-1,\An}_K\).  But by
\eqref{eq:ch17-orthogonal-normal-form} the weight is the pullback of
\(\max_i(\log a_i+t_i)\) under the moment map
\((\log|f_1|,\dots,\log|f_r|)\), and the \((r-1)\)st Monge--Amp\`ere measure
of that convex piecewise affine function on the quotient
\(\mathbf R^{r}/\mathbf R(1,\dots,1)\simeq\mathbf R^{r-1}\) is a Dirac mass.
Indeed the subgradient of \(t\mapsto\max_i(\log a_i+t_i)\) takes values in
the standard simplex \(\{\lambda\ge0,\sum\lambda_i=1\}\), equals a proper
face of it away from the point where all \(r\) terms agree, and equals the
whole simplex at that point; so the Alexandrov measure
\eqref{eq:appJ-alexandrov} is
\(\vol(\Delta^{r-1})=1/(r-1)!\) times a Dirac mass there, while the
superform Monge--Amp\`ere measure is \((r-1)!\) times the Alexandrov
measure and hence is the unit Dirac mass
(\Cref{ex:ch16-fan-p2} is the case \(r=3\)).  A smooth form does not
represent a Dirac mass: by \Cref{cor:ch13-measure} the measure attached to a
top-degree \(G\)-form with locally integrable coefficients gives no mass to
a weakly analytic subset of dimension \(<r-1\), whereas a point of
\(\mathbf P^{r-1,\An}_K\) is such a subset when \(r\ge2\).  This
contradiction proves the claim.
\end{proof}

\begin{example}[A rank-two bundle, carried to numbers]
\label{ex:ch17-rank-two}
\index{metric!rank-two example}
Take \(X=\mathcal M(K)\), \(E=\cO^2\) with coordinates \(x_1,x_2\), and
constants \(a_1,a_2>0\).  By \Cref{prop:ch17-no-smooth-higher-rank}, every
metric on \(E\) has the form
\[
 \|x\|=\max\{a_1|x_1|,\,a_2|x_2|\},
\]
after a change of basis, and no metric is smooth.  Let us see exactly what
its curvature is.

Put \(t=\log|x_2/x_1|\), a coordinate on the tropicalization of
\(\mathbf P^{1,\An}_K=\mathbf P(E)\) through the moment map \(x_2/x_1\).  The
tautological subbundle \(\cO(-1)\subseteq\pi^*E\) inherits the metric, and
on the chart \(x_1\ne0\) the section \(e=x/x_1\) is a frame of \(\cO(-1)\)
with
\[
 \log\|e\|=\log\|x\|-\log|x_1|
 =\max\{\log a_1,\ \log a_2+t\}=:\lambda(t).
\]
Thus \(-\lambda\) is the additive weight of the induced metric on
\(\cO(-1)\), and \(\lambda\) is the additive weight of the dual metric on
\(\cO(1)\).  It is a convex piecewise affine function of \(t\) with a single
corner at \(t_0=\log(a_1/a_2)\), where the slope jumps from \(0\) to \(1\).
By the one-dimensional Monge--Amp\`ere computation of
\Cref{chap:cld-currents}---the current of \(\max\{0,t\}\) is
\(\delta_0\)---the descended curvature on \(\mathbf P^{1,\An}_K\) is
\[
 c_1\bigl(\cO(1),\|\cdot\|^{\vee}\bigr)=\dprimeddprime\lambda=\delta_{\zeta}
 =-\,c_1\bigl(\cO(-1),\|\cdot\|\bigr),
\]
the Dirac mass at the point \(\zeta\) of \(\mathbf P^{1,\An}_K\) with
\(|x_2/x_1|(\zeta)=a_1/a_2\), that is, the Gauss point of the disc of radius
\(a_1/a_2\) about the origin.  Its total mass is
\[
 \int_{\mathbf P^{1,\An}_K}\dprimeddprime\lambda=1
 =\deg\cO_{\mathbf P^1}(1),
\]
consistently with \Cref{cor:ch17-curve-total-mass} below.  So the
determinant-free content of the rank-two metric is: one atom, at a point
determined by the ratio \(a_1/a_2\) of the two orthogonal norms, of mass
one.  When \(a_1=a_2\) the atom sits at the Gauss point of the unit disc;
when \(a_1/a_2\notin|K^\times|\) it sits at a point of type \(3\).

Two consequences deserve emphasis.  First, this single atom is what
\Cref{prop:ch17-no-smooth-higher-rank} used: it cannot be the measure of a
smooth form.  Second, the determinant \(\det E=\cO\) receives the metric
\(\|e_1\wedge e_2\|=a_1a_2\), whose weight \(-\log(a_1a_2)\) is
\emph{constant}, so \(c_1(\det\overline E)=0\).  The determinant therefore
sees none of the above.  This is the precise sense in which, for
\(r\ge2\), the metric carries strictly more information than its
determinant, and it is why no Chern--Weil theory is attempted here.
\end{example}

\subsection{Orthogonal frames}
\label{sec:ch17-orthogonal}

For higher-rank bundles one must remember more than a determinant.  The
classes most often used in the CLD theory are locally orthogonal.

\begin{definition}[Locally orthogonal metric]
A metric on \(E\) is \emph{locally orthogonal} if locally there are a frame
\(e_1,\ldots,e_r\) and positive continuous functions \(\rho_i\) such that
\[
 \left\|\sum_i f_i e_i\right\|_x
       =\max_i |f_i(x)|\rho_i(x).
\]
It is smooth, piecewise linear, or integral piecewise linear when the
functions \(-\log\rho_i\) have the corresponding property in tropical
charts.
\end{definition}

\begin{warning}{Two uses of the word ``smooth'' for higher-rank metrics}
The last sentence of the definition just given uses ``smooth'' in a weaker
sense than \Cref{def:ch17-smooth-metric}: it asks that the \(r\) functions
\(-\log\rho_i\) attached to one orthogonal frame be smooth, not that the
norm function on the total space \(V(E)\setminus\{0\}\) be smooth.  For
\(r=1\) the two agree (\Cref{prop:ch17-smooth-line}).  For \(r\ge2\) they do
not, and cannot: by \Cref{prop:ch17-no-smooth-higher-rank} no metric on a
bundle of rank \(\ge2\) is smooth in the sense of
\Cref{def:ch17-smooth-metric}, whereas the metric
\(\max\{a_1|x_1|,a_2|x_2|\}\) of \Cref{ex:ch17-rank-two} has constant, hence
smooth, functions \(-\log\rho_i\).  Whenever ``smooth metric'' occurs below
without qualification, \Cref{def:ch17-smooth-metric} is meant, and the
bundle is a line bundle.  The weaker frame-wise conditions are used only in
the parlinear and integral parlinear cases, where they are the operative
ones and where no conflict arises because those classes are not required to
be smooth on the total space.
\end{warning}

The exterior power of an orthogonal frame gives the determinant line bundle
the metric
\[
          \|e_1\wedge\cdots\wedge e_r\|
             =\prod_i\|e_i\|.
\]
Changes between orthogonal frames preserve this construction.  In this
chapter, Chern currents are needed only for line bundles and determinants;
we do not postulate a general Chern--Weil theory for arbitrary continuous
norms on vector bundles.

Three elementary facts about orthogonal frames are used repeatedly and are
worth isolating, because they are what makes the class of parlinear metrics
manageable.

\begin{lemma}
\label{lem:ch17-orthogonal-frames}
Let \(E\) be a metrized vector bundle of rank \(r\) on \(X\).
\begin{enumerate}
\item\label{itm:ch17-of:unique}
 If \((e_i)\) is an orthogonal frame, the functions \(\rho_i\) in the
 definition are determined: \(\rho_i=\|e_i\|\).
\item\label{itm:ch17-of:perturb}
 If \((e_i)\) is orthogonal and \((\varepsilon_i)\) are sections with
 \(\|\varepsilon_i\|<\|e_i\|\) pointwise on \(X\), then
 \((e_i+\varepsilon_i)\) is an orthogonal frame and
 \(\|e_i+\varepsilon_i\|=\|e_i\|\).
\item\label{itm:ch17-of:normalize}
 If \(\|e_i\|=c_i|g_i|\) for constants \(c_i>0\) and invertible analytic
 \(g_i\), then \((g_i^{-1}e_i)\) is an orthogonal frame with constant norms
 \(c_i\).  If moreover \(c_i=1\) for all \(i\), that frame is orthonormal.
\item\label{itm:ch17-of:ops}
 Orthogonal direct sums, tensor products of orthogonal frames, duals of
 orthogonal frames, and pullbacks of orthogonal frames are orthogonal, with
 \(\|e_i\otimes f_j\|=\|e_i\|\,\|f_j\|\) and
 \(\|e_i^\vee\|=\|e_i\|^{-1}\).
\end{enumerate}
\end{lemma}

\begin{proof}
\cref{itm:ch17-of:unique}: apply the defining identity to
\(s=e_i\), that is, \(f_j=\delta_{ij}\).

\cref{itm:ch17-of:perturb}: the hypothesis gives, at every point and for
every \((f_i)\),
\(\bigl\|\sum f_i\varepsilon_i\bigr\|\le\max_i|f_i|\,\|\varepsilon_i\|
 <\max_i|f_i|\,\|e_i\|=\bigl\|\sum f_ie_i\bigr\|\)
whenever some \(f_i\ne0\); the ultrametric inequality in the sharp form
``\(\|a+b\|=\|a\|\) when \(\|b\|<\|a\|\)'' then yields
\(\bigl\|\sum f_i(e_i+\varepsilon_i)\bigr\|=\bigl\|\sum f_ie_i\bigr\|
 =\max_i|f_i|\,\|e_i\|\).
In particular \((e_i+\varepsilon_i)\) spans, since a relation
\(\sum f_i(e_i+\varepsilon_i)=0\) with some \(f_i\ne0\) would force
\(\max_i|f_i|\,\|e_i\|=0\); so it is a frame, and it is orthogonal with the
same \(\rho_i\).

\cref{itm:ch17-of:normalize}: substituting \(f_i\mapsto f_ig_i^{-1}\) in the
defining identity gives
\(\bigl\|\sum f_ig_i^{-1}e_i\bigr\|=\max_i|f_i||g_i|^{-1}c_i|g_i|
 =\max_i|f_i|c_i\).

\cref{itm:ch17-of:ops}: for the direct sum and the dual this is a direct
computation with the defining identity.  For the tensor product, the
asserted norm is the tensor-product norm of two orthogonally decomposed
normed spaces, which is again orthogonally decomposed with the product
basis; independence of the chosen orthogonal frames follows because the
tensor-product norm is intrinsic.  Pullback is compatible with the defining
identity by \Cref{prop:ch17-metric-picard}\,\cref{itm:ch17-pic:pullback}.
\end{proof}

Part \cref{itm:ch17-of:perturb} of the lemma is more useful than it looks:
it says that an orthogonal frame may be modified by anything strictly
smaller without changing the norms.  This is the mechanism by which
residual data in \Cref{sec:ch17-residual} lift to analytic data.

\section{Parlinear, formal, and residual metrics}
\label{sec:ch17-parlinear}

There are three closely related, but differently presented, kinds of
integral metric.  The relations among them are: \emph{integral parlinear}
is an intrinsic analytic condition; \emph{formal} is a condition referring
to a model; and \emph{residual} is not a class of metrics at all but the
pointwise linear-algebra invariant that measures how singular an integral
parlinear metric is at a given point.  The first two coincide in the strict
case (\Cref{thm:cld-formal-pl-comparison}); the third is what makes the
curvature computable (\Cref{sec:ch17-formal}).

\subsection{The four parlinear classes}
\label{sec:ch17-four-classes}

Recall from \Cref{chap:tropical-spaces} that a function on a \(K\)-analytic
space is \emph{parlinear} if it is locally, for the \(G\)-topology, of the
form \(\max_i(c_i+\log|g_i|)\) with \(c_i\in\mathbf R\) and
\(g_i\in\cO^\times\); it is \emph{integral parlinear} if one may take a
single term on each piece of a \(G\)-covering, that is, if it is locally of
the form \(c+\log|g|\); and it is \emph{strict} if the covering may be taken
by \emph{strictly} analytic domains.  Combining the two adjectives gives
four classes, and all four occur below.

\begin{definition}[Parlinear metrics]
\label{def:ch17-parlinear-metric}
\index{metric!parlinear}\index{metric!integral parlinear}
A metrized vector bundle \(\overline E\) on \(X\) is \emph{parlinear}
(resp.\ \emph{integral parlinear}) if there is a \(G\)-covering \((X_i)\) of
\(X\) and, for each \(i\), an orthogonal frame \((e_{ij})_j\) of
\(E|_{X_i}\) whose norm functions \(\log\|e_{ij}\|\) are parlinear (resp.\
integral parlinear).  If \(X\) is strict, the \emph{strictly} parlinear and
\emph{strictly integral parlinear} classes are defined in the same way with
\(G_{\mathrm{str}}\)-coverings and strict parlinear functions.
\end{definition}

The reader should keep in mind the smallest instances.  The trivial line
bundle \(\cO_X\) with \(\|f\|=|f|\) is strictly integral parlinear.  More
generally, \(\cO_X\) with \(\|f\|=\rme^{-u}|f|\) is parlinear, integral
parlinear, strictly parlinear, or strictly integral parlinear exactly when
\(u\) is.  Thus the parlinear conditions on metrics are, in the trivialized
case, literally the corresponding conditions on weights, and no new
phenomenon has been introduced.

\begin{lemma}
\label{lem:ch17-orthonormal-refine}
Let \(\overline E\) be integral parlinear.  Then \(X\) admits a
\(G\)-covering on each member of which \(E\) has an orthogonal frame with
\emph{constant} norms.  If \(\overline E\) is strictly integral parlinear,
the covering may be taken by strictly analytic domains and the frame taken
\emph{orthonormal}.
\end{lemma}

\begin{proof}
Let \((X_i,(e_{ij})_j)\) be as in \Cref{def:ch17-parlinear-metric}.  For
fixed \(i\) there are finitely many functions \(\log\|e_{ij}\|\), each
integral parlinear; refining \(X_i\) by a \(G\)-covering \((X_{i\ell})\) we
may assume that on each \(X_{i\ell}\) every one of them equals
\(\log c_{i\ell}+\log|g_{ij\ell}|\) for some \(c_{i\ell}>0\) and invertible
\(g_{ij\ell}\).  By
\Cref{lem:ch17-orthogonal-frames}\,\cref{itm:ch17-of:normalize} the frame
\((g_{ij\ell}^{-1}e_{ij})_j\) is orthogonal on \(X_{i\ell}\) with constant
norms \(c_{i\ell}\).  In the strict case the constants may be absorbed:
strict integrality means \(c_{i\ell}\in|K^\times|\) after refinement, so
dividing by a scalar of that absolute value makes all norms equal to \(1\).
\end{proof}

We now recall the definition already in force, which the lemma shows to be
the special case \(r=1\) of \Cref{def:ch17-parlinear-metric}.

\begin{definition}[Integral PL metric]
A metric on a line bundle \(L\) is \emph{integral PL} if, locally for the
G-topology, there are a frame \(e\) and a finite subdivision by strict
analytic domains on each of which \(\|e\|=|f|\) for an invertible analytic
function \(f\), equivalently the weight of \(e\) is \(-\log|f|\).  Replacing
\(e\) there by \(f^{-1}e\) produces an orthonormal frame.  Equivalently, after
a suitable tropicalization the weight is integral affine on every cell.
An \(\mathbf R\)-PL metric is locally a finite real linear combination of
integral PL weights.
\end{definition}

\begin{proposition}
\label{prop:ch17-parlinear-stability}
Each of the four classes of \Cref{def:ch17-parlinear-metric} is stable
under orthogonal direct sums, tensor products, duals of line bundles, and
pullback along a \(K\)-morphism \(Y\to X\) (allowing \(Y\) to live over a
complete valued extension of \(K\)).  Moreover, if \(\overline E\) is
parlinear, \(s\) is a section of \(E\), and \(r>0\), then
\(\{x\in X:\|s(x)\|\le r\}\) is an analytic domain of \(X\), strictly
analytic if \(\overline E\) is strictly parlinear and \(r\) lies in the
divisible closure of \(|K^\times|\).
\end{proposition}

\begin{proof}
Stability under the four operations is
\Cref{lem:ch17-orthogonal-frames}\,\cref{itm:ch17-of:ops} together with the
observation that sums, differences, maxima and pullbacks of (integral)
parlinear functions are (integral) parlinear---for maxima and sums by
definition, for pullbacks by \Cref{prop:ch15-pl-pullback}.  For tensor
products one must note that a product of two integral parlinear local models
\(c\log|g|\) and \(c'\log|g'|\) is \(\log|gg'|+\log(cc')\), again integral
parlinear.

For the last assertion, work \(G\)-locally and choose an orthogonal frame
\((e_i)\).  Writing \(s=\sum a_ie_i\) with \(a_i\) analytic and
\(\rho_i=\|e_i\|\), the condition \(\|s\|\le r\) reads
\(\max_i|a_i|\rho_i\le r\), that is, the finite conjunction of the
conditions \(|a_i|\le r/\rho_i\).  Each \(\rho_i\) is the exponential of a
parlinear function, so on a further \(G\)-refinement each is of the form
\(c|g|\), and the condition becomes \(|a_ig^{-1}|\le r/c\), which cuts out
an analytic domain.  Strictness under the stated hypothesis on \(r\) is the
same computation with \(c\in|K^\times|\).
\end{proof}

\subsection{Formal metrics}
\label{sec:ch17-formal-metrics}

Let \(\mathfrak X\) be an admissible formal \(K^\circ\)-scheme with generic
fibre \(X\), and let \(\mathfrak E\) be a locally free sheaf on
\(\mathfrak X\).  On the generic fibre, a formal frame is declared
orthonormal.  This gives the \emph{formal metric} on \(E=\mathfrak E_\eta\).
The definition is independent of the formal frame because transition
matrices and their inverses have power-bounded coefficients.

Let us write that argument out, since it is the only content of the
construction and it is used again in \Cref{sec:ch17-residual}.  Let
\(\mathfrak U\) be a formal open trivializing \(\mathfrak E\), with basis
\((\mathfrak e_1,\dots,\mathfrak e_r)\), and let \(\varphi:Y\to X\) be a
morphism with \(\varphi(Y)\subseteq\mathfrak U_\eta\).  Every section \(s\)
of \(\varphi^*E\) is uniquely \(s=\sum f_i\varphi^*\mathfrak e_i\) with
\(f_i\in\cO(Y)\), and one sets
\(\|s\|=\max_i|f_i|\).  If \((\mathfrak e'_i)\) is a second formal basis on
\(\mathfrak U'\ni\varphi(Y)\), the transition matrix
\(A\in\mathrm{GL}_r(\cO(\mathfrak U\cap\mathfrak U'))\) has all entries of
\(A\) and of \(A^{-1}\) in \(\cO^\circ\), hence of absolute value \(\le1\)
at every point of the generic fibre.  Writing \((f'_i)={}^{\mathrm t}\!A\cdot(f_i)\)
gives \(\max_i|f'_i|\le\max_i|f_i|\) and, applying the same to \(A^{-1}\),
the reverse inequality.  Thus \(\max_i|f_i|=\max_i|f'_i|\) and the recipe is
independent of the formal basis.  It manifestly satisfies the four axioms of
\cite[Remark~17.1.5]{CLD25}, hence defines a metric; and by construction
every formal basis is an orthonormal frame, so a formal metric is strictly
integral parlinear.

\begin{example}[The Weil metric on \(\cO(m)\)]
\label{ex:ch17-weil-metric}
\index{metric!Weil}\index{metric!canonical formal on projective space}
Let \(\mathbf P^n_{K^\circ}=\Proj K^\circ[T_0,\dots,T_n]\) with its formal
completion, and give \(\cO(m)\) the resulting formal metric, called the
\emph{canonical formal metric} or \emph{Weil metric}.  Let \(\mathfrak U_i\)
be the formal open \(T_i\ne0\); its generic fibre is the analytic domain of
\(\mathbf P^{n,\An}_K\) where \(|x_i|=\max_j|x_j|\), and \(s_{T_i}\) is a
formal frame there, so \(\|s_{T_i}\|\equiv1\) on it.  For a homogeneous
\(P\) of degree \(m\), writing \(s_P=(P/T_i^m)s_{T_i}^{\otimes m}\) on
\(\mathfrak U_{i,\eta}\) gives the familiar formula
\begin{equation}
\label{eq:ch17-weil-metric}
 \|s_P\|(x)=\frac{|P(x_0,\dots,x_n)|}{\max(|x_0|,\dots,|x_n|)^m}.
\end{equation}
The weight of the frame \(s_{T_0}\) of \(\cO(1)\) on the chart \(x_0\ne0\)
is therefore
\(\lambda_0=\max\{0,\log|x_1/x_0|,\dots,\log|x_n/x_0|\}\), the pullback of
the maximum of \(0\) and the coordinates under the moment map
\((x_1/x_0,\dots,x_n/x_0)\).  This is the toric weight of
\Cref{ex:ch16-fan-p2} for \(n=2\), and it is exactly the weight
\eqref{eq:p1-standard-weight} of the standard model metric for \(n=1\).
The metric is strictly integral parlinear, and it is \emph{not} smooth: its
curvature is computed in \Cref{ex:ch17-weil-curvature}.
\end{example}

\begin{theorem}
\label{thm:cld-formal-pl-comparison}
Let \(X\) be a paracompact strict \(K\)-analytic space over a nontrivially
valued field.  A metric on a vector bundle is formal if and only if it is
strictly integral PL.  In particular, formal metrics exist.  For line
bundles on a locally holomorphically separated space, such metrics are dense
among continuous metrics for uniform convergence on compact subsets.
\end{theorem}
\begin{proof}
If $(\mathfrak X,\mathfrak E)$ represents a formal metric, trivialize
$\mathfrak E$ on formal affine opens.  The induced frames on the generic
fibres are orthonormal, so the metric is strictly integral PL.

Conversely, take a locally finite strict analytic cover carrying orthonormal
frames for the integral PL metric.  Refine it by compact analytic domains.
On an overlap, the change-of-frame matrix $A$ is an isometry.  Thus the
entries of both $A$ and $A^{-1}$ are power bounded.  After replacing the
chosen formal models by common admissible formal blowups, these entries are
formal functions.  The matrix is then invertible over the formal structure
sheaf and the cocycle identity still holds.  It glues the trivial formal
bundles to a locally free sheaf $\mathfrak E$ whose formal metric is the
given one.  Paracompactness and local finiteness permit the local blowups to
be dominated by one admissible formal model.  This proves the equivalence.

The same gluing construction, starting from arbitrary analytic
trivializations and enlarging the local lattices until every transition
matrix and its inverse are power bounded on a compact refinement, produces
at least one formal metric.  For a line bundle, fix such a metric
$\|\cdot\|_0$.  The logarithmic ratio of any continuous metric to
$\|\cdot\|_0$ is a continuous function.  Smooth functions of finitely many
$-\log|f|$ approximate it uniformly on a compact set by the tropical
Stone--Weierstrass argument, and a locally finite PL partition glues the
approximants.  Exponentiating the resulting integral-PL functions gives
formal metrics converging uniformly on compact subsets.

Formal gluing for vector bundles and the compact-open density statement are
proved in \cite[\S17.5]{CLD25}.  For line bundles, the equivalence is also
\cite[Proposition~2.10]{GM19}; on a proper algebraic variety the
algebraization step is \cite[Theorem~1.1]{GM19}.  We do not use that last
step on a nonproper analytic space.
\end{proof}

\section{The residual bundle}
\label{sec:ch17-residual}

Integral PL metrics possess a pointwise reduction finer than the special
fibre of any one model.  Let \(x\in X\), let
\(\widetilde{\mathcal H(x)}^{\bullet}\) be Temkin's graded residue field,
and let \(\operatorname{red}(X,x)\) be the reduction of the germ.

\subsection{From a normed fibre to a graded vector space}
\label{sec:ch17-residual-fibre}

We recall the local linear algebra, in the form in which it will be used.
Let \((V,\|\cdot\|)\) be a finite-dimensional normed vector space over
\(\cH(x)\).  Its \emph{graded reduction} is
\[
 \widetilde V^{\bullet}
 =\bigoplus_{r>0}
   \{v\in V:\|v\|\le r\}\big/\{v\in V:\|v\|<r\},
\]
a graded vector space over the graded residue field
\(\widetilde{\cH(x)}^{\bullet}\) of \Cref{sec:appG-graded-reduction}, with
the class of \(v\) in degree \(\|v\|\) written \(\widetilde v\).  The
comparison with the norm is: \(\widetilde v\ne0\) exactly when
\(\|v\|\) is realized, and \(\widetilde{v+w}=\widetilde v\) when
\(\|w\|<\|v\|\).

\begin{lemma}
\label{lem:ch17-orthogonal-reduces}
If \(V\) has an orthogonal basis \(e_1,\dots,e_r\), with
\[
 \Bigl\|\sum_i\lambda_ie_i\Bigr\|
 =\max_i|\lambda_i|\,\|e_i\|,
\]
then
\(\widetilde e_1,\dots,\widetilde e_r\) is a homogeneous basis of
\(\widetilde V^{\bullet}\) over \(\widetilde{\cH(x)}^{\bullet}\), with
\(\deg\widetilde e_i=\|e_i\|\).  In particular
\(\dim\widetilde V^{\bullet}=\dim V\).
\end{lemma}

\begin{proof}
Homogeneity and the degrees are the definition.  For independence, suppose
\(\sum\widetilde\lambda_i\widetilde e_i=0\) with
\(\widetilde\lambda_i\) homogeneous of degrees \(r_i\) such that all
products \(r_i\|e_i\|\) are equal to a common \(r>0\), and not all
\(\widetilde\lambda_i\) zero.  Lift to \(\lambda_i\in\cH(x)\) with
\(|\lambda_i|\le r_i\), the inequality strict exactly when
\(\widetilde\lambda_i=0\).  Then
\(\bigl\|\sum\lambda_ie_i\bigr\|=\max_i|\lambda_i|\,\|e_i\|=r\), because
some index attains \(|\lambda_i|=r_i\); hence
\(\sum\widetilde\lambda_i\widetilde e_i=\widetilde{\sum\lambda_ie_i}\ne0\),
a contradiction.  For spanning, let \(v\in V\) have \(\|v\|=r\) and write
\(v=\sum\lambda_ie_i\); then \(\max_i|\lambda_i|\,\|e_i\|=r\), and
\(\widetilde v=\sum_{i\in I}\widetilde\lambda_i\widetilde e_i\) where \(I\)
is the set of indices attaining the maximum.
\end{proof}

\subsection{The residual bundle of an integral parlinear metric}
\label{sec:ch17-residual-bundle}

\begin{definition}[Residual bundle]
For an integral PL metrized vector bundle \(\overline E\), the
\emph{residual bundle} \(\widetilde E_x\) is obtained from sections of norm
at most one modulo those of norm strictly less than one.  It is a graded
vector bundle on the graded reduction of the germ \((X,x)\).  In the strict
case its degree-one part descends to the ordinary reduction.
\end{definition}

The definition needs unwinding, because ``sections of norm at most one'' is
a condition that depends on a domain, and it is precisely that dependence
that produces a \emph{bundle} rather than a vector space.  Write
\(\widetilde{(X,x)}=P_{\widetilde{\cH(x)}/\widetilde K}\{\dots\}\) for the
graded Zariski--Riemann space of the germ, whose points are the graded
valuations of \(\widetilde{\cH(x)}\) trivial on \(\widetilde K\) that are
``visible'' from \((X,x)\); this is the space introduced in
\Cref{sec:appG-graded-reduction}.  Given a nonzero homogeneous
\(\sigma\in\widetilde{E(x)}^{\bullet}\) of degree \(r\) and a point \(v\) of
\(\widetilde{(X,x)}\), choose an analytic domain \(V\ni x\) with
\(v\in\widetilde{(V,x)}\), an orthogonal frame \((e_i)\) of \(E|_V\) with
constant norms \(r_i\) (\Cref{lem:ch17-orthonormal-refine}), and a section
\(s=\sum a_ie_i\) of \(E\) over \(V\) with \(\widetilde{s(x)}=\sigma\).  Set
\begin{equation}
\label{eq:ch17-residual-valuation}
 v(\sigma)=\max_{i\in I}v\bigl(\widetilde{a_i(x)}\bigr),
 \qquad
 I=\{i:|a_i(x)|r_i=r\}.
\end{equation}

\begin{proposition}
\label{prop:ch17-residual-welldefined}
The quantity \eqref{eq:ch17-residual-valuation} does not depend on the
choices of \(V\), \((e_i)\) and \(s\); the assignment
\(\sigma\mapsto(v\mapsto v(\sigma))\) is a metric on the graded vector space
\(\widetilde{E(x)}^{\bullet}\) over \(\widetilde{(X,x)}\); and the
associated graded sheaf \(\widetilde E_x\) is a graded vector bundle of rank
\(r=\operatorname{rk}E\).  If \(V\) carries an orthogonal frame with
constant norms, the reductions \(\widetilde{e_i(x)}\) form an orthonormal
basis of \(\widetilde E_x\) over \(\widetilde{(V,x)}\).
\end{proposition}

\begin{proof}
Fix an affinoid \(W\ni x\) and a section \(t\) of \(E\) over \(W\) with
\(\widetilde{t(x)}=\sigma\); such data exist because the germ is almost
good, so that a homogeneous class of degree \(r\) lifts to a section after
shrinking, and by \Cref{lem:ch17-orthogonal-reduces} the lift may be taken
with \(\|t(x)\|=r\).  For an analytic function \(a\) on \(W\) not vanishing
at \(x\), let \(W'\subseteq W\) be the analytic domain
\(\{\|at\|\le r|a(x)|\}\), which is a domain by
\Cref{prop:ch17-parlinear-stability}.

We claim: with \(V,(e_i),s\) as in the construction, \(v\in\widetilde{(W',x)}\)
if and only if \(v(\widetilde{a(x)})\cdot\max_{i\in I}v(\widetilde{a_i(x)})\le1\).
Indeed, we may replace \(W\) by \(V\cap W\) and, since
\(\widetilde{s(x)}=\widetilde{t(x)}\) forces \(\|s(x)-t(x)\|<r\), we may
replace \(t\) by \(s\) without changing \(W'\) near \(x\); shrinking \(V\)
so that \(\|s\|=\max_{i\in I}|a_i|r_i\) identically, the domain \(W'\) is
cut out by the finitely many inequalities \(|aa_i|\le|a(x)a_i(x)|\) for
\(i\in I\).  By the description of germ reductions in
\Cref{sec:appG-graded-reduction}, \(\widetilde{(W',x)}\) is the locus
\(\{v:v(\widetilde{a_i(x)}\,\widetilde{a(x)})\le1\ \forall i\in I\}\), which
is the claim.

The claim shows that the set
\(\Gamma=\{\gamma\in v(\widetilde{\cH(x)}^\times):
 \gamma\cdot\max_{i\in I}v(\widetilde{a_i(x)})\le1\}\)
is determined by \(\sigma\) and \(v\) alone, since \(W'\) is.  As
\(v(\sigma)\) is the inverse of the largest element of \(\Gamma\), it is
determined too.  This proves independence of the choices.

That \(\sigma\mapsto(v\mapsto v(\sigma))\) is a metric over
\(\widetilde{(X,x)}\) in the sense of \Cref{sec:appG-graded-reduction}
follows from \eqref{eq:ch17-residual-valuation}: it is multiplicative in
scalars, ultrametric in \(\sigma\), and takes the value \(0\) only at
\(\sigma=0\).  Finally, if \(V\) carries an orthogonal frame with constant
norms, then applying \eqref{eq:ch17-residual-valuation} to
\(\sigma=\widetilde{e_i(x)}\) gives \(v(\sigma)=v(1)=1\) for every
\(v\in\widetilde{(V,x)}\), so the basis of
\Cref{lem:ch17-orthogonal-reduces} is orthonormal over \(\widetilde{(V,x)}\);
and since integrality of the metric provides finitely many such \(V\) whose
germ reductions cover \(\widetilde{(X,x)}\), the metrized graded space is
locally trivial of rank \(r\), that is, a graded vector bundle.
\end{proof}

\begin{remark}[Why a bundle and not a vector space]
\label{rem:ch17-residual-bundle-not-space}
The underlying graded vector space \(\widetilde{E(x)}^{\bullet}\) is the
generic fibre of \(\widetilde E_x\) and carries no information beyond its
dimension.  All the content is in the metric over \(\widetilde{(X,x)}\),
that is, in \emph{which} lattices of \(\widetilde{E(x)}^{\bullet}\) are
distinguished at each graded valuation.  For a line bundle this is a graded
line bundle on a graded scheme, and it has Chern classes and intersection
numbers.  Those numbers are the coefficients of the curvature measure
(\Cref{thm:ch17-atomic-masses}).
\end{remark}

\subsection{Functoriality, and the formal case}
\label{sec:ch17-residual-functorial}

\begin{proposition}
\label{prop:ch17-residual-functorial}
Let \(\overline E,\overline F\) be integral parlinear metrized bundles on
\(X\) and \(x\in X\).
\begin{enumerate}
\item\label{itm:ch17-res:ops}
 \(\widetilde{(E\oplus F)}_x=\widetilde E_x\oplus\widetilde F_x\),
 \(\widetilde{(E\otimes F)}_x=\widetilde E_x\otimes\widetilde F_x\), and
 for line bundles \(\widetilde{(L^{\vee})}_x=(\widetilde L_x)^{\vee}\) and
 \(\widetilde{(L^{\otimes m})}_x=(\widetilde L_x)^{\otimes m}\).
\item\label{itm:ch17-res:pull}
 If \(f:(Y,y)\to(X,x)\) is a morphism of germs over a complete valued
 extension, then \(\widetilde{(f^*E)}_y\) is the pullback of
 \(\widetilde E_x\) along the induced map
 \(\widetilde{(Y,y)}\to\widetilde{(X,x)}\).
\item\label{itm:ch17-res:iso}
 If two integral parlinear metrics on \(E\) have isomorphic residual
 bundles at \(x\) by an isomorphism of the underlying graded spaces
 preserving degrees, then the two metrized bundles are isometric on a
 neighbourhood of \(x\).
\end{enumerate}
\end{proposition}

\begin{proof}
\cref{itm:ch17-res:ops} and \cref{itm:ch17-res:pull} follow from
\eqref{eq:ch17-residual-valuation} and
\Cref{lem:ch17-orthogonal-frames}\,\cref{itm:ch17-of:ops}: choose orthogonal
frames with constant norms for \(E\) and \(F\) on a common domain; the
induced frames of the direct sum, tensor product and dual are again such,
their reductions are the corresponding bases of the residual bundles by
\Cref{lem:ch17-orthogonal-reduces}, and orthonormality over the germ
reduction is inherited.  For \cref{itm:ch17-res:pull} one uses in addition
that \(\widetilde{f^*a(y)}\) is the image of \(\widetilde{a(x)}\).

For \cref{itm:ch17-res:iso}, let \(\theta\) be the given isomorphism.  Choose
a basis of \(\widetilde{E(x)}^{\bullet}\); by
\Cref{lem:ch17-orthogonal-frames}\,\cref{itm:ch17-of:perturb} and the lifting
argument of \Cref{prop:ch17-residual-welldefined} there are, after shrinking,
frames \((e_i)\) and \((f_i)\) of \(E\) near \(x\), orthogonal for the two
metrics respectively, with \(\widetilde{f_i(x)}=\theta(\widetilde{e_i(x)})\)
and \(\|f_i(x)\|'=\|e_i(x)\|\).  The isomorphism \(u\) sending \(e_i\) to
\(f_i\) is then an isometry at \(x\), and both metrics being determined near
\(x\) by their orthogonal frames, \(u\) is an isometry on a neighbourhood.
This is \cite[Lemma~17.3.2 and Theorem~17.4.9]{CLD25}.
\end{proof}

The next statement is the bridge between the abstract residual construction
and the concrete geometry of a model, and it is used in every computation of
\Cref{sec:ch17-formal}.

\begin{lemma}
\label{lem:ch17-residual-formal}
Let \(\mathfrak X\) be a separated, quasi-paracompact formal
\(K^\circ\)-scheme, topologically locally of finite type, let
\(\mathfrak E\) be a locally free sheaf of rank \(r\) on \(\mathfrak X\),
and give \(E=\mathfrak E_\eta\) the formal metric.  Let \(x\in X\), let
\(\xi=\operatorname{red}_{\mathfrak X}(x)\), and let
\(V=\overline{\{\xi\}}\subseteq\mathfrak X_s\) with its reduced structure.
Then \(V\) is of finite type over \(\widetilde K\),
\(\widetilde{(X,x)}^{1}=P_{\widetilde{\cH(x)}^1/\widetilde K^1}\{V\}\), and
\((V,\mathfrak E|_V)\) is a premodel of
\(\bigl(\widetilde{(X,x)}^1,\widetilde E_x^{\,1}\bigr)\).
\end{lemma}

\begin{proof}
Choose a locally finite open cover \((\mathfrak U_i)\) of \(\mathfrak X\)
with bases \((\mathfrak e_{ij})_j\) of \(\mathfrak E|_{\mathfrak U_i}\), and
let \(X_i=\mathfrak U_{i,\eta}\).  By construction the \(e_{ij}\) are an
orthonormal frame of \(E|_{X_i}\), and by anticontinuity of the reduction
map (\Cref{prop:berk-reduction-anticontinuous}) the union of those \(X_i\)
with \(x\in X_i\) is a neighbourhood of \(x\); write \(I\) for that finite
set of indices and \(V_i=\mathfrak U_{i,s}\cap V\), an open cover of \(V\).

Let \(\mathfrak U\ni\xi\) be affine open and \(\sigma\) a section of
\(\mathfrak E|_{\mathfrak U_s}\).  Lift \(\sigma\) to a section
\(\mathfrak s\) of \(\mathfrak E\) on \(\mathfrak U\), and regard
\(\mathfrak s\) as a section \(s\) of \(E\) on \(\mathfrak U_\eta\).  If
\(\mathfrak t\) is a second lift whose reduction \(\tau\) satisfies
\(\tau-\sigma=0\) at \(\xi\), then all coordinates of
\(\mathfrak t-\mathfrak s\) in a formal basis lie in the maximal ideal at
\(\xi\), hence have absolute value \(<1\) at \(x\), so
\(\|(t-s)(x)\|<1\) and \(\widetilde{t(x)}^{1}=\widetilde{s(x)}^{1}\).  Thus
\(\widetilde{s(x)}^{1}\) depends only on \(\sigma(\xi)\), compatibly with
restriction, and we obtain a natural map
\(F:=\mathfrak E\otimes\kappa(\xi)\to\widetilde{E(x)}^{1}\).  For
\(i\in I\), the \(\mathfrak e_{ij}(\xi)\) form a \(\kappa(\xi)\)-basis of
\(F\), while the \(\widetilde{e_{ij}(x)}^{1}\) form a
\(\widetilde{\cH(x)}^1\)-basis of \(\widetilde{E(x)}^{1}\) by
\Cref{lem:ch17-orthogonal-reduces}; so
\(F\otimes_{\kappa(\xi)}\widetilde{\cH(x)}^1\to\widetilde{E(x)}^{1}\) is an
isomorphism.

Finally, the description of the germ reduction of a formal generic fibre
(\Cref{sec:berk-reduction}) gives
\(\widetilde{(X_i,x)}^1=P_{\widetilde{\cH(x)}^1/\widetilde K^1}\{V_i\}\),
hence the asserted description of \(\widetilde{(X,x)}^1\) by taking the
union over \(i\in I\).  Since the \(\widetilde{e_{ij}(x)}^1\) are an
orthonormal basis of \(\widetilde E_x^{\,1}\) over
\(\widetilde{(X_i,x)}^1\) by
\Cref{prop:ch17-residual-welldefined}, and they are the images of the
\(\mathfrak e_{ij}(\xi)\), the pair \((V,\mathfrak E|_V)\) satisfies the two
conditions characterizing a premodel: generic fibre \(F\), and a basis over
each \(V_i\).  This is \cite[Lemma~17.7.10]{CLD25}.
\end{proof}

\begin{example}[The residual bundle of the Weil metric]
\label{ex:ch17-residual-weil}
Continue \Cref{ex:ch17-weil-metric}.  Let \(x\in\mathbf P^{n,\An}_K\) and
let \(\xi\) be its reduction in \(\mathbf P^n_{\widetilde K}\), with
\(V=\overline{\{\xi\}}\).  By \Cref{lem:ch17-residual-formal}, the residual
line bundle \(\widetilde{\cO(m)}_x\) has premodel
\((V,\cO_{\mathbf P^n_{\widetilde K}}(m)|_V)\).  Two consequences, both used
later:
\begin{itemize}
\item if \(x\) is the unique point reducing to the generic point of
 \(\mathbf P^n_{\widetilde K}\)---the Gauss point of the standard
 simplex---then \(V=\mathbf P^n_{\widetilde K}\) and the residual bundle has
 a model on an \(n\)-dimensional proper variety, but on nothing smaller;
\item if \(\dim V=d<n\), the residual bundle already has a model on the
 \(d\)-dimensional \(V\).
\end{itemize}
This is the source of the dimension count of the next subsection, and of the
support description in \Cref{ex:ch17-support-projective}.
\end{example}

\section{The dimension of definition}
\label{sec:ch17-dimdef}

The residual bundle at a point is a graded bundle on a germ reduction whose
transcendence degree may be as large as \(\dim_xX\).  What matters for
curvature is how small a graded variety already carries the bundle.

\begin{definition}[Dimension of definition]
The \emph{dimension of definition} of a residual line bundle is the least
dimension of a proper graded variety on which it has a model.  For several
residual line bundles, their simultaneous dimension of definition is
defined by requiring one common model.
\end{definition}

Following \cite[\S17.6]{CLD25} we distinguish two variants.  The
\emph{residual} dimension of definition at \(x\) of a family
\((\overline E_i)\) of integral parlinear bundles is the least transcendence
degree over \(\widetilde K\) of a graded field of definition of
\(\bigl(\widetilde{(X,x)},(\widetilde{E_{i,x}})\bigr)\); equivalently, the
least dimension of a common premodel.  The \emph{analytic} dimension of
definition is the least \(d\) such that there are an almost good germ
\((Y,y)\) of dimension \(d\), integral parlinear metrized bundles
\((\overline F_i)\) on it, and a flat boundaryless morphism of germs
\(f:(X,x)\to(Y,y)\) with \(E_i\simeq f^*F_i\) as metrized bundles.  The
analytic dimension is always at least the residual one, and equals it when
\(\widetilde K\) is perfect; in general one must pass to a \(p\)-power
tensor operation, which is why the statement of
\Cref{thm:ch17-atomic-support} is insensitive to that subtlety.

The following finiteness statement is the one used everywhere.

\begin{proposition}
\label{prop:ch17-dimdef-discrete}
Let \(X\) be a separated \(K\)-analytic space of dimension \(\le n\) and let
\((\overline E_i)\) be a finite family of integral parlinear metrized
bundles on \(X\).  For \(d\in\mathbf N\), the set \(W^d\) of points at which
the simultaneous residual dimension of definition is \(\ge d\) is contained
in a weakly analytic subset of dimension \(\le n-d\).  In particular
\(W^n\) is closed and discrete in \(X\).
\end{proposition}

\begin{proof}
Fix \(x\in X\) and, using
\Cref{lem:ch17-orthonormal-refine}, a finite family \((X_j)_{j\in J}\) of
affinoid domains containing \(x\), whose union is a neighbourhood of \(x\),
such that each \(E_i|_{X_j}\) has an orthogonal frame
\((e_{ij\alpha})_\alpha\) of constant norms.  For \(\emptyset\ne P\subseteq
J\) write \(X_P=\bigcap_{j\in P}X_j\) and let
\(\rho_P:X_P\to\widetilde{X_P}\) be the reduction of the affinoid.  Let
\(U\) be the topological interior of \(\bigcup_jX_j\).  It suffices to prove
that \(W^d\cap U\) is weakly analytic of dimension \(\le n-d\).

Let \(y\in U\) and let \(J(y)=\{j:y\in X_j\}\).  Put
\(A=\cO(X_{J(y)})\), \(M_i=E_i(X_{J(y)})\) and \(\eta=\rho_{J(y)}(y)\), with
graded residue field \(\kappa(\eta)\).  Evaluation at \(y\) is a contracting
morphism \(A\to\cH(y)\), so it induces
\(\widetilde A\to\widetilde{\cH(y)}\) factoring through \(\kappa(\eta)\).
Equipping \(M_i\) with the spectral seminorm makes evaluation a contracting
morphism of normed modules, whence a morphism of graded
\(\widetilde A\)-modules \(\widetilde{M_i}\to\widetilde{E_i(y)}\) factoring
through \(N_i=\widetilde{M_i}\otimes_{\widetilde A}\kappa(\eta)\).

Fix \(j\in J(y)\).  Orthogonality of \((e_{ij\alpha})\) gives
\(\|\sum_\alpha\lambda_\alpha e_{ij\alpha}\|
 =\max_\alpha\|\lambda_\alpha\|\,\|e_{ij\alpha}\|\)
for \(\lambda_\alpha\in A\), so the \(\widetilde{e_{ij\alpha}}\) form a
basis of \(\widetilde{M_i}\) over \(\widetilde A\) and their images
a basis of \(N_i\) over \(\kappa(\eta)\).  Since the \(e_{ij\alpha}(y)\)
form an orthogonal basis of \(E_i(y)\), their reductions form a basis of
\(\widetilde{E_i(y)}\) by \Cref{lem:ch17-orthogonal-reduces}.  Hence
\(N_i\otimes_{\kappa(\eta)}\widetilde{\cH(y)}\to\widetilde{E_i(y)}\) is an
isomorphism, and each \(\widetilde{e_{ij\alpha}(y)}\) lies in \(N_i\).

Because \(y\) lies in the interior of \(\bigcup_jX_j\), the union of the
\(X_j\) for \(j\in J(y)\) is a neighbourhood of \(y\), so
\(\widetilde{(X,y)}=\bigcup_{j\in J(y)}\widetilde{(X_j,y)}\), and
\(\widetilde{(X_j,y)}=P_{\widetilde{\cH(y)}/\widetilde K}\{\widetilde{B_j}\}\)
with \(B_j=\cO(X_j)\).  The morphism \(\widetilde{B_j}\to\widetilde{\cH(y)}\)
factors through \(\widetilde A\), hence through \(\kappa(\eta)\); so
\(\widetilde{B_j}\subseteq\kappa(\eta)\) for all \(j\in J(y)\).  Together
with \(\widetilde{e_{ij\alpha}(y)}\in N_i\), this shows that
\(\widetilde{(X,y)}\) is the preimage of an open subset of
\(P_{\kappa(\eta)/\widetilde K}\) and each \(\widetilde{E_{i,y}}\) is a
pullback from there.  Therefore the simultaneous residual dimension of
definition at \(y\) is at most
\(\operatorname{tr.deg}(\kappa(\eta)/\widetilde K)\).

Consequently
\[
 W^d\cap U\subseteq
 \bigcup_{\emptyset\ne P\subseteq J}
  \bigl\{z\in X_P:
   \operatorname{tr.deg}\bigl(\kappa(\rho_P(z))/\widetilde K\bigr)\ge d
  \bigr\},
\]
and each set on the right is weakly analytic of dimension at most \(n-d\)
in the sense of \Cref{def:trop-weakly-analytic}.  That last assertion is the
semicontinuity of the graded transcendence degree along an affinoid
reduction map, \cite[Proposition~11.10.4]{CLD25}; it is the one input of
this proof that is not proved in this book, and it is used nowhere else.
Intersecting with \(U\) preserves both properties, since a closed subset of
a weakly analytic set of dimension \(\le n-d\) is again one.  For \(d=n\), a
weakly analytic set of dimension \(\le0\) is closed and discrete.
\end{proof}

This invariant controls where a top curvature product can live.  On a pure
\(n\)-dimensional boundaryless space, the top product of \(n\) integral PL
metrics vanishes near a point at which their simultaneous residual dimension
is less than \(n\); see \cite[\S\S17.3--17.4, 18.4]{CLD25}.  The precise
statement is \Cref{thm:ch17-atomic-support} below, and
\Cref{prop:ch17-dimdef-discrete} is what makes its conclusion---a locally
finite sum of Dirac masses---meaningful.

\begin{example}[Dimension of definition on the projective line]
\label{ex:ch17-dimdef-p1}
Let \(K\) be discretely valued and let \(L=\cO(1)\) on
\(\mathbf P^{1,\An}_K\) with the Weil metric.  Let \(x\in\mathbf P^{1,\An}_K\).
\begin{itemize}
\item If \(x=\zeta_{0,1}\) is the Gauss point, its reduction is the generic
 point of \(\mathbf P^1_{\widetilde K}\), so
 \(V=\mathbf P^1_{\widetilde K}\) and, by
 \Cref{lem:ch17-residual-formal}, \(\widetilde L_x\) has premodel
 \((\mathbf P^1_{\widetilde K},\cO(1))\).  No zero-dimensional premodel
 exists, because \(\cO(1)\) is not trivial; the dimension of definition is
 \(1\).
\item If \(x\ne\zeta_{0,1}\), its reduction is a closed point \(\xi\), so
 \(V=\{\xi\}\) is a point and \(\widetilde L_x\) is defined over a
 zero-dimensional graded variety: the dimension of definition is \(0\).
\end{itemize}
Thus \(W^1=\{\zeta_{0,1}\}\), a single point, in accordance with
\Cref{prop:ch17-dimdef-discrete}; and
\Cref{thm:ch17-atomic-support} will say that \(c_1(\overline L)\) is
supported there.  Comparing with the computation
\(c_1(\overline{\cO(1)})=\delta_{\zeta_{0,1}}\) of
\Cref{ex:ch17-weil-curvature}, the support is exactly \(W^1\), not merely
contained in it---as it must be, since the mass is \(1\ne0\).
\end{example}

\section{First Chern forms and currents}
\label{sec:ch17-chern}

We use the normalization
\[
                              \dc=\ddprime,\qquad \ddc=\dprimeddprime.
\]
Let \(\overline L=(L,\|\cdot\|)\) be smoothly metrized.  In a local frame
\(e\), define
\begin{equation}
             c_1(\overline L)|_U=\dprimeddprime\phi_e,
             \qquad \phi_e=-\log\|e\|.
\label{eq:cld-c1-local}
\end{equation}
If \(e'=ge\) with \(g\) invertible, Poincar\'e--Lelong gives
\(\dprimeddprime\log|g|=0\) on \(U\).  Hence
\eqref{eq:cld-c1-local} is independent of the frame and defines a closed
smooth \((1,1)\)-form.

For a locally integrable metric the same formula defines a current:
\[
                   c_1(\overline L)|_U=\dprimeddprime[\phi_e].
\]
It is additive under tensor products, changes sign under duality, and is
compatible with pullback whenever pullback of the relevant currents is
defined.

Let us record the construction once, carefully, because the sheaf-theoretic
formulation is the one that makes the frame independence a triviality rather
than a computation.

\begin{proposition}
\label{prop:ch17-c1-construction}
Let \(X\) be boundaryless of pure dimension \(n\) and let \(L\) be a line
bundle on \(X\) with a continuous metric.
\begin{enumerate}
\item\label{itm:ch17-c1:section}
 The local functions \(\log\|s\|^{-1}\), for \(s\) an invertible section of
 \(L\) over an open \(U\), glue to a global section of the quotient sheaf
 \(C^0_X/\log|\cO_X^\times|\); if the metric is smooth they glue to a
 section of \(A^0_X/\log|\cO_X^\times|\), and if it is parlinear to a
 section of \(\PL_X/\log|\cO_X^\times|\).
\item\label{itm:ch17-c1:current}
 There is a unique current \(c_1(\overline L)\) of bidegree \((1,1)\) on
 \(X\) whose restriction to any such \(U\) is
 \(\dprimeddprime[\log\|s\|^{-1}]\).  It is \(\dprime\)-closed and
 \(\ddprime\)-closed, and its formation commutes with restriction.
\item\label{itm:ch17-c1:regularity}
 If the metric is smooth, \(c_1(\overline L)\) is (the current of) a smooth
 closed \((1,1)\)-form, the \emph{first Chern form}.  If the metric is psh,
 \(c_1(\overline L)\) is a positive current.
\item\label{itm:ch17-c1:seminorm}
 The same construction applies verbatim to a locally integrable
 semi-metric, that is, one for which \(\log\|s\|\) is locally integrable for
 every local invertible \(s\).
\end{enumerate}
\end{proposition}

\begin{proof}
\cref{itm:ch17-c1:section}.  If \(s,t\) are two invertible sections of \(L\)
over \(U\), then \(f=t/s\in\cO^\times(U)\) and
\(\log\|s\|^{-1}=\log\|t\|^{-1}+\log|f|\).  Thus the local data agree in the
quotient sheaf and glue.  For the refined statements, a smooth metric has
\(\log\|s\|\) smooth by \Cref{prop:ch17-smooth-line}, and a parlinear metric
has \(\log\|s\|\) parlinear by \Cref{def:ch17-parlinear-metric} together
with \Cref{lem:ch17-orthogonal-frames}\,\cref{itm:ch17-of:unique}.

\cref{itm:ch17-c1:current}.  The map \(u\mapsto\dprimeddprime[u]\) from
\(C^0_X\) to currents of bidegree \((1,1)\) kills the subsheaf
\(\log|\cO_X^\times|\): its sections are smooth---indeed
\(\log|f|\) with \(f\) invertible is smooth by
\Cref{prop:cld-local-normal-form}---and satisfy
\(\dprimeddprime\log|f|=0\), which is
\Cref{thm:cld-poincare-lelong} on the locus where \(f\) is invertible, that
is, where \(\divis(f)=0\).  Hence \(u\mapsto\dprimeddprime[u]\) factors
through \(C^0_X/\log|\cO_X^\times|\), and applying it to the global section
of \cref{itm:ch17-c1:section} gives a well-defined current, unique because
currents form a sheaf (\Cref{prop:ch14-currents-sheaf}).  Closedness is
\(\dprime\dprimeddprime=\ddprime\dprimeddprime=0\), which holds on
currents by \Cref{prop:current-differential-calculus}.

\cref{itm:ch17-c1:regularity}.  For a smooth metric the local expression is
the current of the smooth form \(\dprimeddprime\log\|s\|^{-1}\), and these
smooth forms glue by the same cocycle computation.  For a psh metric the
local weight \(\log\|s\|^{-1}\) is psh, so
\(\dprimeddprime[\log\|s\|^{-1}]\ge0\) by
\Cref{prop:cld-psh-positive-current}, and positivity of a current is local
by \Cref{prop:current-positivity-calculus}.

\cref{itm:ch17-c1:seminorm}.  The only point used above about \(\log\|s\|\)
was that it defines a current; local integrability supplies that, and the
cocycle relation is unchanged because \(\log|f|\) is continuous.
\end{proof}

\begin{proposition}
\label{prop:cld-c1-metric-change}
If \(\|\cdot\|=\rme^{-u}\|\cdot\|_0\) with \(u\) locally integrable, then
\[
       c_1(L,\|\cdot\|)=c_1(L,\|\cdot\|_0)+\dprimeddprime[u].
\]
\end{proposition}
\begin{proof}
In any frame the new weight is \(\phi_e=\phi_{0,e}+u\).  Apply
\(\dprimeddprime\) and use linearity of currents.
\end{proof}

\begin{corollary}
\label{cor:ch17-c1-functorial}
For metrized line bundles \(\overline L,\overline M\) on \(X\) and a
morphism \(f:Y\to X\) of boundaryless spaces,
\[
 c_1(\overline L\otimes\overline M)
  =c_1(\overline L)+c_1(\overline M),
 \qquad
 c_1(\overline L^{\vee})=-c_1(\overline L),
\]
and, when the metrics are smooth or when \(f\) is such that the relevant
pullback of currents is defined,
\(c_1(f^*\overline L)=f^*c_1(\overline L)\).  Thus \(c_1\) is a group
homomorphism \(\widehat{\Pic}(X)\to D^{1,1}(X)\) killing exactly the metrics
whose relative weights are pluriharmonic.
\end{corollary}

\begin{proof}
Additivity and duality follow from
\(\phi_{L\otimes M}=\phi_L+\phi_M\) and \(\phi_{L^\vee}=-\phi_L\) and
linearity of \(\dprimeddprime\) on currents.  For pullback, use
\(\phi_{f^*e}=f^*\phi_e\) from
\Cref{prop:ch17-metric-picard}\,\cref{itm:ch17-pic:pullback}: in the smooth
case both sides are the currents of pulled-back smooth forms and the
identity is \Cref{prop:cld-form-pullback}; in general it is the definition
of the pullback of a current where that is available, as delimited in
\Cref{prop:ch14-no-pullback} and the discussion following it.  The kernel
description is \Cref{prop:ch17-metric-picard}\,\cref{itm:ch17-pic:kernel}
together with \Cref{prop:cld-c1-metric-change}.
\end{proof}

\subsection{Poincar\'e--Lelong for a metrized line bundle}
\label{sec:ch17-pl-bundle}

The sign can be checked against a meromorphic section.  If \(s=fe\), then
\(-\log\|s\|=\phi_e-\log|f|\).  With the Poincar\'e--Lelong convention of
\cref{chap:cld-poincare-lelong},
\[
             \dprimeddprime[\log|f|]=\delta_{\operatorname{div}(f)},
\]
and therefore
\begin{equation}
 c_1(\overline L)
   =\dprimeddprime[-\log\|s\|]+\delta_{\operatorname{div}(s)}.
\label{eq:cld-pl-section-sign}
\end{equation}
For a nowhere-vanishing section the divisor term disappears.

Because \eqref{eq:cld-pl-section-sign} is used constantly---it is the only
way in which a global section enters a curvature computation---we state it
as a theorem with its hypotheses displayed.

\begin{theorem}
\label{thm:ch17-pl-metrized}
\index{Poincare--Lelong formula@Poincar\'e--Lelong formula!for metrized line bundles}
Let \(X\) be a boundaryless \(K\)-analytic space of pure dimension \(n\),
let \(L\) be a line bundle on \(X\) with a locally integrable semi-metric,
and let \(s\) be a regular meromorphic section of \(L\) on \(X\).  Then
\(\log\|s\|^{-1}\) is defined and locally integrable on the complement of
\(\Supp\divis(s)\), defines a current \([\log\|s\|^{-1}]\) on \(X\), and
\[
 \dprimeddprime[\log\|s\|^{-1}]+\delta_{\divis(s)}=c_1(\overline L).
\]
\end{theorem}

\begin{proof}
The assertion is local, and currents form a sheaf
(\Cref{prop:ch14-currents-sheaf}), so we may work on an open \(U\) carrying
an invertible section \(\varepsilon\) of \(L\).  Write \(s|_U=f\varepsilon\)
with \(f\) a regular meromorphic function on \(U\), which is possible and
unique because \(\varepsilon\) trivializes \(L|_U\); then
\(\divis(s)|_U=\divis(f)\) by definition of the divisor of a meromorphic
section, and
\[
 \log\|s\|^{-1}=\log\|\varepsilon\|^{-1}-\log|f|
\]
on the locus where \(f\) is invertible.  Both terms are locally integrable
there---the first by hypothesis, the second by
\Cref{prop:ch15-log-current}---so \(\log\|s\|^{-1}\) defines a current on
\(U\), namely \([\log\|\varepsilon\|^{-1}]-[\log|f|]\).  Applying
\(\dprimeddprime\) and using \Cref{thm:cld-poincare-lelong} for \(f\) gives
\[
 \dprimeddprime[\log\|s\|^{-1}]
 =\dprimeddprime[\log\|\varepsilon\|^{-1}]-\delta_{\divis(f)}
 =c_1(\overline L)|_U-\delta_{\divis(s)}|_U,
\]
which is the assertion on \(U\).  Regularity of \(s\) is what makes
\(\divis(s)\) a well-defined Cartier divisor of pure codimension one and
what allows \Cref{thm:cld-poincare-lelong} to be applied; without it the
left-hand side need not be defined, as \Cref{rem:ch15-nonregular} explains.
\end{proof}

\begin{corollary}
\label{cor:ch17-curve-total-mass}
\index{degree!and total curvature mass}
Let \(X\) be a proper \(K\)-analytic curve and let \(L\) be a line bundle on
\(X\) with a locally integrable semi-metric.  Then
\[
 \int_{X}c_1(\overline L)=\deg L.
\]
In particular the total mass of the curvature current does not depend on the
metric.
\end{corollary}

\begin{proof}
Suppose first that \(X\) is integral.  Then \(L\) has a nonzero global
meromorphic section \(s\), which is automatically regular, and
\Cref{thm:ch17-pl-metrized} gives
\[
 \langle c_1(\overline L),1\rangle
 =\langle\dprimeddprime[\log\|s\|^{-1}],1\rangle
  +\langle\delta_{\divis(s)},1\rangle.
\]
The first term vanishes: by \Cref{prop:current-differential-calculus} it
equals \(\langle[\log\|s\|^{-1}],\dprimeddprime1\rangle=0\), the pairing
being legitimate because \(X\) is compact and boundaryless, so that the
constant function \(1\) is an admissible test form and Green's formula
(\Cref{cor:analytic-stokes-green}) carries no boundary term.  For the second
term, write \(\divis(s)=\sum_in_i[P_i]\); then
\(\langle\delta_{P_i},1\rangle=[k(P_i):K]\) by
\Cref{prop:ch14-delta-Y-current} applied to a rigid point, so
\(\langle\delta_{\divis(s)},1\rangle=\sum_in_i[k(P_i):K]=\deg L\).

In general, let \(X_1,\dots,X_r\) be the irreducible components of \(X\)
with generic multiplicities \(m_i\), let \(Y=\coprod_iX_i\) and
\(p:Y\to X\) the canonical finite morphism.  Choose a finite open cover
\((U_j)\) of \(X\) with invertible sections \(s_j\) of \(L|_{U_j}\) and a
smooth partition of unity \((\lambda_j)\) subordinate to it
(\Cref{thm:cld-smooth-partition}).  Then
\[
 \langle c_1(\overline L),1\rangle
 =\sum_j\langle c_1(\overline L),\lambda_j\rangle
 =\sum_j\int_{U_j}\log\|s_j\|^{-1}\,\dprimeddprime\lambda_j.
\]
Let \(V_i\subseteq X\) be the complement of the union of the components
other than \(X_i\); the \(V_i\) are disjoint and the complement of their
union is a finite set of rigid points, near which the smooth form
\(\dprimeddprime\lambda_j\) vanishes, so each integral splits as a sum over
\(i\) of integrals over \(U_j\cap V_i\).  On \(p^{-1}(V_i)\subseteq X_i\)
the map \(p\) is an isomorphism onto \(V_{i,\mathrm{red}}\) away from
finitely many rigid points, and the generic multiplicity contributes the
factor \(m_i\) through the calibration
(\Cref{thm:current-fundamental-cycle}); hence
\[
 \int_{U_j\cap V_i}\log\|s_j\|^{-1}\dprimeddprime\lambda_j
 =m_i\int_{p^{-1}(U_j)\cap X_i}
   \log\|p^*s_j\|^{-1}\dprimeddprime(p^*\lambda_j).
\]
Summing over \(j\) and reversing the computation on each \(X_i\) gives
\[
 \langle c_1(\overline L),1\rangle
 =\sum_im_i\deg(L|_{X_i})
 =\deg L.
\]
Metric independence is now automatic, the right-hand side involving no
metric.  This is \cite[Corollary~18.1.7]{CLD25}.
\end{proof}

\begin{example}[Chern currents on a line bundle over the Berkovich line]
\label{ex:ch17-weil-curvature}
\index{Berkovich line!Chern currents on}
Let \(K\) be complete and nontrivially valued and let
\(X=\mathbf P^{1,\An}_K\) with coordinate \(T=x_1/x_0\).  Write
\(\zeta_{a,r}\) for the point of \(\mathbf A^{1,\An}_K\) attached to the
closed disc of centre \(a\) and radius \(r\), as in
\Cref{thm:berk-a1-classification}.  We compute the curvature of four
metrics on \(L=\cO(1)\), each time by writing the weight of the frame
\(s_{T_0}=x_0\) on the chart \(\{x_0\ne0\}\) as a function of
\(t=\log|T|\in[-\infty,+\infty]\) on the standard apartment, and then
applying the one-dimensional rule
\(\dprimeddprime[\max\{0,t\}]=\delta_{0}\) of
\Cref{ex:ch14-line-model-function}, in the form that a convex piecewise
affine weight has curvature \(\sum_v(\text{slope jump at }v)\,\delta_v\).

\emph{(a) The Weil metric.}  By \eqref{eq:ch17-weil-metric}, the weight is
\(\lambda(t)=\max\{0,t\}\).  Its slope jumps by \(1\) at \(t=0\), that is,
at \(\zeta_{0,1}\), and by \(0\) elsewhere; the branches at
\(T=0\) and \(T=\infty\), and every branch off the standard apartment, carry
constant weight.  Hence
\[
 c_1(\overline{\cO(1)}_{\mathrm{Weil}})=\delta_{\zeta_{0,1}},
 \qquad
 \int_Xc_1=1=\deg\cO(1),
\]
in agreement with \Cref{cor:ch17-curve-total-mass}.  For \(\cO(d)\) the
weight is \(d\lambda\) and \(c_1=d\,\delta_{\zeta_{0,1}}\), of mass \(d\).
This is the \(n=1\) case of \Cref{ex:ch17-support-projective} and reproves
\eqref{eq:p1-standard-measure}.

\emph{(b) A shifted formal metric.}  Fix \(\pi\in K^{\times}\) with
\(|\pi|<1\) and put \(\|\cdot\|'=\rme^{-u}\|\cdot\|_{\mathrm{Weil}}\) with
\(u=\min\{1,t^{+}\}\cdot v_K(\pi)^{-1}\cdot v_K(\pi)=\min\{1,t^{+}\}\); that
is, \(u\) is the hat function of \eqref{eq:metrics-p1-blowup-hat}.  The new
weight is \(\max\{0,t\}+\min\{1,t^+\}\); its slope is \(0\) for \(t\le0\),
\(2\) on \([0,1]\), and \(1\) for \(t\ge1\).  So
\[
 c_1(\overline L')=2\,\delta_{\zeta_{0,1}}-\delta_{\zeta_{0,|\pi|}},
\]
of total mass \(2-1=1=\deg\cO(1)\).  The measure is a signed measure: this
metric is not psh, and indeed the model line bundle
\(\rho^*\cO(1)\otimes\cO(E)\) is not relatively nef.  The example is
included precisely to show that \Cref{cor:ch17-curve-total-mass} is a
statement about total mass, not about positivity.

\emph{(c) The blowup metric of \Cref{ex:metrics-p1-blowup}.}  Take instead
\(\|\cdot\|''=\rme^{+u}\|\cdot\|_{\mathrm{Weil}}\), so that the weight is
\(\max\{0,t\}-\min\{1,t^+\}=\max\{-1,t-1\}+\) constant; concretely the
weight is \(\max\{-1,t\}\) after adding the constant \(1\), which does not
change the curvature.  The slope is \(0\) for \(t\le-1\) and \(1\) for
\(t\ge-1\), with a single jump of \(1\) at \(t=-1\).  Hence
\[
 c_1(\overline L'')=\delta_{\zeta_{0,|\pi|}},
\]
of mass \(1\).  This is exactly \Cref{ex:metrics-cl-p1-blowup}, computed
there from intersection numbers on the blown-up model
\(\mathcal X'\); the two computations agree, which is the one-dimensional
case of \Cref{thm:cld-recovers-cl}.

\emph{(d) A metric that is not parlinear.}  Let
\(\|\cdot\|_\theta=\rme^{-\theta u}\|\cdot\|_{\mathrm{Weil}}\) for
\(\theta\in\mathbf R\).  Its weight is \(\max\{0,t\}+\theta\min\{1,t^+\}\)
and
\[
 c_1(\overline L_\theta)
 =(1+\theta)\delta_{\zeta_{0,1}}-\theta\,\delta_{\zeta_{0,|\pi|}},
\]
of mass \(1\) for every \(\theta\).  For \(\theta\notin\mathbf Q\) this
metric is \(\mathbf R\)-parlinear but not \(\mathbf Q\)-parlinear, hence not
a model metric in the sense of \Cref{chap:metrics}; nonetheless
\(c_1\) is defined, is a signed atomic measure, and has the right mass.  The
metric is psh exactly for \(-1\le\theta\le0\), which is the condition that
both atoms be nonnegative---and, by \Cref{thm:ch17-psh-residual-nef}, the
condition that the residual line bundles be nef.
\end{example}

\begin{remark}[Reading the four computations]
\label{rem:ch17-four-computations}
All four weights in \Cref{ex:ch17-weil-curvature} are piecewise affine on
the segment \([-\infty,+\infty]\) and constant on every branch off it, so
all four curvatures are supported on the two points \(\zeta_{0,1}\) and
\(\zeta_{0,|\pi|}\).  What changes between them is only the distribution of
the mass \(1\).  This is the one-dimensional shadow of two later theorems:
that the total mass is an intersection number independent of the metric
(\Cref{thm:ch17-degree}), and that the individual atoms are the
special-fibre degrees of a model
(\Cref{thm:cld-formal-intersection}).  The reader who has followed the
present example has, in miniature, the entire content of
\Crefrange{sec:ch17-products}{sec:ch17-cl}.
\end{remark}

\section{Psh metrics}
\label{sec:ch17-psh}

\begin{definition}[CLD-psh metric]
A continuous metric is \emph{CLD-psh} on an open set if its local weights
are psh in the CLD sense.  For a smooth metric this is equivalent to
positivity of \(c_1(\overline L)\).  A metric is \emph{dpsh} if locally it
is a quotient of two psh metrics.
\end{definition}

Following \Cref{chap:cld-bedford-taylor} we also use the tropical variants:
a metric is \emph{tropically convex} (resp.\ \emph{tropically
\(d\)-convex}) if its local weights \(\log\|s\|^{-1}\) are, in the sense of
\Cref{prop:ch16-tropically-convex}.  The implications are
\[
 \begin{gathered}
 \text{tropically convex}\Rightarrow\text{psh}\Rightarrow\text{dpsh},\\
 \text{tropically convex}\Rightarrow\text{tropically }d\text{-convex}
 \Rightarrow\text{dpsh},
 \end{gathered}
\]
and, by \Cref{prop:ch16-dpsh-smooth}, every \(G\)-smooth metric---in
particular every parlinear metric---is tropically \(d\)-convex.  That last
point is worth pausing on: it is why the Bedford--Taylor machinery of
\Cref{chap:cld-bedford-taylor} applies to formal metrics at all, even
though those metrics are neither smooth nor, in general, psh.

\begin{proposition}
\label{prop:ch17-globally-generated-psh}
Let \(X\) be separated and almost good, and let \(L\) be a line bundle on
\(X\) with an integral parlinear metric, locally trivial for the usual
topology.  Let \(x\in X\) be such that some positive tensor power of the
residual bundle \(\widetilde L_x\) is generated by its global sections.
Then the metric of \(L\) is tropically convex, hence psh, on a
neighbourhood of \(x\).
\end{proposition}

\begin{proof}
Replacing \(L\) by a positive power---which changes neither the conclusion,
by homogeneity of convexity, nor the hypothesis---we may assume
\(\widetilde L_x\) globally generated.  By
\Cref{prop:ch17-residual-functorial}\,\cref{itm:ch17-res:iso} and the lifting
argument of \Cref{prop:ch17-residual-welldefined}, global generation
translates into the following analytic statement: if \(s\) is an invertible
section of \(L\) near \(x\), there are finitely many
\(a_1,\dots,a_m\in\cO_{X,x}^\times\) with
\[
 \|s\|^{-1}=\max\{|a_1|,\dots,|a_m|\}
\]
on a neighbourhood of \(x\).  Indeed, global sections
\(\sigma_1,\dots,\sigma_m\) of \(\widetilde L_x\) generating it lift, after
shrinking, to sections whose reductions are the \(\sigma_j\), and the
generation statement says exactly that the maximum of the corresponding
absolute values computes the norm.  Now
\(\log\max_j|a_j|\) is the pullback of the convex function
\(\max_j t_j\) under the moment map \((a_1,\dots,a_m)\), so it is tropically
convex by \Cref{prop:ch16-tropically-convex}; hence the metric is
tropically convex near \(x\), and a tropically convex function is psh.
\end{proof}

\begin{theorem}
\label{thm:ch17-psh-residual-nef}
\index{nef!residual line bundle}
Let \(X\) be boundaryless and separated, let \(L\) be a line bundle on \(X\)
with an integral parlinear metric, and let \(x\in X\).  The following are
equivalent:
\begin{enumerate}
\item\label{itm:ch17-nef:psh} the metric of \(L\) is psh near \(x\);
\item\label{itm:ch17-nef:pos} the current \(c_1(\overline L)\) is positive
 near \(x\);
\item\label{itm:ch17-nef:nef} the residual line bundle \(\widetilde L_x\)
 is numerically effective.
\end{enumerate}
\end{theorem}

\begin{proof}
\cref{itm:ch17-nef:psh}\(\Rightarrow\)\cref{itm:ch17-nef:pos} is
\Cref{prop:ch17-c1-construction}\,\cref{itm:ch17-c1:regularity}.

\cref{itm:ch17-nef:pos}\(\Leftrightarrow\)\cref{itm:ch17-nef:nef} is the
positivity half of \Cref{thm:ch17-formal-positivity} below, applied with
\(d=1\); we have separated the statements only because the general case
requires the residual intersection theory of \Cref{sec:ch17-formal}.

\cref{itm:ch17-nef:nef}\(\Rightarrow\)\cref{itm:ch17-nef:psh}.  Let
\((V,M)\) be a premodel of \(\widetilde L_x\); since \(X\) is boundaryless,
\(V\) is proper, and by Chow's lemma we may assume \(V\) projective.  Choose
an ample \(M'\) on \(V\).  There is a compact neighbourhood \(U\) of \(x\)
and a strictly parlinear metric on the trivial bundle \(\cO_U\) whose
residual bundle at \(x\) is \(M'\); write \(h=\log\|1\|^{-1}\) for its
weight, a tropically convex, hence psh, function by
\Cref{prop:ch17-globally-generated-psh} applied to a power of \(M'\).

For \(\varepsilon>0\) let \(\|\cdot\|_\varepsilon
=\rme^{-\varepsilon h}\|\cdot\|\) on \(L|_U\).  As
\(\varepsilon\downarrow0\) these converge uniformly on \(U\) to the given
metric, so by \Cref{prop:ch15-uniform-limit} it suffices to show that
\(\|\cdot\|_\varepsilon\) is psh for rational \(\varepsilon\).  Write
\(\varepsilon=m/N\) with \(m,N\ge1\).  Then
\((\|\cdot\|_\varepsilon)^{\otimes N}\) is a strictly parlinear metric on
\(L^{\otimes N}|_U\) whose residual bundle at \(x\) has premodel
\((V,M^{\otimes N}\otimes(M')^{\otimes m})\).  Since \(M\) is nef and
\(M'\) ample, the Nakai--Moishezon criterion makes
\(M^{\otimes N}\otimes(M')^{\otimes m}\) ample, in particular globally
generated after a further power; \Cref{prop:ch17-globally-generated-psh}
then gives pshness of \((\|\cdot\|_\varepsilon)^{\otimes N}\) near \(x\),
hence of \(\|\cdot\|_\varepsilon\).  This is
\cite[Proposition~17.8.4 and Corollary~18.8.8]{CLD25}.
\end{proof}

\subsection{Global psh metrics, and the gap}
\label{sec:ch17-global-psh}

A continuous metric on \(L\) is \emph{globally psh} if there is a sequence
of smooth psh metrics \(\|\cdot\|_j\) on \(L\) with
\(\log(\|\cdot\|/\|\cdot\|_j)\to0\) uniformly on \(X\); it is
\emph{globally dpsh} if \(L\) is a quotient of two globally psh metrized
line bundles.  This is the smooth-approximation definition of
\cite[\S17.8.5]{CLD25}.  It must be distinguished from uniform
approximation by nef \(\mathbf Q\)-model metrics on a fixed line bundle,
as used in \cref{chap:metrics}.  The implications below have explicit
positivity hypotheses; the two approximation classes are not identified
by definition.

\begin{theorem}
\label{thm:ch17-weil-globally-psh}
On \(\mathbf P^{n,\An}_K\), the line bundle \(\cO(1)\) with its canonical
formal metric is tropically convex and globally psh.  Consequently:
\begin{enumerate}
\item\label{itm:ch17-gpsh:semiample}
 if \(X\) is a proper strictly \(K\)-analytic space, \(\mathfrak X\) a
 proper \(K^\circ\)-model, \(\mathfrak L\) a semiample line bundle on
 \(\mathfrak X\), and \(L\) its generic fibre with the formal metric, then
 \(\overline L\) is globally psh;
\item\label{itm:ch17-gpsh:dpsh}
 if \(X\) is the analytification of a projective \(K\)-variety and
 \(\overline L\) carries a parlinear metric, then \(\overline L\) is
 globally dpsh.
\end{enumerate}
\end{theorem}

\begin{proof}
Every residual bundle of \(\cO(1)\) admits, by
\Cref{lem:ch17-residual-formal}, a premodel of the form
\((V,\cO_{\mathbf P^n_{\widetilde K}}(1)|_V)\) with
\(V\subseteq\mathbf P^n_{\widetilde K}\) closed, hence semiample; so the
Weil metric is tropically convex by
\Cref{prop:ch17-globally-generated-psh}.

For global pshness, let \(M_\eta:\mathbf R^{p}\to\mathbf R\)
(\(\eta>0\)) be Demailly's regularized maximum: smooth, convex,
nondecreasing in each variable, with
\(\max t_i\le M_\eta(t)\le\max t_i+\eta\), invariant under adding a constant
to all variables, and independent of a variable that is smaller than the
others by more than \(2\eta\).  Set
\[
 \log\|s_P\|_\eta^{-1}([x_0:\cdots:x_n])
  =-\log|P(x)|+M_\eta(\log|x_0|,\dots,\log|x_n|)
\]
for \(P\) homogeneous of degree one.  Translation invariance makes this
well defined, and smoothness of \(M_\eta\) makes \(\|\cdot\|_\eta\) a smooth
metric.  Comparing with \eqref{eq:ch17-weil-metric},
\(\log\|s_P\|_\eta^{-1}-\log\|s_P\|^{-1}\in[0,\eta]\), so
\(\|\cdot\|_\eta\to\|\cdot\|\) uniformly.  Pshness of \(\|\cdot\|_\eta\) is
local: at a point with \(x_0,\dots,x_p\ne0\) and \(x_{p+1}=\cdots=x_n=0\),
take \(P=T_0\) and \(f_j=T_j/T_0\); on a small enough neighbourhood the
last \(n-p\) arguments are smaller than the others by more than \(2\eta\)
once \(\eta<1\), so
\(\log\|s_{T_0}\|_\eta^{-1}=M_\eta(0,\log|f_1|,\dots,\log|f_p|)\),
the pullback of a smooth convex function under a moment map, which is
smooth psh.

\cref{itm:ch17-gpsh:semiample}.  Semiampleness gives \(d\ge1\) and
\(p:\mathfrak X\to\mathbf P^N_{K^\circ}\) with
\(p^*\cO(1)\simeq\mathfrak L^{\otimes d}\); the induced metric on
\(L^{\otimes d}\) is then \(p^*\) of the Weil metric, so it is globally psh,
and hence so is \(\overline L\).

\cref{itm:ch17-gpsh:dpsh}.  A parlinear metric on \(L^{\An}\) comes from a
formal model \((\mathfrak X,\mathfrak L)\) by
\Cref{thm:cld-formal-pl-comparison}, and \(\mathfrak X\) may be taken
projective over \(K^\circ\); formal GAGA then makes \(\mathfrak L\)
algebraic, so \(\mathfrak L\simeq\mathfrak L_1\otimes\mathfrak L_2^{-1}\)
with \(\mathfrak L_i\) ample.  By \cref{itm:ch17-gpsh:semiample} each
\(\overline L_i\) is globally psh, and the isomorphism is an isometry
because the formal metric of a tensor product is the tensor product of the
formal metrics.  This is \cite[Corollaries~17.8.7 and 17.8.9]{CLD25}.
\end{proof}

For a general continuous metric, positivity of the current is a necessary
condition in all comparison results below, but one must not silently replace
CLD's approximation-defined psh class by ``positive curvature current'' in
every singularity class.  The continuous smooth-proper-curve exception is
\Cref{thm:p063-current-psh-ordinary-regularization}; the distinctions in
general are made explicit in
\Cref{chap:delta-forms}.

For a smooth proper curve, the comparison with model metrics separates into
two cases.  Positive degree provides a strictly positive curvature margin
on a finite skeleton, which permits rational-slope approximation while
preserving the metric cocycle.  In degree zero that margin disappears:
approximation is equivalent to the existence of one nef rational model
metric on the fixed line bundle.  The following proof treats an arbitrary
complete nontrivially valued ground field and keeps this distinction
throughout.

\begin{theorem}
\label{thm:p081-corrected-curve-comparison}
Let $K$ be a complete non-Archimedean field with nontrivial absolute
value, let $C$ be a smooth proper geometrically connected algebraic
$K$-curve, and let $L$ be a line bundle on $C$.  Write the same
symbols for their analytifications.  Suppose that a continuous metric
$\lVert\cdot\rVert$ on $L$ has CLD-psh weights
$-\log\lVert s\rVert$ in local invertible frames.  Then its Thuillier
curvature measure $\mu$ is positive and has total mass $\deg L$.
In particular, $\deg L\geq0$.
\begin{enumerate}
\item\label{itm:p081-corrected-positive}
  If $\deg L>0$, the metric is a uniform limit of semipositive
  $\mathbf Q$-formal metrics over $K$.  The approximants may in fact
  be chosen to be algebraic $\mathbf Q$-model metrics from relatively
  ample models.  Consequently the metric lies in the uniform
  Fubini--Study closure.  These approximation assertions remain true
  after every complete scalar extension of $K$.
\item\label{itm:p081-corrected-zero}
  If $\deg L=0$, the metric is harmonic.  In this case it is a uniform
  limit of semipositive $\mathbf Q$-formal metrics if and only if
  $L$ admits at least one semipositive $\mathbf Q$-formal metric over
  $K$.  Whenever this condition holds, the approximants may be taken
  numerically trivial on the special fibre and may be obtained from
  one fixed such model metric by constants in
  $\mathbf Q(-\log|K^\times|)$.  Their pullbacks give the same
  approximation after every complete scalar extension.
\end{enumerate}
No model-existence assertion is made for an arbitrary degree-zero line
bundle.
\end{theorem}

% Unverified replacement proof.  The real-PL sheaf is never identified
% with the real span of global rational model functions in this argument.

\begin{proof}
We use additive weights.  The distance between two continuous metrics is
\[
 d_\infty(\lVert\cdot\rVert_1,\lVert\cdot\rVert_2)
  :=\sup_{C^{\mathrm{an}}}
       \left|\log\frac{\lVert\cdot\rVert_1}
                       {\lVert\cdot\rVert_2}\right|;
\]
the quotient is independent of the frame.

\emph{1. Curvature and degree.}
By \cite[\S1.2 and \S4.1.2]{GR25Classical}, CLD-psh weights are
Thuillier-subharmonic on a smooth strict analytic curve over any complete
nontrivially valued field.  If \(s=f s'\), then
\[
 -\log\lVert s\rVert=-\log\lVert s'\rVert-\log|f|,
\]
and \(\log|f|\) is harmonic.  Thus the weights define one global
subharmonic metric and a positive curvature measure \(\mu\).
The metrized Poincar\'e--Lelong formula and its degree consequence
\cite[Theorem~18.1.6 and Corollary~18.1.7]{CLD25} give
\[
 \mu=\mathrm d'\mathrm d''[-\log\lVert\sigma\rVert]
          +\delta_{\operatorname{div}(\sigma)},\qquad
 \mu(C^{\mathrm{an}})=\deg L
\]
for a nonzero meromorphic section \(\sigma\).  Closed points and component
degrees are counted with their residue-field degrees.  This proves the
first assertion.  If the degree is zero, the positive measure vanishes.
For a local continuous weight \(h\), both \(\mathrm d'\mathrm d''h\)
and \(\mathrm d'\mathrm d''(-h)\) are then zero.  Thuillier's current
criterion \cite[Theorem~3.4.12]{Thu05} says that \(h\) and \(-h\) are
subharmonic.  By \cite[\S3.12]{GR25Classical}, this is exactly
harmonicity, so the metric is harmonic.

\emph{2. Two elementary facts about a finite metric graph.}
Let \(G\) be a finite connected metric graph.  For a continuous function
affine on each edge, use the Laplacian convention
\[
 (\Delta_G f)(\{v\})=\sum_{\nu\in T_vG}\partial_\nu f.
\]
A break in an oriented edge with left slope \(r_-\) and right slope
\(r_+\) contributes \(r_+-r_-\).  Loops contribute both incident
directions.  In this convention the energy identity is
\[
 \int_G |f'|^2=-\int_G f\,\Delta_G f.
\]
It follows by integration by parts separately on the finitely many
edges.  In particular, the graph Laplacian on functions affine between
specified vertices is a symmetric matrix whose kernel consists exactly
of the constants: zero energy forces equality along every edge, and
connectedness then forces constancy.  Its image is the space of
vertex-mass vectors of sum zero.  Hence any finite atomic signed measure
of total mass zero has a piecewise affine potential, unique modulo
constants, after its support has been inserted among the vertices.

We will also use an explicit continuity estimate.  If
\(\Delta_G g=m(\delta_y-\delta_x)\), \(m\geq0\), and \(g(p)=0\) at a
fixed point \(p\), then
\begin{equation}
\label{eq:p081-corrected-transport-bound}
                 \lVert g\rVert_\infty\leq m\,d_G(x,y).
\end{equation}
Indeed, the maximum principle on an edge and at a vertex of zero
Laplacian shows that the maximum and minimum occur at \(x\) and \(y\),
respectively, unless \(g\) is constant.  The energy identity yields
\(\int_G|g'|^2=m\,\operatorname{osc}(g)\).
Along a shortest path from \(x\) to \(y\), the Cauchy--Schwarz inequality
therefore gives
\[
 \operatorname{osc}(g)^2
 \leq d_G(x,y)\int_G|g'|^2
 =m\,d_G(x,y)\operatorname{osc}(g).
\]
This proves the estimate, since the normalization puts zero between
the extreme values.  Potentials transporting finitely many masses
satisfy the sum of the corresponding bounds.

\emph{3. Reduction to a skeletal metric over an algebraically closed
field.}
For the moment suppose that \(K\) is algebraically closed and that
\(d:=\deg L>0\).  By \cite[Theorem~3.4.19]{Thu05}, a continuous
subharmonic metric is the increasing envelope, in norm notation, of
global lisse subharmonic metrics on the same line bundle.  Here
\emph{lisse} is Thuillier's skeletal notion, not ordinary CLD
smoothness.  The proof of that theorem constructs the metrics globally
and checks their transition functions.  The approximation is uniform:
the logarithms of the quotients of the target norm by the approximating
norms are nonnegative continuous functions decreasing to zero on the
compact curve.  For each positive accuracy their sublevel sets form an
increasing open cover.  A finite subcover and directedness give uniform
convergence.  A choice for each accuracy \(2^{-j}\) gives a sequence.
It suffices to approximate one lisse subharmonic metric.

Choose a formal reference metric \(\lVert\cdot\rVert_0\) on \(L\).
Existence and the formal--PL correspondence are
\cite[\S2.2 and Proposition~2.10]{GM19}.  After a common semistable
refinement, the reference curvature \(\mu_0\) is supported on the
type-II vertices of a finite skeleton, and the relative weight of the
lisse metric has the form \(\Phi\circ\tau_G\), with \(\Phi\) piecewise
affine on a finite metric graph \(G\).  The required factorization and
common refinement are \cite[Proposition~3.2.4 and Lemma~2.2.22]{Thu05}.
On the graph its curvature is
\[
                    \mu_\ell=\mu_0+\Delta_G\Phi.
\]
There may initially be type-III breakpoints.  For each such atom
\(m_x\delta_x\) of the positive measure \(\mu_\ell\), move the mass to
a nearby type-II point \(y_x\) on the same edge.  Type-II points are
dense there because \(\Gamma:=-\log|K^\times|\) is nonzero and divisible.
Solve
\[
 \Delta_G g=\sum_xm_x(\delta_{y_x}-\delta_x),\qquad g(p)=0.
\]
By \eqref{eq:p081-corrected-transport-bound}, the norm of \(g\) is at most
\(\sum_xm_xd_G(x,y_x)\), so this change can be made arbitrarily small.
The new metric, with relative weight
\((\Phi+g)\circ\tau_G\), has positive curvature supported only at
type-II points.  An old type-III vertex has valence two and zero
remaining Laplacian, so it is no longer a breakpoint.  Subdivide at
the remaining type-II breakpoints; by \cite[Lemma~2.2.22]{Thu05} this is
a strict semistable refinement over \(K\).

We have therefore reduced to a finite semistable skeleton \(G\), with
vertex set \(V\), edge lengths in \(\Gamma\), and a function \(\Phi\)
affine on every edge, such that
\begin{equation}
\label{eq:p081-corrected-vertex-curvature}
 \mu_0+\Delta_G\Phi=\sum_{v\in V}b_v\delta_v,\qquad
 b_v\geq0,\qquad\sum_{v\in V}b_v=d>0.
\end{equation}
The reference model is included in the refinement.  We have not
identified \(\Phi\) with a real linear combination of global model
functions.

\emph{4. Rational-slope approximation with positive vertex margin.}
Let
\[
        \nu=\frac{d}{|V|}\sum_{v\in V}\delta_v.
\]
The first graph fact gives a potential \(h\) with
\(\Delta_G h=\nu-\sum_v b_v\delta_v\).  For \(0<\lambda<1\), put
\(\Phi_\lambda=\Phi+\lambda h\).  It is still affine on the edges, and
\[
 \mu_0+\Delta_G\Phi_\lambda
   =\sum_{v\in V}\left((1-\lambda)b_v+
                         \lambda\frac{d}{|V|}\right)\delta_v.
\]
Every vertex thus has curvature at least
\(\sigma:=\lambda d/|V|>0\), while
\(\lVert\Phi_\lambda-\Phi\rVert_\infty
 \leq\lambda\lVert h\rVert_\infty\).

We describe explicitly a rational-slope approximation \(\Psi\) of
\(\Phi_\lambda\).  If the graph is a point, approximate its one value by
an element of \(\Gamma\), and there is nothing further to prove.
Otherwise choose
\[
 0<\delta<\frac{\sigma}{2\max(1,\max_{v\in V}|T_vG|)}.
\]
Choose \(q_v\in\Gamma\), close enough to \(\Phi_\lambda(v)\) that, on
each oriented edge \(e\) from \(v\) to \(w\), of length \(\ell_e\), the
number
\[
 s'_e=\frac{q_w-q_v}{\ell_e}
\]
differs from the slope \(s_e\) of \(\Phi_\lambda\) by less than
\(\delta/2\).  This is possible by density of \(\Gamma\) and finiteness
of the graph.  Choose rational numbers
\[
 r_e^-<s'_e<r_e^+,\qquad |r_e^\pm-s'_e|<\delta/2,
\]
and set
\begin{equation}
\label{eq:p081-corrected-breakpoint}
 t_e=\frac{r_e^+\ell_e-(q_w-q_v)}{r_e^+-r_e^-}.
\end{equation}
Then \(0<t_e<\ell_e\).  Divisibility of \(\Gamma\) implies that
\(t_e\in\Gamma\).  On this edge define \(\Psi\) to start at \(q_v\),
have slope \(r_e^-\) until \(t_e\), and have slope \(r_e^+\)
afterwards.  Formula \eqref{eq:p081-corrected-breakpoint} makes its final
value exactly \(q_w\).  The construction also works on a loop, for
which \(q_w=q_v\).  It therefore gives one continuous global function
on the graph, not independently chosen edgewise weights.

The new interior breakpoint has positive Laplacian
\(r_e^+-r_e^-\).  At an original vertex each outgoing slope has changed
by less than \(\delta\).  Thus its curvature is at least
\(\sigma-|T_vG|\delta>\sigma/2\).
It follows that
\[
                         \mu_0+\Delta_G\Psi\geq0.
\]
All vertex and breakpoint values belong to \(\Gamma\), all subdivision
lengths belong to \(\Gamma\), and all slopes are rational.  After
multiplication by one positive integer, \(\Psi\) is integral
\(\Gamma\)-affine.  By \cite[Proposition~3.5(c)]{GJKM19},
\(\Psi\circ\tau_G\) is a \(\mathbf Q\)-model function.  The model slope
criterion \cite[Proposition~3.7]{GJKM19} shows that
\(\lVert\cdot\rVert_0\mathrm e^{-\Psi\circ\tau_G}\) is semipositive.
In that criterion the residue weights are one in the present
algebraically closed setting.

For the approximation bound, put
\(\eta=\max_v|q_v-\Phi_\lambda(v)|\).  The slopes on either piece of
an edge differ from \(s_e\) by less than \(\delta\), so
\[
 \lVert\Psi-\Phi_\lambda\rVert_\infty
       \leq\eta+\delta\max_e\ell_e.
\]
First choose \(\lambda\) small, then \(\delta\) small subject to the
vertex-margin inequality, and finally the \(q_v\) sufficiently close.
The resulting semipositive model metrics approximate the skeletal
metric arbitrarily well.  Together with Step 3 this proves uniform
\(\mathbf Q\)-formal approximation for positive degree over an
algebraically closed field.  The use of \(d>0\) occurs precisely in the
strict vertex margin; it cannot be suppressed at degree zero.

\emph{5. Descent of the positive-degree approximation to \(K\).}
Return to an arbitrary \(K\), still with \(\deg L>0\), and put
\(F=\widehat{K^{\mathrm a}}\).  Smoothness and geometric connectedness
imply that \(C\) is geometrically integral.  Its base change to \(F\)
is a smooth proper connected curve.  Scalar extension preserves
CLD-psh weights, including their transfinite stages, by
\cite[Proposition~8.1.5 and \S11.1.9]{CLD25}.  Step 4 therefore supplies
semipositive \(\mathbf Q\)-formal approximants over \(F\).

Choose a formal reference metric over \(K\), and write
\[
 \theta=c_1(L,\lVert\cdot\rVert_0),\qquad
 u=-\log\frac{\lVert\cdot\rVert}{\lVert\cdot\rVert_0}.
\]
The metric--function dictionary is \cite[\S4.7]{GM19}.  Choose upstairs
\(\theta_F\)-psh model functions \(\varphi_i^F\) with
\(\|\varphi_i^F-u_F\|_\infty<i^{-1}\).
By \cite[Proposition~2.18(b)]{GM19}, each upstairs
\(\mathbf Q\)-PL metric is defined over a finite subextension
\(K_i/K\) of \(F/K\).  It is semipositive over \(K_i\): after clearing
denominators choose a representing formal model using
\cite[Proposition~2.10]{GM19}; its base change is nef by
\cite[Proposition~3.5]{GM19}, and nefness descends by
\cite[Lemma~3.3]{GM19}.  Denote the descended model function by
\(\varphi_i^{K_i}\).

Choose
\[
 c_i\in\mathbf Q(-\log|K_i^\times|),\qquad
                     i^{-1}\leq c_i\leq2i^{-1},
\]
and put \(\psi_i^{K_i}=\varphi_i^{K_i}-c_i\).  Such constants exist
because a nonzero divisible subgroup of \(\mathbf R\) is dense.
They are \(\mathbf Q\)-model constants, so the shift does not change
semipositivity.  Upstairs,
\[
                         u_F-3i^{-1}\leq\psi_i^F\leq u_F.
\]
Take a finite normal extension \(N_i/K\) in \(K^{\mathrm a}\)
containing \(K_i\), for example the splitting field of finitely many
generating minimal polynomials.  Separability is not required.
Write \(q_i:C_{N_i}^{\mathrm{an}}\to C^{\mathrm{an}}\).
Surjectivity of scalar-extension projections
\cite[\S11.1.9]{CLD25} gives
\[
 q_i^*u-3i^{-1}\leq\psi_i^{N_i}\leq q_i^*u.
\]
Consequently
\[
 q_i^*u-3i^{-1}
 \leq P_{q_i^*\theta}(q_i^*u)\leq q_i^*u.
\]
For the finite normal extension, and only for that extension,
\cite[Proposition~2.11]{GJKM19} identifies the middle expression with
\(q_i^*P_\theta(u)\).  Surjectivity of \(q_i\) and passage to the limit
give
\begin{equation}
\label{eq:p081-corrected-envelope}
                              P_\theta(u)=u.
\end{equation}
Here the class \(c_1(L)\) is ample, since \(\deg L>0\).
The ample curve envelope theorem
\cite[Theorem~1.1 and Proposition~2.9(viii)]{GJKM19}, together with
\eqref{eq:p081-corrected-envelope}, gives uniform approximation by
\(\theta\)-psh model functions over \(K\).  Their metrics are the
required semipositive \(\mathbf Q\)-formal metrics on the fixed
line bundle \(L\).  This use of the envelope theorem is exclusively in
positive degree; no nonemptiness assertion is inferred from a
degree-zero numerical class.

\emph{6. Relatively ample models and scalar extensions.}
By \cite[Theorem~1.1]{GM19}, semipositive formal and algebraic model
metrics agree on a proper algebraic variety.  Represent one approximant
by a nef \(\mathbf Q\)-model \(\mathcal N_j\).  Proposition~4.11 of
\cite{GM19} provides a dominating algebraic model on which a positive
tensor power of \(L\) has an ample extension.  Divide by that tensor
power and call the resulting ample \(\mathbf Q\)-model \(\mathcal A_j\).
The pullback of \(\mathcal N_j\) is nef.

For a sufficiently small rational \(0<\varepsilon_j<1\), the
\(\mathbf Q\)-model
\[
            (1-\varepsilon_j)\mathcal N_j+
                           \varepsilon_j\mathcal A_j
\]
has metric at distance less than \(j^{-1}\) from the original
approximant.  Its degree on each special-fibre component is strictly
positive.  Hence it is relatively ample by the fibrewise criterion
used in \cite[Proposition~4.13]{GM19}.  These relatively ample metrics
still converge to the target.  Clear denominators and take a tensor
power that is relatively generated.  If \(s_0,\ldots,s_N\) are model
sections generating that power, its model metric is characterized by
\[
                           \max_\alpha\|s_\alpha(x)\|=1.
\]
This is the Fubini--Study metric of the generated \(K^\circ\)-lattice.
Taking the rational root proves the Fubini--Study assertion.

Under any complete extension \(E/K\), model metrics remain model
metrics and nefness is preserved
\cite[Proposition~2.18(a) and Lemma~3.3]{GM19}; relative ampleness is
also preserved by base change.  Surjectivity of
\(C_E^{\mathrm{an}}\to C^{\mathrm{an}}\) shows that pullback preserves
the uniform metric distance.  The pulled-back sequence therefore
proves all the extension assertions in
\cref{itm:p081-corrected-positive}.

\emph{7. The exact degree-zero criterion.}
Suppose that \(\deg L=0\).
If a sequence of semipositive \(\mathbf Q\)-formal metrics exists,
then any member supplies the model metric required in the statement.
Conversely, let \(\lVert\cdot\rVert_*\) be one such metric.
The model positivity criterion and its curve subharmonicity
interpretation \cite[Proposition~3.5]{GM19},
\cite[Proposition~5.1]{GR25PL} give a positive curvature measure.
Its mass is zero, so it too is harmonic.  Hence
\[
                  a=-\log
                    \frac{\lVert\cdot\rVert}{\lVert\cdot\rVert_*}
\]
is a global continuous harmonic function on the proper connected curve.
It is constant: it attains a maximum by compactness, and the harmonic
maximum principle makes its maximum locus open as well as closed.
Choose constants
\(a_j\in\mathbf Q(-\log|K^\times|)\) tending to \(a\).
The metrics
\(\lVert\cdot\rVert_*\mathrm e^{-a_j}\) are semipositive
\(\mathbf Q\)-formal metrics converging uniformly to the given one.

They are numerically trivial on the special fibre.  To see this without
assuming a discrete valuation or reduced special fibre, choose an
algebraic nef model of a positive tensor power, say \(\mathcal L\)
extending \(L^{\otimes m}\).  For \(0\ne\rho\in K^{\circ\circ}\), the
rank-one divisorial degree formula of \cite[\S4.8]{GM19} gives
\[
 (-\log|\rho|)m\deg L
      =\sum_Y\lambda_Y(\rho)\deg_Y(\mathcal L),\qquad
                            \lambda_Y(\rho)>0.
\]
Here \(Y\) runs over the reduced irreducible special-fibre components;
the coefficients come from
\(\operatorname{cyc}(\operatorname{div}(\rho))\), not an assumption
that the special fibre is itself a Cartier divisor.  The left side
vanishes.  Nefness makes each component degree nonnegative, so each is
zero.  Every proper integral curve in the special fibre is one of
these components, proving numerical triviality.  The constant shifts
preserve these degrees.  The same scalar-extension argument as in
Step 6 completes \cref{itm:p081-corrected-zero}.
\end{proof}


The model-existence condition in degree zero is essential.  The next example
has two edges on a Tate skeleton whose lengths have an irrational ratio.
Zero curvature forces the correction to be affine along each of these
edges, but its slopes cannot then be rational.  Adding vertices or attached
trees cannot change this obstruction.  At the same time, the canonical
cubist metric is locally smooth and harmonic in the ordinary CLD sense.

% The target is fixed-line-bundle Q-formal nef approximation, not an enlarged CLD cone.
\begin{theorem}
\label{thm:p081-irrational-tate-period-obstruction}
Let $K$ be an algebraically closed complete non-Archimedean field with
nontrivial real-valued absolute value, put $v=-\log|\cdot|$, and choose
$q,u\in K^\times$ such that
\[
  \ell:=v(q)>0,\qquad 0<a:=v(u)<\ell,
  \qquad a/\ell\notin\mathbf Q.
\]
Let $E_q/K$ be the Tate elliptic curve, let
$\pi:\mathbf G_{m,K}^{\mathrm{an}}\to E_q^{\mathrm{an}}$
be its uniformization with kernel $q^{\mathbf Z}$, let
$O=\pi(1)$, $P=\pi(u)$, and set
$L=\mathcal O_{E_q}(P-O)$.
Then the following assertions hold.
\begin{enumerate}
\item
$E_q$ is smooth, projective and geometrically connected, and $\deg L=0$.
For every integer $m>0$, there is no nef formal model of
$L^{\otimes m}$ on any proper admissible formal model of
$E_q^{\mathrm{an}}$. Here nef means nonnegative degree on every
irreducible component of the special fibre.
Consequently $L$ admits neither a semipositive
$\mathbf Q$-formal metric nor a numerically trivial
$\mathbf Q$-formal metric.
\item
After choosing a rigidification at the origin, the canonical cubist
metric $\|\cdot\|_{\mathrm{cub}}$ on $L$ is continuous and locally
CLD-psh, and its CLD curvature current satisfies
\[
 c_1(L,\|\cdot\|_{\mathrm{cub}})=0.
\]
In fact its local weights in the Tate uniformization are real affine
functions of $-\log|z|$, up to logarithms of analytic units, and are
therefore locally CLD-smooth and harmonic.
\item
The metric $\|\cdot\|_{\mathrm{cub}}$ is not a uniform limit of
semipositive $\mathbf Q$-formal metrics on this fixed line bundle
$L$. Thus local CLD-pshness does not imply membership in the
fixed-$L$ uniform closure of nef $\mathbf Q$-formal metrics over
arbitrary complete nontrivially valued fields, even when the field is
algebraically closed and the curve is smooth and projective.
\end{enumerate}
The hypotheses occur, for example, by choosing an irrational
$\alpha\in(0,1)$, equipping $\mathbf Q_p(T)$ with the weighted
Gauss absolute value extending $|p|=p^{-1}$ and satisfying
$|T|=p^{-\alpha}$, and taking the completed algebraic closure of its
completion. One may then take $q=p$, $u=T$.
\end{theorem}

% Source abbreviations and precise versions are specified in ../SOURCE_INTERFACE.md.
\begin{proof}
We use the model-metric convention of Gubler--Martin \cite[Introduction and Sections 2--4]{GM19}: a $\mathbf Q$-formal metric becomes a
formal metric after a positive integral tensor power, and semipositivity
means that a representing formal line bundle is nef on the special fibre.
All slopes below are measured using the intrinsic logarithmic annulus
length, not an arbitrarily rescaled coordinate on an edge.

\emph{Step 1: the field, the Tate curve, and its marked skeleton.}
For the field example in the statement, the weighted Gauss norm on
$\mathbf Q_p[T]$ is
\[
 \left|\sum_i c_iT^i\right|
   =\max_i |c_i|_p p^{-\alpha i}.
\]
Distinct nonzero monomials have distinct valuations because
$\alpha\notin\mathbf Q$; hence the term of largest norm is unique,
and the same observation proves multiplicativity. Extend the norm to
$\mathbf Q_p(T)$, complete, and take a completed algebraic closure.
The resulting complete algebraically closed non-Archimedean field
contains $p,T$ with valuations $\log p,\alpha\log p$.

Tate's construction applies to every complete nontrivially valued field
and every $q$ with $0<|q|<1$: the Weierstrass equation and
nonvanishing discriminant in \cite[equations (13)--(18)]{Tate95} give the
elliptic curve, and \cite[Theorem 1]{Tate95} gives the multiplicative
uniformization on rational points. Its analytic uniformization and
skeletal description are \cite[Remark 4.24]{BPR14}. In particular, the minimal
skeleton is
\[
 \Sigma=\mathbf R/\ell\mathbf Z,
\]
with coordinate induced by $v(z)$. Its retraction $\tau$ sends
$O$ to $x_0=0$ and $P$ to $x_a=a$.
The two points $x_0,x_a$ have type 2, since their coordinates belong
to $v(K^\times)$. By \cite[Proposition 3.13(2) and Theorem 4.11]{BPR14},
$\{x_0,x_a\}$ is a semistable vertex set corresponding to a proper
semistable formal model $\mathfrak E$. Its two edges have lengths
$a$ and $\ell-a$; neither is a loop edge, so the model is strictly
semistable. There are two special-fibre components $C_0,C_a$, and
all component multiplicities and nodal residue-field degrees are one.
If algebraic models are preferred, the formal model algebraizes as in
\cite[Remark 3.3]{GJKM19}.

The points $O$ and $P$ reduce to smooth points of $C_0$ and
$C_a$, respectively. Indeed, a node has open-annular tube whose
skeleton is the interior of an edge, whereas a residue disc at a smooth
point retracts to the vertex; this is the formal-fibre description in
\cite[Section 1.3 and Proposition 4.10]{BPR14}. Concretely, the reductions
are the unit point in the annular coordinates $z$ and $z/u$, not
either of the nodal directions $0,\infty$.
Their horizontal closures are Cartier divisors: in a smooth formal
coordinate at either section the closure is cut out by that coordinate
minus the coordinate of the section. Reduction gives a simple
rational point on the indicated component and meets no other component.
Thus
\[
 \mathfrak L_0
   :=\mathcal O_{\mathfrak E}(\overline P-\overline O)
\]
is a formal model of $L$ with component degrees
\begin{equation}
\label{eq:tate-reference-degrees}
 \deg(\mathfrak L_0|C_0)=-1,
 \qquad \deg(\mathfrak L_0|C_a)=1.
\end{equation}
Equivalently its model curvature measure is
$\delta_{x_a}-\delta_{x_0}$.
This last formulation uses only the usual curve model-curvature
formula; the rest of the obstruction can instead be read entirely in
terms of the two component degrees in
\eqref{eq:tate-reference-degrees}.
Since $P,O$ are $K$-points, $\deg L=1-1=0$.

\emph{Step 2: an arbitrary hypothetical nef model can be tested on a
common semistable refinement.}
Suppose that some $m>0$ and some proper admissible formal model
carry a nef model $\mathfrak N$ of $L^{\otimes m}$.
Let $\|\cdot\|_N$ be its formal metric and
$\|\cdot\|_0$ the metric of $\mathfrak L_0$, and put
\[
 \varphi=-\log\frac{\|\cdot\|_N}{\|\cdot\|_0^{\otimes m}}.
\]
This is a $\mathbf Z$-model function; the metric--function dictionary
is \cite[Section 4.7]{GM19}. Take a common model and then a strictly
semistable model $\mathfrak E'$ dominating $\mathfrak E$ on which
$\varphi$ is determined. The semistable domination used here is the
one recalled in the proof of \cite[Proposition 3.7]{GJKM19}, citing
Bosch--L\"utkebohmert \cite[Section 7]{BL85}; it applies over our already
algebraically closed complete field. Formal/algebraic interchange on a
proper curve is allowed by \cite[Theorem 1.1]{GM19} and \cite[Remark 3.3]{GJKM19}.
Nefness is preserved by pullback to the common model, by the projection
formula; alternatively \cite[Proposition 3.5]{GM19} tests semipositivity on
any representing model.

Write $\Sigma'=S(\mathfrak E')$ and let
$\tau':E_q^{\mathrm{an}}\to\Sigma'$ be its retraction.
The skeleton $\Sigma'$ is obtained by subdividing $\Sigma$ and
attaching finitely many finite trees, each at one point of $\Sigma$.
This follows from \cite[Proposition 3.13]{BPR14} and the open-ball components
of the complement of the minimal skeleton. By
\cite[Proposition 3.5]{GJKM19},
\[
 \varphi=F\circ\tau',
\]
where $F$ is continuous, affine on the edges of a finite subdivision
of $\Sigma'$, and all its slopes are integers. Additional subdivision
points can be incorporated in the semistable vertex set.
Pullback of $\mathfrak L_0^{\otimes m}$ retains degree $-m$ and
$m$ at the vertices $x_0,x_a$; its degree on every new component
is zero. In fact a new component maps to a closed point of the original
special fibre, and the old components map birationally to the old ones.

\emph{Step 3: positivity leaves no curvature on a degree-zero graph.}
Use the outgoing-slope convention of \cite[Proposition 3.7]{GJKM19}. For
each vertex $x$ of the refined graph set
\[
 b_x=m\,\mathbf 1_{\{x_a\}}(x)-m\,\mathbf 1_{\{x_0\}}(x).
\]
The slope formula and nefness of $\mathfrak N$ give
\begin{equation}
\label{eq:tate-nef-vertex-inequality}
 b_x+\sum_{\nu\in T_x\Sigma'}\partial_\nu F\ \geq\ 0.
\end{equation}
All tangent-direction weights are one because the residue field is
algebraically closed. Summing the left-hand sides over all vertices
gives zero: $\sum_x b_x=0$, and each edge contributes its slope and
its negative at its two endpoints. A finite sum of nonnegative numbers
is zero only if every summand is zero. Therefore
\begin{equation}
\label{eq:tate-zero-vertex-curvature}
 b_x+\sum_{\nu\in T_x\Sigma'}\partial_\nu F=0
 \quad\hbox{at every vertex }x.
\end{equation}
No global CLD-current positivity theorem or approximation theorem is
needed for this finite-graph conclusion.

On every tree attached to the circle, $F$ is constant. To prove this,
take a leaf away from the attachment point. Here $b_x=0$, and its
single outgoing slope is zero by
\eqref{eq:tate-zero-vertex-curvature}. Remove that edge and leaf and
repeat. This terminates because the tree is finite. In particular,
attached trees contribute no slope at their roots.
At a subdivision vertex of either original circle edge,
\eqref{eq:tate-zero-vertex-curvature} now says that the two oriented
slopes agree. Hence the restriction of $F$ to the original circle
has a constant integer slope $s$ on $(0,a)$ and a constant integer
slope $t$ on $(a,\ell)$, with the same orientation of the
coordinate on both intervals.

At $x_0$ the two outgoing slopes are $s$ and $-t$, and its
reference degree is $-m$. Thus
\begin{equation}
\label{eq:tate-jump}
 s-t=m.
\end{equation}
The continuity and single-valuedness of $F$ on the circle give
\begin{equation}
\label{eq:tate-period}
 as+(\ell-a)t=0.
\end{equation}
Solving \eqref{eq:tate-jump}--\eqref{eq:tate-period} yields
\[
 t=-m\frac{a}{\ell},\qquad
 s=m\left(1-\frac{a}{\ell}\right).
\]
Both numbers are irrational, contradicting their integrality. This
rules out the arbitrary model $\mathfrak N$, for every $m>0$.
Equivalently, working directly with a $\mathbf Q$-model correction
would require the rational slopes
$t=-a/\ell$ and $s=1-a/\ell$, which is the same contradiction.
Since a numerically trivial model is nef, the asserted nonexistence of
numerically trivial $\mathbf Q$-formal metrics follows as well.

\emph{Step 4: the canonical metric satisfies the CLD hypothesis.}
An elliptic curve is an algebraic abeloid variety, and our ground-field
absolute value is nontrivial. Fix any rigidification of $L$ at the
origin. \cite[Lemma 19.1.6]{CLD25} supplies its continuous canonical cubist
metric. On an algebraic curve numerical triviality of a line bundle is
equivalent to degree zero, so both $L$ and $L^{-1}$ are nef.
\cite[Proposition 19.3.1]{CLD25} states, under a nontrivial absolute value,
that for an algebraic abeloid the cubist metric is psh if and only if
the underlying line bundle is nef. \cite[Corollary 19.3.2]{CLD25} states
that its curvature current vanishes if and only if the line bundle is
numerically trivial. Applied here, these give
\[
 \|\cdot\|_{\mathrm{cub}}\ \text{locally CLD-psh},
 \qquad c_1(L,\|\cdot\|_{\mathrm{cub}})=0.
\]
The split-uniformization setting used in those results is satisfied
without any field extension: for a Tate curve the torus is
$\mathbf G_m$, the good-reduction abeloid is a point, and the lattice
is $q^{\mathbf Z}$.

For completeness the stronger local regularity asserted in the
statement follows from the same source. By \cite[Proposition 19.2.5]{CLD25},
after trivializing the pullback of $L$ to $\mathbf G_m^{\mathrm{an}}$,
the cubist weight differs from a constant by a quadratic function of
the tropical coordinate, with quadratic part $q_L/2$.
\cite[Corollary 19.3.2]{CLD25} gives $q_L=0$ for our numerically trivial
line bundle, so this weight is real affine in $-\log|z|$.
The latter is a CLD-smooth function furnished by an analytic unit and
has zero $\mathrm d'\mathrm d''$. The Tate covering has local analytic sections, so
these weights descend locally. Changing an invertible local frame
adds the logarithm of the absolute value of an analytic unit, again
smooth and harmonic. Thus the metric satisfies the local CLD-psh
hypothesis in the ordinary analytic topology, not just a formal
statement about its global curvature mass.

\emph{Step 5: the fixed-line-bundle closure is empty.}
In the convention of \cite[Introduction]{GM19}, a continuous semipositive
metric on $L$ is a uniform limit of semipositive model metrics on
the same $L$. Step 3 shows that there are no such model metrics at
all. Hence there cannot be an approximating sequence for the metric
constructed in Step 4. This proves the stated obstruction, without
making any assertion about definitions that allow an ample twist,
varying generic-fibre line bundles, or larger cones of real local
potentials.
\end{proof}


In \cref{thm:p081-irrational-tate-period-obstruction}, the cubist metric is
itself smooth and psh, so its constant sequence satisfies the definition of
globally psh at the beginning of this subsection.  Thus the example
separates that definition from the fixed-line-bundle rational-model closure;
it does not disprove the first question of
\cite[Remark~17.8.10]{CLD25}.  Nor is it a counterexample to the usual
ample-class envelope property: its line bundle has degree zero.

% The superseded unrestricted degree-zero theorem, its strict-analytic
% transfer, and the dependent alternative curve factorization are retained
% in the dated editorial archive, not as live proof inputs. Their corrected
% extensions have separate independent and production verification gates.

The proper strict analytic extension proposed in the second paragraph of
\cite[Remark~17.8.10]{CLD25} is false.  The next exact statement and proof
construct the counterexample on the Raynaud generic fibre of a very general
formal twistor deformation of a projective K3 surface.  The analytic Picard
group of that fibre is zero, while an admissible blow-up produces a
nonconstant formal metric on the trivial line bundle.  The maximum principle
forces every globally dpsh metric on that trivial bundle to be constant.
This argument concerns parlinear/global-dpsh approximation only; it neither
asserts that the exceptional metric is psh nor answers the continuous
psh/global-psh question.

For auditability, the following are the frozen principal bytes accepted by
the dependency-closed production endpoint.  The statement and proof have
SHA-256 hashes
\begin{center}
\small
\begin{tabular}{@{}c@{}}
\nolinkurl{1e9b7c31a70c55f407de07d76c67d6912e61e1673ef81c66c60789afd2202d2f}\\
\nolinkurl{78fab726abcdfb3086988150359568ccd8a9b29d63c3a94e0c513f14b9fc148c}.
\end{tabular}
\end{center}
The run used \texttt{gpt-5.6-sol} at \texttt{xhigh} and returned
\texttt{correct}, with zero critical errors and zero gaps.  The direct
endpoint assigned no run or item identifier; its request and response have
SHA-256 hashes
\begin{center}
\small
\begin{tabular}{@{}c@{}}
\nolinkurl{7916d3750482946cd933591a3f836d6e3b308dcbddb8a24a26c836de41d73e1b}\\
\nolinkurl{f5e780184515b39bbb7d4fd099366da98b96a74896fa1d693b43f66cce19d57e}.
\end{tabular}
\end{center}
The internal quarantine comments are retained solely to preserve the exact
verified bytes.  No novelty or priority claim is made.

\begin{proposition}\label{prop:p082-formal-twistor-counterexample}
Let $S$ be a complex projective K3 surface, let $\omega$ be a very
general K\"ahler class in the sense of
\cite[Definition~2.8]{HMS09Orientation}, and let

\[
        \mathfrak X\longrightarrow\operatorname{Spf}\mathbf C[[t]]
\]

be the induced formal twistor deformation.  Equip
$K=\mathbf C((t))$ with a complete $t$-adic absolute value satisfying
$0<|t|<1$, and let $X=\mathfrak X_\eta$ be the Raynaud generic fibre.
Then the following assertions hold.

\begin{enumerate}
\item The space $X$ is a connected smooth proper strict $K$-analytic
space, and $\mathfrak X$ is a formal model of $X$ whose special fibre is
the projective surface $S$.

\item One has
\begin{equation}\label{eq:p082-formal-twistor-picard-zero-candidate}
                         \operatorname{Pic}(X)=0.
\end{equation}

\item The trivial line bundle $\mathcal O_X$ admits a formal, hence
strictly integral parlinear, metric for which the function
$x\mapsto\lVert 1(x)\rVert$ is nonconstant.

\item Every globally dpsh metric on $\mathcal O_X$, in the sense of
\cite[\S17.8.5]{CLD25}, is constant.
\end{enumerate}

Consequently, the metric in the third item is not globally dpsh.  Thus the
second question in \cite[Remark~17.8.10]{CLD25}, interpreted as the proposed
extension of \cite[Corollary~17.8.9]{CLD25}, has a negative answer.
\end{proposition}

\begin{proof}[Proof of
\cref{prop:p082-formal-twistor-counterexample}]
Put
\[
                 R=\mathbf C[[t]],\qquad K=\mathbf C((t)).
\]
We separate the four assertions because the passage from formal to analytic
line bundles is load-bearing.

\emph{The analytic space.}
In the terminology of \cite[\S2.1]{HMS09Orientation}, a formal deformation
of $S$ is a smooth proper formal $R$-scheme whose special fibre is
$S$.  The formal twistor space associated to $\omega$ is such a
deformation.  Its Raynaud generic fibre is therefore a smooth proper strict
$K$-analytic space.  More precisely, the generic-fibre construction in
\cite[\S5.2.2]{Temkin12Introduction} produces a strict analytic space from
an admissible formal scheme of finite type, and
\cite[Fact~5.2.2.4]{Temkin12Introduction} takes formal smoothness to
rig-smoothness and formal properness to analytic properness.  A proper
analytic space is boundaryless, so rig-smoothness here is smoothness.  See
also \cite[Introduction and Remark~2.2]{HMS11FormalFibre}.  The original
formal model $\mathfrak X$ has projective special fibre $S$.

We verify connectedness rather than inferring it merely from domains on
formal affine charts.  Write
\[
  \mathfrak X_n
   =\mathfrak X\times_{\operatorname{Spf}R}\operatorname{Spec}
        \bigl(R/(t^{n+1})\bigr).
\]
Flatness gives, for $n\geq 1$, an exact sequence
\begin{equation}\label{eq:p082-infinitesimal-structure-sequence-candidate}
  0\longrightarrow \mathcal O_S
   \xrightarrow{\ t^n\ }
   \mathcal O_{\mathfrak X_n}
   \longrightarrow \mathcal O_{\mathfrak X_{n-1}}
   \longrightarrow 0.
\end{equation}
For a K3 surface,
$H^0(S,\mathcal O_S)=\mathbf C$ and
$H^1(S,\mathcal O_S)=0$.  Taking global sections in
\eqref{eq:p082-infinitesimal-structure-sequence-candidate} and comparing
with
\[
  0\longrightarrow \mathbf C\,t^n
   \longrightarrow R/(t^{n+1})
   \longrightarrow R/(t^n)
   \longrightarrow 0
\]
shows inductively that the structural map is an isomorphism
\begin{equation}\label{eq:p082-infinitesimal-global-functions-candidate}
       R/(t^{n+1})
          \xrightarrow{\ \sim\ }
       H^0(\mathfrak X_n,\mathcal O_{\mathfrak X_n})
       \qquad(n\geq0).
\end{equation}
Indeed, the right-hand restriction map is surjective because the next term
in the cohomology sequence is $H^1(S,\mathcal O_S)=0$, and the kernels in
the two displayed short exact sequences are identified by $t^n\mapsto
t^n\cdot1$.  Passing to inverse limits gives
\begin{equation}\label{eq:p082-formal-global-functions-candidate}
                \Gamma(\mathfrak X,\mathcal O_{\mathfrak X})=R.
\end{equation}
Here global sections commute with the inverse limit defining the structure
sheaf of the noetherian formal scheme.

By \cite[Remark~2.2 and Proposition~2.3]{HMS11FormalFibre}, coherent
sheaves on the Raynaud generic fibre are the categorical generic fibres of
coherent formal sheaves, and their morphisms are obtained by tensoring with
$K$.  Applying this to the structure sheaf and using
\eqref{eq:p082-formal-global-functions-candidate} yields
\begin{equation}\label{eq:p082-analytic-global-functions-candidate}
             \Gamma(X,\mathcal O_X)
              =\operatorname{End}(\mathcal O_X)
              \simeq R\otimes_R K=K.
\end{equation}
A disconnection of $X$ would give a nontrivial idempotent in
$\Gamma(X,\mathcal O_X)$, whereas the field $K$ has no such idempotent.
Thus $X$ is connected.  This proves the first assertion.

\emph{The analytic Picard group.}
Let $L$ be an analytic line bundle on $X$.  The fixed-model existence
theorem \cite[Theorem~B.1.8(a)]{CLD25} gives a coherent
$\mathcal O_{\mathfrak X}$-module $\mathcal F_0$ whose generic fibre is
$L$.  Let $\mathcal T$ be its $t$-power-torsion subsheaf and put
\[
                    \mathcal F=\mathcal F_0/\mathcal T.
\]
Noetherianity makes $\mathcal T$ coherent and locally annihilated by one
power of $t$.  Hence $\mathcal F_\eta=L$, while multiplication by $t$
is injective on $\mathcal F$.

We claim that $\mathcal F$ is torsion-free of rank one as an
$\mathcal O_{\mathfrak X}$-module.  Work in a formal local ring $A$ of
$\mathfrak X$.  The local map $R\to A$ is flat because
$\mathfrak X\to\operatorname{Spf}R$ is smooth, and $A/tA$ is a regular
local ring because it is a local ring of the smooth surface $S$.
Since $R$ is regular, \cite[Tag~031E]{StacksProject} shows that $A$ is
regular.  In particular, it is a local domain.  If
$0\neq a\in A$, $m\in\mathcal F_A$, and $am=0$, then after inverting
$t$ the element $a/1$ remains nonzero and the generic-fibre module is
locally free of rank one.  Therefore $m/1=0$, so $t^N m=0$ for some
$N$.  The $t$-torsion-freeness of $\mathcal F_A$ gives $m=0$.
Its rank is one because its generic fibre is $L$.

Consider the reflexive hull
\[
             \mathcal M:=\mathcal F^{\vee\vee}.
\]
It is coherent and has rank one because finite duals preserve coherence and
localization at the fraction field identifies its generic fibre with a
one-dimensional vector space.  It is reflexive by
\cite[Tag~0AV3]{StacksProject}, applied locally to
$\operatorname{Hom}_A(\mathcal F_A^\vee,A)$, and it can equivalently be
recognized by the normal-domain criteria in
\cite[Tags~0AVB and 0AVC]{StacksProject}.  To spell out the local algebra,
over a normal formal local domain $A$ the torsion-free rank-one module
$\mathcal F_A$ embeds in its one-dimensional vector-space generic fibre
and is therefore a fractional ideal.  Its double dual is the associated
divisorial fractional ideal.  Such
rank-one reflexive modules
are classified by the local Weil divisor class group
\cite[Tag~0EBM]{StacksProject}.  Every formal local ring $A$ of
$\mathfrak X$ is regular and hence a unique factorization domain
\cite[Tag~0AG0]{StacksProject}; its divisor class group is zero.
Consequently every stalk of $\mathcal M$ is free of rank one.  Coherence
then shows that $\mathcal M$ is a line bundle on $\mathfrak X$.

Formation of the dual of a finite module commutes with localization.
Since $\mathcal F[1/t]$ is already locally free, the natural map
$\mathcal F\to\mathcal M$ becomes an isomorphism after inverting $t$.
Consequently
\begin{equation}\label{eq:p082-double-dual-generic-fibre-candidate}
                         \mathcal M_\eta\simeq L.
\end{equation}
For a very general formal twistor deformation, every line bundle on
$\mathfrak X$ is trivial by
\cite[Proposition~2.10(i)]{HMS09Orientation}.  Thus
$\mathcal M\simeq\mathcal O_{\mathfrak X}$, and
\eqref{eq:p082-double-dual-generic-fibre-candidate} gives
$L\simeq\mathcal O_X$.  Since $L$ was arbitrary, this proves
\eqref{eq:p082-formal-twistor-picard-zero-candidate}.

\emph{A nonconstant formal metric on the trivial bundle.}
Choose a closed point $p\in S$, and let
\[
 \mathcal J=\ker\bigl(
       \mathcal O_{\mathfrak X}\longrightarrow (i_p)_*\mathbf C
                         \bigr).
\]
This is a coherent open ideal because it contains $t$.  Let
\[
       \mu\colon\mathfrak X'
          =\operatorname{Bl}_{\mathcal J}(\mathfrak X)
          \longrightarrow\mathfrak X
\]
be the admissible formal blow-up.  Since
$\mathcal J[1/t]=\mathcal O_{\mathfrak X}[1/t]$, its generic-fibre map is
an isomorphism by Raynaud's localization theorem
\cite[Exercise~5.3.2.3 and Fact~5.3.3.1]{Temkin12Introduction}; see also
the direct chart proof
in \cite[\S8.4, Proposition~2]{Bosch14FormalRigid}.  The special-fibre morphism
$\mathfrak X'_s\to S$ is a relative $\operatorname{Proj}$ of the
reduction of the finitely generated Rees algebra of $\mathcal J$
\cite[Tags~052P, 0804, and 01OF]{StacksProject};
hence it is projective.  Since $S$ is projective,
$\mathfrak X'_s$ is projective as well.  Thus even the formal model that
defines the metric below has projective special fibre.

The tautological transform of the centre ideal is invertible; write
\begin{equation}\label{eq:p082-exceptional-line-bundle-candidate}
       \mathcal J\mathcal O_{\mathfrak X'}
                 =\mathcal O_{\mathfrak X'}(-E).
\end{equation}
The inclusion of this ideal in $\mathcal O_{\mathfrak X'}$ identifies the
generic fibre of \eqref{eq:p082-exceptional-line-bundle-candidate} with
$\mathcal O_X$.  Denote its canonical formal metric by
$\lVert\cdot\rVert_E$.  By
\cite[Proposition~17.7.5(b)]{CLD25}, it is a strictly integral parlinear
metric.

We compute both the special fibre and the metric in one chart.  Smoothness
of relative dimension two gives a noncanonical identification
\[
       \widehat{\mathcal O}_{\mathfrak X,p}
             \simeq \mathbf C[[t,x,y]],
       \qquad
       \mathcal J_p=(t,x,y).
\]
Here the completion is the maximal-ideal completion of the noetherian local
ring $\mathcal O_{\mathfrak X,p}$, and the completion map is faithfully
flat by \cite[Tag~00MC]{StacksProject}.  Scheme-theoretic blow-up commutes
with this flat base change by \cite[Lemma~31.33.3, Tag~01OF]{StacksProject}.
Moreover, the formal blow-up is the $(t)$-adic completion of the
scheme-theoretic blow-up by
\cite[\S8.2, Proposition~6]{Bosch14FormalRigid}.  Consequently the affine
blow-up charts and their special-fibre equations may be computed after
passing to $\mathbf C[[t,x,y]]$.  Since the passage is faithfully flat,
the reduced component equations and the multiplicity-one assertions below
are detected there; adic completion does not alter the reduction modulo
$t$.
On the $x$-chart of the blow-up put
\[
                      T=t/x,\qquad Y=y/x.
\]
Thus $t=xT$, $y=xY$, and
$\mathcal J\mathcal O_{\mathfrak X'}=(x)$ on this chart.  Its special
fibre has equation
\[
                         xT=0.
\]
The exceptional component is $E=(x)$, while the strict transform of
$S$ is $S'=(T)$.  Both occur with multiplicity one in
\begin{equation}\label{eq:p082-special-fibre-components-candidate}
                         \operatorname{div}(t)=S'+E.
\end{equation}
Globally, $\mathfrak X'_s$ is the union of
$\operatorname{Bl}_p S$ and $E\simeq\mathbf P^2_{\mathbf C}$, meeting
along $\mathbf P^1_{\mathbf C}$.  Let $x_{S'}$ and $x_E$ be the
divisorial points of $X$ associated to these two components.

At the generic point of $S'$, the centre ideal $\mathcal J$ is the unit
ideal, since that generic point maps outside the closed point $p$.
Hence the ideal in
\eqref{eq:p082-exceptional-line-bundle-candidate} is generated by $1$,
and the formal-frame convention gives
\begin{equation}\label{eq:p082-metric-value-strict-transform-candidate}
                         \lVert1(x_{S'})\rVert_E=1.
\end{equation}
At the generic point of $E$, the function $x$ is a formal frame of
$\mathcal O_{\mathfrak X'}(-E)$.  Under the chosen generic-fibre
identification the trivial section satisfies
\[
                              1=x^{-1}\,x.
\]
Therefore
\begin{equation}\label{eq:p082-metric-value-exceptional-candidate}
                         \lVert1(x_E)\rVert_E=|x(x_E)|^{-1}.
\end{equation}
For a component $V$ of a formal special fibre, its divisorial seminorm is
normalized by
\[
 |f(x_V)|=|t|^{
       \operatorname{ord}_V(f)/\operatorname{ord}_V(t)}.
\]
Since $t=xT$ and $T$ is a unit at the generic point of $E$,
$\operatorname{ord}_E(x)=\operatorname{ord}_E(t)=1$.  It follows that
\begin{equation}\label{eq:p082-two-metric-values-candidate}
        \lVert1(x_E)\rVert_E=|t|^{-1}\neq1
              =\lVert1(x_{S'})\rVert_E.
\end{equation}
For complete sign bookkeeping, the logarithmic norm
$\log\lVert1\rVert_E$ has respective values
$-\log|t|$ and $0$ at $x_E$ and $x_{S'}$.  With the CLD weight
convention $\varphi=\log\lVert1\rVert_E^{-1}$, the respective values are
$\log|t|$ and $0$.  This proves the third assertion.

\emph{Constancy of globally dpsh metrics on $\mathcal O_X$.}
We first show that every globally psh metric on an analytically trivial
line bundle over $X$ is constant in a global trivialization.  Fix a
nowhere-vanishing analytic section $s$.  By
\cite[\S17.8.5]{CLD25}, the metric is a uniform limit of smooth psh metrics
on the same line bundle.  If
\[
                         u_n=\log\lVert s\rVert_n^{-1}
\]
is the weight of such an approximating metric, then $u_n$ is a smooth
CLD-psh function.  CLD-psh functions are classically psh
\cite[\S1.2]{GR25Classical}.  The proper space $X$ is compact and
boundaryless, so $u_n$ attains a maximum.  Its maximum locus is closed by
continuity.  At each point of that locus, the local maximum principle
\cite[Theorem~6.1]{GR25Classical} makes $u_n$ constant on a
neighbourhood; hence the maximum locus is also open.  Connectedness of
$X$ makes it all of $X$.  Thus every $u_n$ is constant, and uniform
convergence shows that the limiting weight is constant as well.

Let now $\lVert\cdot\rVert$ be a globally dpsh metric on
$\mathcal O_X$.  By the definition in \cite[\S17.8.5]{CLD25}, there are
globally psh metrized line bundles
$(L_1,\lVert\cdot\rVert_1)$ and
$(L_2,\lVert\cdot\rVert_2)$, together with an isometric isomorphism
\begin{equation}\label{eq:p082-global-dpsh-decomposition-candidate}
  (\mathcal O_X,\lVert\cdot\rVert)
       \simeq
  (L_1,\lVert\cdot\rVert_1)
       \otimes (L_2,\lVert\cdot\rVert_2)^{-1}.
\end{equation}
By \eqref{eq:p082-formal-twistor-picard-zero-candidate}, choose analytic
trivializations $s_i$ of $L_i$.  The preceding paragraph shows that
\[
                        u_i=\log\lVert s_i\rVert_i^{-1}
\]
is constant for $i=1,2$.  In these trivializations, the isomorphism in
\eqref{eq:p082-global-dpsh-decomposition-candidate} sends $1$ to
\[
                         f\,s_1\otimes s_2^{-1}
\]
for a global analytic unit
$f\in\Gamma(X,\mathcal O_X^\times)$.  Hence the weight
$u=\log\lVert1\rVert^{-1}$ is
\begin{equation}\label{eq:p082-global-dpsh-weight-candidate}
                         u=u_1-u_2-\log|f|.
\end{equation}
The function $\log|f|$ is classically pluriharmonic by
\cite[Proposition~7.2(2)]{GR25Classical}.  It attains a maximum on the
compact space $X$.  Its maximum locus is open by the pluriharmonic local
maximum principle, closed because $\log|f|$ is continuous, and nonempty;
connectedness makes
it all of $X$.  Thus $\log|f|$ is constant.  Equation
\eqref{eq:p082-global-dpsh-weight-candidate} proves that $u$ is constant,
which establishes the fourth assertion.

The formal metric $\lVert\cdot\rVert_E$ is nonconstant by
\eqref{eq:p082-two-metric-values-candidate}, whereas every globally dpsh
metric on $\mathcal O_X$ is constant.  It is therefore not globally
dpsh, which disproves the proposed extension of
\cite[Corollary~17.8.9]{CLD25}.

This conclusion concerns only the second paragraph of
\cite[Remark~17.8.10]{CLD25}.  The constructed metric is parlinear; no
claim that it is psh is made.  In particular, the argument does not answer
the first paragraph's continuous-psh-to-global-psh question, nor does it
identify global dpsh approximation with local psh approximation.
\end{proof}


\section{Curvature products and relative Monge--Amp\`ere}
\label{sec:ch17-products}

Let \(X\) be pure \(n\)-dimensional and boundaryless.  If
\(\overline L_1,\ldots,\overline L_q\) are smooth, their mixed Chern form is
the ordinary wedge
\[
             c_1(\overline L_1)\wedge\cdots\wedge
             c_1(\overline L_q).
\]
For continuous psh weights, \cref{chap:cld-bedford-taylor} constructs the Bedford--Taylor
product recursively.  For dpsh metrics it is defined by multilinearity after
choosing local differences of psh metrics; locality of Bedford--Taylor
products makes the result independent of the decomposition.

\subsection{The construction, and what it inherits}
\label{sec:ch17-product-construction}

\begin{definition}[Mixed curvature product]
\label{def:ch17-mixed-product}
Let \(\overline L_1,\dots,\overline L_q\) be line bundles with dpsh metrics
on \(X\) and let \(T\) be a positive closed current of bidegree \((r,r)\)
with \(q+r\le n\).  On an open \(U\) on which each \(L_i\) has an invertible
section \(s_i\), set \(u_i=\log\|s_i\|^{-1}\) and
\[
 \bigl(c_1(\overline L_1)\wedge\cdots\wedge c_1(\overline L_q)\wedge T
 \bigr)\big|_U
 =\dprimeddprime u_1\wedge\cdots\wedge\dprimeddprime u_q\wedge T,
\]
the right-hand side being the dpsh Bedford--Taylor product of
\Cref{thm:ch16-dpsh-product}.  For \(T=\delta_X\) the factor \(T\) is
omitted from the notation.
\end{definition}

\begin{proposition}
\label{prop:ch17-product-welldefined}
\Cref{def:ch17-mixed-product} is independent of the choices of \(s_i\) and
glues to a global current of bidegree \((q+r,q+r)\).  It is symmetric and
multilinear in the \(\overline L_i\), and:
\begin{enumerate}
\item\label{itm:ch17-prod:closed}
 it is \(\dprime\)- and \(\ddprime\)-closed, and locally
 \(\dprimeddprime\)-exact when \(q\ge1\);
\item\label{itm:ch17-prod:bt}
 it is a Bedford--Taylor current when \(T\) is;
\item\label{itm:ch17-prod:pos}
 it is positive when all the metrics are psh and \(T\ge0\);
\item\label{itm:ch17-prod:skel}
 it is skeletal when all metrics are parlinear and \(T\) is skeletal, and
 in general, if all the metrics are tropically \(d\)-convex, its support is
 locally contained in a skeletal subset of \(X\).
\end{enumerate}
\end{proposition}

\begin{proof}
Changing \(s_i\) to \(f_is_i\) with \(f_i\in\cO^\times(U)\) replaces
\(u_i\) by \(u_i-\log|f_i|\).  Since \(\log|f_i|\) is smooth and
pluriharmonic---\(\dprimeddprime\log|f_i|=0\) on \(U\) by
\Cref{thm:cld-poincare-lelong}, \(f_i\) being invertible---the
multilinear map of \Cref{thm:ch16-dpsh-product} is unchanged: it descends,
as recorded in \eqref{eq:ch16-harmonic-quotient}, to the quotient by the
subsheaf of \emph{smooth} pluriharmonic functions, which is where
\(\log|f_i|\) lives.  Locality of the Bedford--Taylor product
(\Cref{thm:bt-symmetry-locality}) then lets the local currents be glued, and
symmetry and multilinearity are inherited from the same theorem.

\cref{itm:ch17-prod:closed}: closedness is
\Cref{prop:ch16-iterate-basic}\,\cref{itm:ch16-iter:posclosed}; local
exactness is the identity \eqref{eq:ch16-dpsh-outer} read backwards,
\(\dprimeddprime u_1\wedge\Theta=\dprimeddprime(u_1\Theta)\), valid since
\(u_1\) is dpsh.
\cref{itm:ch17-prod:bt} and \cref{itm:ch17-prod:pos} are
\Cref{prop:bt-current-stability} and
\Cref{prop:ch16-iterate-basic}\,\cref{itm:ch16-iter:posclosed}.
\cref{itm:ch17-prod:skel}: for parlinear metrics this is
\Cref{thm:ch16-skeletal-product}, and the general support statement is
\Cref{prop:ch16-skeletal-support}, applicable because a tropically
\(d\)-convex metric has tropically \(d\)-convex weights.
\end{proof}

\begin{proposition}
\label{prop:ch17-product-pushforward}
Let \(\overline L_1,\dots,\overline L_q\) be dpsh metrized line bundles on
\(X\).
\begin{enumerate}
\item\label{itm:ch17-pf:proper}
 If \(j:Y\to X\) is topologically proper between boundaryless,
 topologically separated spaces and \(S\) is a positive closed current on
 \(Y\), then
 \[
  c_1(\overline L_1)\wedge\cdots\wedge c_1(\overline L_q)\wedge j_*S
  =j_*\bigl(c_1(j^*\overline L_1)\wedge\cdots\wedge
    c_1(j^*\overline L_q)\wedge S\bigr).
 \]
\item\label{itm:ch17-pf:base}
 If \(K'/K\) is a complete valued extension and \(p:X_{K'}\to X\) the base
 change, then
 \(p_*\bigl(\bigwedge_ic_1(p^*\overline L_i)\bigr)
   =\bigwedge_ic_1(\overline L_i)\).
\item\label{itm:ch17-pf:finite}
 If \(f:Y\to X\) is a finite morphism of boundaryless \(n\)-dimensional
 spaces whose calibrated degree is constantly \(d\) on \(X^{\mathrm{gen}}\),
 then
 \(f_*\bigl(\bigwedge_ic_1(f^*\overline L_i)\bigr)
   =d\cdot\bigwedge_ic_1(\overline L_i)\).
\end{enumerate}
\end{proposition}

\begin{proof}
\cref{itm:ch17-pf:proper}: the statement is local on \(X\), and by
multilinearity we may assume all metrics psh; then it is
\Cref{prop:bt-projection-formula} applied to the local weights, using
\(j^*\log\|s_i\|^{-1}=\log\|j^*s_i\|^{-1}\), which is
\Cref{prop:ch17-metric-picard}\,\cref{itm:ch17-pic:pullback}.  Items
\cref{itm:ch17-pf:base} and \cref{itm:ch17-pf:finite} follow from
\cref{itm:ch17-pf:proper} together with the calibrated degree computations of
\Cref{chap:cld-currents}: \Cref{prop:current-pushforward-formal} and
\Cref{cor:ch14-pushforward-functorial} give
\(p_*p^*=\mathrm{id}\) on currents for a base change, and
\(f_*f^*=d\) for a finite morphism of constant calibrated degree \(d\).
\end{proof}

\subsection{The total mass is an intersection number}
\label{sec:ch17-degree}

The following two results are the reason the whole construction is worth
making: the top product has a total mass that does not see the metric, and
on an algebraic variety that mass is the classical intersection number.

\begin{proposition}
\label{prop:ch17-mass-independent}
Let \(X\) be boundaryless of pure dimension \(n\) and let
\(L_1,\dots,L_n\) be line bundles, each equipped with two smooth (resp.\
two dpsh) metrics.  Then the difference of the two top products is
\(\dprimeddprime\alpha\) for a smooth (resp.\ a symmetric closed)
\((n-1,n-1)\)-form (resp.\ current) \(\alpha\).  If \(X\) is in addition
proper, the number
\(\int_Xc_1(\overline L_1)\wedge\cdots\wedge c_1(\overline L_n)\)
depends only on \(L_1,\dots,L_n\).
\end{proposition}

\begin{proof}
Let \(u_i\) be the global function with
\(\log\|s\|=u_i+\log\|s\|'\) for every local invertible section \(s\) of
\(L_i\); it is smooth (resp.\ dpsh) by hypothesis, and
\(c_1(\overline L_i)=\dprimeddprime u_i+c_1(\overline L_i')\) by
\Cref{prop:cld-c1-metric-change}.  Expanding the product,
\(\bigwedge_ic_1(\overline L_i)-\bigwedge_ic_1(\overline L_i')\) is a sum of
terms \(T_1\wedge\cdots\wedge T_n\) in which each \(T_i\) is either
\(c_1(\overline L_i)\) or \(\dprimeddprime u_i\), with at least one index
\(j\) of the second kind.  For such a term,
\eqref{eq:ch16-dpsh-outer} gives
\[
 T_1\wedge\cdots\wedge T_n
 =\dprimeddprime\bigl(u_j\,T_1\wedge\cdots\widehat{T_j}\cdots\wedge T_n
   \bigr),
\]
so the difference is \(\dprimeddprime\alpha\) with \(\alpha\) a sum of such
brackets, smooth in the smooth case and a symmetric closed
\((n-1,n-1)\)-current in general.  If \(X\) is proper,
\(\int_X\dprimeddprime\alpha=\langle\dprimeddprime\alpha,1\rangle
 =\langle\alpha,\dprimeddprime1\rangle=0\)
by \Cref{prop:current-differential-calculus} and
\Cref{cor:analytic-stokes-green}, the constant \(1\) being an admissible
test form on a compact boundaryless space.
\end{proof}

\begin{theorem}
\label{thm:ch17-degree}
\index{intersection number!as total curvature mass}
Let \(X\) be a proper \(K\)-variety of pure dimension \(n\) and let
\(L_1,\dots,L_n\) be line bundles on \(X\), each equipped with a dpsh metric
on \(X^{\An}\)---for instance a smooth or a formal one.  Then
\[
 \int_{X^{\An}}c_1(\overline L_1)\wedge\cdots\wedge c_1(\overline L_n)
 =(L_1\cdots L_n)_X,
\]
the intersection number computed on the scheme \(X\).  In particular the
left-hand side is an integer, nonnegative when the \(L_i\) are nef.
\end{theorem}

\begin{proof}
By \Cref{prop:ch17-mass-independent} and
\Cref{prop:ch17-smooth-existence} we may assume the metrics smooth.  Let
\(X_j\) be the reduced irreducible components of \(X\) with generic
multiplicities \(m_j\).  The \(X_j^{\An}\) are the irreducible components of
\(X^{\An}\) with the same generic multiplicities, so both sides are
\(\sum_jm_j\) times the corresponding quantity on \(X_j\)---on the analytic
side by \Cref{thm:current-fundamental-cycle}, on the algebraic side by the
definition of the intersection number.  We may therefore assume \(X\)
integral, and induct on \(n\).

If \(n=0\) then \(X=\Spec F\) for a finite extension \(F/K\), and both sides
equal \([F:K]\): analytically by \Cref{prop:ch14-delta-Y-current}, and
algebraically by definition.

Let \(n>0\) and assume the statement in dimension \(n-1\).  Choose a nonzero
rational section \(s\) of \(L_n\), with divisor \(D=\sum_jm_jD_j\) on \(X\),
whose analytification is the divisor \(D^{\An}\) of \(s\) viewed as a
meromorphic section on \(X^{\An}\).  Put
\(\alpha=c_1(\overline L_1)\wedge\cdots\wedge c_1(\overline L_{n-1})\), a
smooth \((n-1,n-1)\)-form.  Locally \(\alpha=\dprimeddprime u_1\wedge\cdots
\wedge\dprimeddprime u_{n-1}\) with \(u_i\) smooth, so
\(\dprimeddprime\alpha=0\).  Pairing
\Cref{thm:ch17-pl-metrized} with \(\alpha\) gives
\[
 \underbrace{\langle[\log\|s\|^{-1}],\dprimeddprime\alpha\rangle}_{=0}
 +\int_{D^{\An}}\alpha
 =\int_{X^{\An}}c_1(\overline L_1)\wedge\cdots\wedge c_1(\overline L_n).
\]
On the other hand, by the definition of integration over a cycle and the
inductive hypothesis applied to each \(D_j\),
\[
 \int_{D^{\An}}\alpha
 =\sum_jm_j\int_{D_j^{\An}}\alpha
 =\sum_jm_j\,(L_1\cdots L_{n-1})_{D_j}
 =(L_1\cdots L_n)_X,
\]
the last equality because \(c_1(L_n)\cap[X]=[D]\).  This proves the formula.
Integrality and the nef statement now follow because the right-hand side is
an intersection number.  This is \cite[Theorem~18.3.11]{CLD25}.
\end{proof}

\begin{remark}[The nonalgebraic case]
\label{rem:ch17-integrality-analytic}
Integrality of \(\int_Xc_1(\overline L_1)\wedge\cdots\wedge
c_1(\overline L_n)\) also holds on an arbitrary proper \(K\)-analytic space,
without algebraicity.  The argument is: the number is metric-independent by
\Cref{prop:ch17-mass-independent}; it may be computed after a complete
scalar extension by
\Cref{prop:ch17-product-pushforward}\,\cref{itm:ch17-pf:base}, so one may
assume \(|K^\times|\ne1\); by
\Cref{thm:cld-formal-pl-comparison} the \(L_i\) then carry strictly integral
parlinear metrics; and for those the product is the finite atomic measure
\Cref{thm:ch17-atomic-support} whose masses are integers by
\Cref{thm:ch17-atomic-masses}.  We shall not need this refinement, and we
record it only so that the reader is not left with the impression that
\Cref{thm:ch17-degree} is the strongest available statement.  It is
\cite[Corollary~18.4.10]{CLD25}.
\end{remark}

\subsection{The top product of parlinear metrics is atomic}
\label{sec:ch17-atomic}

Integral PL metrics admit an additional construction.  Their first Chern
objects are singular, but the residual intersection theory defines their
mixed product directly.  For \(q=n\) the product is a locally finite atomic
signed measure.  If the residual bundles are nef, every atom has
nonnegative mass.

We now prove that assertion.  Recall from \Cref{chap:analytic-tropicalization}
that \(X^{\#}\) denotes the set of points \(x\in X\) with
\(d_K(x)=n\), the maximal Abhyankar dimension; at such a point
\(\widetilde{\cH(x)}\) has transcendence degree \(n\) over \(\widetilde K\)
and is finitely generated over \(\widetilde{\cH(K)}\).

\begin{theorem}
\label{thm:ch17-atomic-support}
Let \(X\) be boundaryless of pure dimension \(n\) and let
\(\overline L_1,\dots,\overline L_n\) be integral parlinear metrized line
bundles.  Let \(S\subseteq X\) be the set of points at which the
simultaneous residual dimension of definition of the
\(\widetilde{L_{i,x}}\) equals \(n\).  Then \(S\) is closed and discrete,
\(S\subseteq X^{\#}\), each \(\overline L_i\) is dpsh, and there are real
numbers \((\lambda_x)_{x\in S}\) with
\[
 \bigwedge_{i=1}^nc_1(\overline L_i)=\sum_{x\in S}\lambda_x\delta_x.
\]
\end{theorem}

\begin{proof}
That \(S\) is closed and discrete is
\Cref{prop:ch17-dimdef-discrete} with \(d=n\); that \(S\subseteq X^{\#}\)
is the inequality ``residual dimension of definition
\(\le\operatorname{tr.deg}\widetilde{\cH(x)}\le d_K(x)\le n\)'' together
with the definition of \(S\).  Each \(\overline L_i\) is tropically
\(d\)-convex, hence dpsh, by \Cref{prop:ch16-dpsh-smooth}.

Fix \(x\in X\); we show that near \(x\) the product is \(\lambda_x\delta_x\)
for some real \(\lambda_x\), vanishing if \(x\notin S\).  Choose a premodel
\((V,M_1,\dots,M_n)\) of \((\widetilde{L_{1,x}},\dots,\widetilde{L_{n,x}})\)
with \(\dim V\) minimal, that is, equal to the simultaneous residual
dimension of definition.  By multilinearity and the polarization identity we
may assume all the \(M_i\) equal to one globally generated \(M\), so that
near \(x\) all the \(\overline L_i\) are one metrized bundle
\(\overline L\).  By the global-generation description used in the proof of
\Cref{prop:ch17-globally-generated-psh}, there are an invertible section
\(s\) of \(L\) near \(x\), reals \(r_1,\dots,r_m>0\) and units
\(a_1,\dots,a_m\in\cO_{X,x}^\times\) with
\(\widetilde{a_i}/\widetilde{a_j}\in\widetilde K(V)\) for all \(i,j\), such
that
\[
 \|s\|^{-1}=\max\{r_1|a_1|,\dots,r_m|a_m|\}
\]
near \(x\).  Thus the weight is the pullback of a convex piecewise affine
function under the moment map \(a=(a_1,\dots,a_m)\), and the local
Monge--Amp\`ere computation for such a function
(\Cref{ex:ch16-fan-p2} and the general statement quoted there) shows that
\(c_1(\overline L)^n\) is \(\lambda_x\delta_x\) near \(x\), the coefficient
being the tropical-cone invariant of the moment \(a\) at \(x\).

It remains to see that \(\lambda_x=0\) when \(x\notin S\).  In that case
\(\dim V<n\), and the relations \(\widetilde{a_i}/\widetilde{a_j}\in
\widetilde K(V)\) force the tropical dimension at \(x\) of the family
\((a_i/a_1)_i\) to be \(<n\), by the tropical-dimension formula
\Cref{thm:trop-dimension-formula}.  A convex piecewise affine function
pulled back from a tropical set of dimension \(<n\) has vanishing \(n\)-fold
Monge--Amp\`ere measure, since its subgradient image is contained in a
proper affine subspace; hence \(c_1(\overline L)^n=0\) near \(x\).  This is
\cite[Proposition~18.4.3]{CLD25}.
\end{proof}

\begin{theorem}
\label{thm:ch17-atomic-masses}
In the situation of \Cref{thm:ch17-atomic-support}, the mass of any
\(x\in X^{\#}\) is the integer
\begin{equation}
\label{eq:ch17-atomic-mass}
 \lambda_x=m(x)\cdot\mathrm{def}_K(x)\cdot
 \int_{P_{\widetilde{\cH(x)}/\widetilde K}}
   c_1(\widetilde{L_{1,x}})\cdots c_1(\widetilde{L_{n,x}}),
\end{equation}
where \(m(x)\) is the multiplicity of the unique irreducible component of
\(X\) containing \(x\), \(\mathrm{def}_K(x)\) is the geometric defect of
\(x\), and the integral is a graded intersection number.
\end{theorem}

\begin{proof}
We only sketch the reduction, the graded intersection theory being quoted
from \Cref{sec:appG-graded-reduction}; the full argument is
\cite[Proposition~18.4.5 and Theorem~18.4.7]{CLD25}.

\emph{Step 1: reduction to a strictly integral parlinear metric at a
strictly Abhyankar point.}  Choose finitely many affinoid domains
\(V_j\ni x\) covering a neighbourhood of \(x\) on which each \(L_i\) has a
section whose norm lies in \(\rho_{ij}|\cO^\times(V_j)|\).  Let \(r\) be a
finite polyradius containing all the \(\rho_{ij}\) and a generating system
of \((|\cH(x)^\times|/|K^\times|)\otimes\mathbf Q\), and let \(K_r\) be the
corresponding complete extension.  Pulling back to \(X_r=X\times_KK_r\) and
using \Cref{prop:ch17-product-pushforward}\,\cref{itm:ch17-pf:base}, the mass
\(\lambda_x\) equals the mass at the Shilov section \(\sigma(x)\) of the
product upstairs; and the metrics become strictly integral parlinear.  The
factors \(m(x)\) and \(\mathrm{def}_K(x)\) are unchanged by this extension,
which is analytically separable.

\emph{Step 2: the strict case over an algebraically closed base.}  Assume
\(K\) algebraically closed, \(X\) generically reduced, and the metrics
strictly integral parlinear.  Then \(m(x)=\mathrm{def}_K(x)=1\), and one
must show that \(\lambda_x\) is the graded degree.  As in
\Cref{thm:ch17-atomic-support} one reduces to a single \(\overline L\) with
globally generated residual bundle and
\(\|s\|^{-1}=\max_i|f_i|\), \(|f_i(x)|=1\).  Because \(d_K(x)=n\), the local
ring \(\cO_{X,x}\) is an artinian reduced ring, hence a field, so \(X\) is
smooth near \(x\) and may be replaced by the analytification of a smooth
affine \(K\)-variety on which the \(f_i\) are units.  Let \((V,M)\) be a
proper model of \(\bigl(\widetilde{(X,x)},\widetilde L_x\bigr)\), and let
\(\xi\in V^{\An}\) be the point given by the trivial absolute value on
\(\widetilde K(V)\), so that \(\widetilde{\cH(\xi)}=\widetilde{\cH(x)}\).
The tropical cones of the moment \((f_i)\) at \(x\) and of the moment
\((\widetilde{f_i(x)})\) at \(\xi\) coincide
(\Cref{thm:trop-local-germ}), so \(\lambda_x=\lambda_\xi\), where
\(c_1(\overline{M_\eta})^n=\lambda_\xi\delta_\xi\) on \(V^{\An}\) for the
formal metric of \(M\) over the trivially valued \(\widetilde K\).  But
\(\xi\) is the only point of \(V^{\An}\) with residue transcendence degree
\(n\), so that measure has total mass \(\lambda_\xi\), and by
\Cref{thm:ch17-degree} that mass is \(\int_Vc_1(M)^n\).

\emph{Step 3: descent.}  For general \(K\), pass to the completion
\(\widehat{\overline K}\) and to the reduction \(X_{K'\!,\mathrm{red}}\).
The calibrated degree of the closed immersion
\(X_{K'\!,\mathrm{red}}\to X_{K'}\) is \(1/m_{K'}\), so
\Cref{prop:ch17-product-pushforward} turns the measure downstairs into
\(m_{K'}\) times the pushforward of the measure upstairs.  Applying Step 2
upstairs and summing over the fibre of \(x\), the combinatorial count of the
graded degrees over the fibre produces exactly the factor
\(m(x)\mathrm{def}_K(x)\) and converts the sum of graded degrees over
\(\widetilde{\cH(y)}\) into the single graded degree over
\(\widetilde{\cH(x)}\).  Integrality is then visible from
\eqref{eq:ch17-atomic-mass}, all three factors being integers.
\end{proof}

\begin{corollary}
\label{cor:ch17-atoms-nonneg}
If every residual bundle \(\widetilde{L_{i,x}}\) is nef, then
\(\lambda_x\ge0\) for all \(x\).  In particular, if each
\(\overline L_i\) has positive curvature current, the top product is a
positive measure.
\end{corollary}

\begin{proof}
The graded intersection number in \eqref{eq:ch17-atomic-mass} is the degree
of a product of nef classes, hence nonnegative by Kleiman's theorem in the
graded setting; \(m(x)\) and \(\mathrm{def}_K(x)\) are positive.  The last
sentence follows from \Cref{thm:ch17-psh-residual-nef}.
\end{proof}

\subsection{The relative Monge--Amp\`ere operator}
\label{sec:ch17-relative-ma}

The relative version is particularly transparent for PL metrics.  The
statement to keep in mind is: pulling back \(m\) parlinear metrics from an
\(m\)-dimensional base and taking the top product turns each atom of the
base measure into the integration current over the corresponding fibre, with
the same coefficient.  In other words, the relative Monge--Amp\`ere operator
of parlinear metrics computes fibres.

Before the theorem we record the fibre-integration statement it rests on,
which is \Cref{thm:cld-fibre-integration} in the form used here.

\begin{proposition}
\label{prop:ch17-relative-pairing}
Let \(p:Y\to X\) be as in \Cref{thm:cld-relative-ma}, let
\(\overline L_1,\dots,\overline L_r\) be tropically \(d\)-convex metrized
line bundles on \(X\), let \(\omega\) be a \(G\)-smooth \((d,d)\)-form on
\(Y\) with \(p\)-compact support, and let
\(u_\omega(x)=\int_{Y_x}\omega\) when \(\dim Y_x=d\), and \(0\) otherwise.
Then for every compactly supported smooth \((m-r,m-r)\)-form \(\alpha\) on
\(X\),
\[
 \bigl\langle c_1(p^*\overline L_1)\wedge\cdots\wedge
   c_1(p^*\overline L_r),\ \omega\wedge p^*\alpha\bigr\rangle
 =\int_Xu_\omega\;
   c_1(\overline L_1)\wedge\cdots\wedge c_1(\overline L_r)\wedge\alpha .
\]
Equivalently,
\(p_*\bigl(\bigwedge_ic_1(p^*\overline L_i)\wedge p^*\alpha\wedge\omega\bigr)
 =u_\omega\cdot\bigwedge_ic_1(\overline L_i)\wedge\alpha\).
\end{proposition}

\begin{proof}
Using a partition of unity on \(X\)
(\Cref{thm:cld-smooth-partition}) we may assume the \(L_i\) trivial and
each \(c_1(\overline L_i)\) a difference of currents \(\dprimeddprime u_i\)
with \(u_i\) \(d\)-convex on \(X\); by multilinearity we may assume each
\(u_i\) convex.  The statement then becomes the projection formula for
Bedford--Taylor products along \(p\), which is
\Cref{thm:cld-fibre-projection}.
\end{proof}


\begin{theorem}
\label{thm:cld-relative-ma}
Let \(p:Y\to X\) be a generically flat morphism between boundaryless,
Hausdorff, equidimensional analytic spaces, with
\(\dim X=m\), \(\dim Y=n\), and \(d=n-m\).  Let
\(\overline L_1,\ldots,\overline L_m\) be integral PL metrized line bundles
on \(X\).  Write
\[
 c_1(\overline L_1)\wedge\cdots\wedge c_1(\overline L_m)
       =\sum_{x\in X^\#}\lambda_x\delta_x,
\]
where \(X^\#\) is the set of points of maximal Abhyankar dimension.  Then,
as a current of bidimension \((d,d)\) on \(Y\),
\[
 \bigwedge_{i=1}^m c_1(p^*\overline L_i)
        =\sum_{x\in X^\#}\lambda_x\,[Y_x].
\]
The nonzero coefficients form a closed discrete family, so the right-hand
side is locally finite.
\end{theorem}
\begin{proof}
Put the integral PL metrics on a common formal model of $X$ and dominate its
pullback by a formal model of $Y$.  On a residual chart, the mixed Chern
product is computed by intersecting the residual line bundles.  Its
top-dimensional components lie over the generic points represented by
$X^\#$; the coefficient over such a point $x$ is precisely $\lambda_x$.

Generic flatness identifies the pullback of the residual fundamental cycle
at $x$ with the fundamental cycle of the analytic fibre $Y_x$, including
component lengths and residue-field degrees.  Thus the residual intersection
calculation gives, chartwise,
\[
       \bigwedge_{i=1}^m c_1(p^*\overline L_i)
          =\sum_x\lambda_x[Y_x].
\]
To verify the equality as currents, pair both sides with a compactly
supported $(d,d)$-form $\omega$.  Calibrated fibre integration gives
\[
 \left\langle\bigwedge_i c_1(p^*\overline L_i),\omega\right\rangle
   =\sum_x\lambda_x\int_{Y_x}\omega,
\]
which is the desired pairing.  On a relatively compact set only finitely
many residual components meet the support; hence the nonzero atoms are a
closed discrete family and the sum is locally finite.  This is the residual
intersection and fibre-integration proof of \cite[\S18.6]{CLD25}.
The boundaryless, generically flat, and integral-PL hypotheses are used,
respectively, in Stokes, the fibre-cycle identity, and residual
intersection.
\end{proof}

\begin{proof}[Second proof, by fibre integration alone]
The displayed formula for the base measure is
\Cref{thm:ch17-atomic-support} together with
\Cref{thm:ch17-atomic-masses}, and the closed discreteness of the support
is \Cref{prop:ch17-dimdef-discrete}.  Let \(\omega\) be a smooth
\((d,d)\)-form with compact support on \(Y\); we must show
\[
 \Bigl\langle\bigwedge_{i=1}^mc_1(p^*\overline L_i),\omega\Bigr\rangle
 =\sum_{x\in X^{\#}}\lambda_x\int_{Y_x}\omega .
\]
Apply \Cref{prop:ch17-relative-pairing} with \(r=m\) and
\(\alpha=1\)---legitimate because the support of \(\omega\) is compact,
hence \(p\)-compact, and \(p_*\omega\) has compact support.  The left-hand
side becomes \(\int_Xu_\omega\,\bigwedge_ic_1(\overline L_i)\), which by
atomicity is \(\sum_{x}\lambda_xu_\omega(x)\).  For \(x\) in the support one
has \(d_K(x)=m\), so \(\dim Y_x\le n-m=d\); generic flatness forces
equality, so \(u_\omega(x)=\int_{Y_x}\omega=\langle\delta_{Y_x},\omega\rangle\).
This is exactly the required identity.  Note that this argument uses
generic flatness only through the equality \(\dim Y_x=d\), and that it never
mentions a formal model.
\end{proof}

The following corollary is the computational engine of
\Cref{sec:ch17-bm-levine}: it identifies the \(d\)-fold Monge--Amp\`ere
current of a truncated maximum of absolute values with an integration
current over a fibre.

\begin{corollary}
\label{cor:ch17-truncated-max-fibre}
Let \(X\) be boundaryless, topologically separated and equidimensional, and
let \(f=(f_1,\dots,f_d):X\to\mathbf A^{d,\An}\) be given by \(d\) analytic
functions.  For real numbers \(t,r_1,\dots,r_d>0\) put
\[
 \varphi=\log\max\{t,\ |f_1|/r_1,\ \dots,\ |f_d|/r_d\}.
\]
Then \((\dprimeddprime[\varphi])^d=\delta_{f^{-1}(\eta_\rho)}\), the
integration current over the fibre of \(f\) above the Gauss point
\(\eta_\rho\) of polyradius \(\rho=(tr_1,\dots,tr_d)\).
\end{corollary}

\begin{proof}
\emph{Step 1: the support.}  Let \(x\in X\).  Suppose first that
\(|f_i(x)|/r_i\ne t\) for some \(i\).  Then near \(x\) we may write
\(\varphi=\log\max\{|g_1|/s_1,\dots,|g_p|/s_p\}\) with \(p\le d\), the
\(g_j\) invertible near \(x\), the \(s_j>0\), and \(|g_j(x)|/s_j=\rho'\)
independent of \(j\)---each pair \((g_j,s_j)\) being one of the
\((f_i,r_i)\) attaining the maximum, or \((1,1/t)\).  Setting
\(\psi=\varphi-\log|g_1|\) gives \(\dprimeddprime\varphi=\dprimeddprime\psi\)
near \(x\), and \(\psi\) is the pullback of a convex piecewise affine
function of \(p-1\le d-1\) variables under the moment
\((g_2/g_1,\dots,g_p/g_1)\).  Its tropical dimension is at most \(d-1\), so
the \(d\)-fold Monge--Amp\`ere current vanishes near \(x\), as in the last
step of the proof of \Cref{thm:ch17-atomic-support}.

Suppose next that \(|f_i(x)|=\rho_i\) for all \(i\).  The same argument with
the moment \((f_1,\dots,f_d)\) shows that
\((\dprimeddprime\varphi)^d\) vanishes near \(x\) unless the tropical
dimension of \(f\) at \(x\) is \(d\).  Because \(X\) is boundaryless, that
tropical dimension is the transcendence degree of the family
\((\widetilde{f_i(x)})\) over \(\widetilde K\)
(\Cref{thm:trop-dimension-sigma}), so it equals \(d\) only when
\(f(x)\) is the Shilov point of the polydisc, that is, \(f(x)=\eta_\rho\).
Hence \(\Supp(\dprimeddprime\varphi)^d\subseteq f^{-1}(\eta_\rho)\).

\emph{Step 2: flatness.}  The local ring
\(\cO_{\mathbf G_{\mathrm m}^{d,\An},\eta_\rho}\) is artinian and
\(\mathbf G_{\mathrm m}^{d,\An}\) is reduced, so that local ring is a field;
since \(X\) is boundaryless, \(f\) is flat at every point of
\(f^{-1}(\eta_\rho)\).  Replacing \(X\) by the flat locus of \(f\), which
by Step~1 does not change the current, we may assume \(f\) flat.

\emph{Step 3: identification of the coefficient.}  Let \(\overline N\) be
the trivial line bundle on \(\mathbf A^{d,\An}\) with the metric determined
by \(\log\|1\|^{-1}=\log\max\{t,|T_1|/r_1,\dots,|T_d|/r_d\}\); it is
integral parlinear, and \(\dprimeddprime\varphi=c_1(f^*\overline N)\) by
construction.  \Cref{thm:cld-relative-ma} applied to \(f\) and to the
\(d\) copies of \(\overline N\) gives
\((\dprimeddprime\varphi)^d=\lambda\,\delta_{f^{-1}(\eta_\rho)}\) with
\(\lambda=\int_{P_{\widetilde{\cH(\eta_\rho)}/\widetilde K}}
 c_1(\widetilde N_{\eta_\rho})^d\).
It remains to see that \(\lambda=1\).  The residual bundle
\(\widetilde N_{\eta_\rho}\) is described by
\(v(\widetilde{1(\eta_\rho)})^{-1}
 =\max\{1,v(\widetilde{T_1(\eta_\rho)}),\dots,
        v(\widetilde{T_d(\eta_\rho)})\}\)
for every graded valuation \(v\); this is the residual line bundle of
\(\cO(1)\) on the graded projective space \(\mathbf P^d\) restricted along
the standard toric embedding, whose \(d\)-fold self-intersection is \(1\).
Hence \(\lambda=1\).  This is \cite[Corollary~18.6.2]{CLD25}.
\end{proof}

\section{Formal metrics and the special-fibre formula}
\label{sec:ch17-formal}

\subsection{The easy case: irreducible special fibre}
\label{sec:ch17-good-reduction}

It is worth doing the simplest case first, because it already contains the
mechanism and because it is the one used most often.

\begin{proposition}
\label{prop:ch17-good-reduction}
Let \(X\) be a proper \(K\)-scheme of dimension \(n\), let \(L\) be a line
bundle on \(X\), let \(\mathfrak X\) be a proper formal
\(K^\circ\)-model of \(X\) that is relatively normal in its generic fibre,
and let \(\mathfrak L\) be a line bundle on \(\mathfrak X\) extending \(L\).
Assume \(\mathfrak X_s\) irreducible, and let \(\eta\in X^{\An}\) be the
unique point whose reduction is the generic point of \(\mathfrak X_s\).
Give \(L\) the formal metric of \(\mathfrak L\).  Then
\[
 c_1(\overline L)^n=\Bigl(\int_Xc_1(L)^n\Bigr)\,\delta_\eta
 =(L^n)\,\delta_\eta .
\]
\end{proposition}

\begin{proof}
Let \(x\in X^{\An}\) with \(x\ne\eta\).  Its reduction has Zariski closure
of dimension \(\le n-1\), so by \Cref{lem:ch17-residual-formal} the residual
bundle \(\widetilde L_x\) has a premodel of dimension \(\le n-1\); hence its
dimension of definition is \(<n\), and \(x\notin S\) in the notation of
\Cref{thm:ch17-atomic-support}.  That theorem therefore gives
\(c_1(\overline L)^n=\lambda\,\delta_\eta\) for a single real number
\(\lambda\).  Its total mass is \(\lambda\), and by
\Cref{thm:ch17-degree} that mass is \(\int_Xc_1(L)^n=(L^n)\).
\end{proof}

\begin{example}[Projective space]
\label{ex:ch17-projective-good-reduction}
Take \(X=\mathbf P^n_K\), \(\mathfrak X\) the formal completion of
\(\mathbf P^n_{K^\circ}\), and \(\mathfrak L=\cO(1)\), so that \(L\) carries
the Weil metric of \Cref{ex:ch17-weil-metric}.  The special fibre is
\(\mathbf P^n_{\widetilde K}\), irreducible and reduced, and there is a
unique point \(\eta_1\in\mathbf P^{n,\An}_K\) reducing to its generic point:
the Gauss point of the standard simplex.  Hence
\[
 c_1\bigl(\overline{\cO(1)}\bigr)^n=\delta_{\eta_1},
 \qquad
 c_1\bigl(\overline{\cO(d)}\bigr)^n=d^n\,\delta_{\eta_1},
\]
of masses \(1\) and \(d^n=(\cO(d)^n)\).  For \(n=1\) this recovers
\Cref{ex:ch17-weil-curvature}\,(a) and, through it,
\eqref{eq:p1-standard-measure}.  For \(n=2\) it recovers the total mass
\(1\) computed in \Cref{ex:ch16-fan-p2} from the fan of \(\mathbf P^2\)
directly, and the reader may check that the two computations locate the same
atom: the vertex of the fan corresponds to the Gauss point.
\end{example}

\subsection{The three defect factors}
\label{sec:ch17-defects}

We now state the current CLD theorem in a form that displays its
hypotheses.  Let \(\mathfrak X\) be a flat formal \(K^\circ\)-scheme,
locally topologically of finite presentation and pure of relative dimension
\(n\).  Assume that its generic fibre \(X\) is boundaryless.  Let
\(\mathfrak L_1,\ldots,\mathfrak L_n\) be line bundles on
\(\mathfrak X\), and give \(L_i=(\mathfrak L_i)_\eta\) their strictly
integral PL formal metrics.

For a point \(x\) specializing to the generic point of an irreducible
component \(V\) of \(\mathfrak X_s\), CLD associate a positive defect
factor \(\epsilon_{\mathfrak X}(x)\).  It is the product of the analytic
multiplicity, the geometric defect, and a graded residue-field degree.  The
factor is needed over a general ground field and for nonreduced formal
models; deleting it would make the formula false.

\begin{definition}[The defect factor and the fundamental cycle]
\label{def:ch17-defect}
\index{defect factor}\index{fundamental cycle!of a formal scheme}
Let \(x\in X\) reduce to \(\xi\in\mathfrak X_s\) with
\(V=\overline{\{\xi\}}\) of dimension \(n\).  Since \(X\) is boundaryless,
\(V\) is proper over \(\widetilde K^1\).  Let \(E=\widetilde K^1(V)\) be its
function field and \(F=E\otimes_{\widetilde K^1}\widetilde K\) the
associated graded field.  Set
\[
 \lambda_{\mathfrak X}(x)=[\widetilde{\cH(x)}:F]
 =\bigl(|\cH(x)^\times|:|K^\times|\bigr)\cdot
   \bigl[\widetilde{\cH(x)}^1:E\bigr],
\]
let \(m(x)\) be the multiplicity of the irreducible component of \(X\)
containing \(x\), and let \(\mathrm{def}_K(x)\) be the geometric defect.
The \emph{defect factor} is
\[
 \epsilon_{\mathfrak X}(x)
 =m(x)\cdot\mathrm{def}_K(x)\cdot\lambda_{\mathfrak X}(x).
\]
If \(\mathfrak X\) is proper, its \emph{fundamental cycle} is the pure
\(n\)-dimensional cycle supported on \(\mathfrak X_s\) given by
\[
 \operatorname{cycle}(\mathfrak X)
 =\sum_x\epsilon_{\mathfrak X}(x)\,[V_x],
\]
the sum being over the points \(x\in X\) whose reduction has
\(n\)-dimensional Zariski closure \(V_x\).  All coefficients are strictly
positive, and every component of \(\mathfrak X_s\) occurs, because the
reduction map is surjective (\Cref{rem:berk-reduction-surjective}).
\end{definition}

\begin{warning}{The fundamental cycle is not always the special fibre}
\(\operatorname{cycle}(\mathfrak X)\) and the cycle \([\mathfrak X_s]\) of
the special fibre agree in the two situations we shall actually use---the
geometrically reduced case (\Cref{prop:ch17-defect-reduced}) and the
discretely valued relatively normal case
(\Cref{prop:ch17-defect-dvr})---but they differ in general.  The three
factors measure three different things: \(m(x)\) measures nonreducedness of
\(X\) itself, \(\mathrm{def}_K(x)\) measures failure of \(K\) to be stable
at \(x\), and \(\lambda_{\mathfrak X}(x)\) measures both the ramification
index \((|\cH(x)^\times|:|K^\times|)\) and the residue degree of
\(\widetilde{\cH(x)}^1\) over the function field of \(V\).  A statement of
\Cref{thm:cld-formal-intersection} with \([\mathfrak X_s]\) in place of
\(\operatorname{cycle}(\mathfrak X)\) would be false, and no argument below
may silently substitute one for the other.
\end{warning}

\begin{theorem}
\label{thm:cld-formal-intersection}
Under the preceding hypotheses,
\[
 \bigwedge_{i=1}^n c_1(\overline L_i)
   =\sum_x \epsilon_{\mathfrak X}(x)
       \left(\int_{V_x}
          c_1(\mathfrak L_{1,s})\cdots
          c_1(\mathfrak L_{n,s})\right)\delta_x.
\]
The sum is over the locally finite set of analytic points specializing to
generic points of components of \(\mathfrak X_s\).  If \(\mathfrak X\) is
proper, its total mass is the intersection of the \(\mathfrak L_{i,s}\)
with CLD's fundamental cycle of \(\mathfrak X\).
\end{theorem}
\begin{proof}
The assertion is local on $\mathfrak X$, so choose a formal affine chart and
a tropicalization on which the formal line bundles are represented by
integer-affine transition functions.  The first Chern delta-current of
$\overline L_i$ is the corner current of that transition datum.  Intersect
the $n$ corner currents using a common subdivision.  A top-dimensional
intersection can occur only at a calibrated point $x$ specializing to the
generic point of a component $V_x$ of the special fibre.

At such a point, pass to Temkin's graded reduction.  The initial forms of
the transition functions define the residual line bundles
$\mathfrak L_{i,s}|_{V_x}$.  The stable-intersection multiplicity of their
corner currents is their residual intersection degree
\[
             \int_{V_x}c_1(\mathfrak L_{1,s})\cdots
                            c_1(\mathfrak L_{n,s}).
\]
The comparison between lattice-normalized tropical volume and the residual
degree contributes three finite factors: the length of the analytic
component, the geometric defect, and the graded residue-field degree.  Their
product is, by definition, $\epsilon_{\mathfrak X}(x)$.  Summing the local
corner products therefore gives the formula in the statement.  Common
refinements show that the computation is independent of the chart and that
the local sums glue.  Local finiteness of the special-fibre components gives
local finiteness of the measure.

If $\mathfrak X$ is proper, integrate the formula.  The residual projection
formula sums the graded local degrees, with the same defect factors, to the
intersection of the $\mathfrak L_{i,s}$ against the fundamental cycle of the
special fibre.  This proves the mass assertion.  The graded-reduction
comparison and its defect ledger are
\cite[Theorem~18.5.5]{CLD25}.  In the geometrically reduced case the defect
and graded residue factors are one; over a discrete valuation ring the
remaining analytic length is the usual component multiplicity.
\end{proof}

\begin{proof}[Second proof of the support and mass statements]
The support statement follows from \Cref{thm:ch17-atomic-support} and
\Cref{lem:ch17-residual-formal}: if \(x\in X\) reduces to \(\xi\) with
\(V=\overline{\{\xi\}}\), then \((V,\mathfrak L_i|_V)\) is a premodel of
\(\widetilde{L_{i,x}}\), so the simultaneous residual dimension of
definition at \(x\) is at most \(\dim V\); if \(\dim V<n\) the measure
vanishes near \(x\).

For the mass at a point \(x\) with \(\dim V_x=n\), start from
\eqref{eq:ch17-atomic-mass}.  Let \(\Lambda_i\) be the graded line bundle on
\(P_{E/\widetilde K^1}\) with premodel \((V,\mathfrak L_i|_V)\) and let
\(M_i\) be its inverse image on \(P_{F/\widetilde K}\).  By
\Cref{lem:ch17-residual-formal} the inverse image of \(M_i\) on
\(P_{\widetilde{\cH(x)}/\widetilde K}\) is \(\widetilde{L_{i,x}}\).  The
graded degree is therefore computed in three steps:
\[
\begin{split}
 \int_{P_{\widetilde{\cH(x)}/\widetilde K}}
   c_1(\widetilde{L_{1,x}})\cdots c_1(\widetilde{L_{n,x}})
 &=\lambda_{\mathfrak X}(x)
   \int_{P_{F/\widetilde K}}c_1(M_1)\cdots c_1(M_n)\\
 &=\lambda_{\mathfrak X}(x)
   \int_{P_{E/\widetilde K^1}}c_1(\Lambda_1)\cdots c_1(\Lambda_n)\\
 &=\lambda_{\mathfrak X}(x)
   \int_{V_x}c_1(\mathfrak L_{1,s})\cdots c_1(\mathfrak L_{n,s}),
\end{split}
\]
where the first equality is the behaviour of a graded degree under a finite
field extension of degree \(\lambda_{\mathfrak X}(x)=[\widetilde{\cH(x)}:F]\),
the second is invariance of the degree under the scalar extension
\(E\rightsquigarrow F=E\otimes_{\widetilde K^1}\widetilde K\), and the third
is the definition of the ungraded intersection number of a premodel.
Multiplying by \(m(x)\mathrm{def}_K(x)\) gives the asserted mass.  Summing
over \(x\) and using \Cref{def:ch17-defect} gives the mass assertion for
proper \(\mathfrak X\).
\end{proof}

\subsection{When the defect factors disappear}
\label{sec:ch17-defect-vanishes}

Two situations in which \(\epsilon_{\mathfrak X}(x)\) simplifies cover every
application in this book.

\begin{proposition}
\label{prop:ch17-defect-reduced}
Let \(\mathfrak X\) be as above and let \(x\in X\) reduce to the generic
point \(\xi\) of an irreducible component \(V\) of \(\mathfrak X_s\).
\begin{enumerate}
\item\label{itm:ch17-def:m}
 If \(\mathfrak X\) is reduced then \(m(x)=1\).
\item\label{itm:ch17-def:lambda}
 If \(\mathfrak X\) and \(\mathfrak X_s\) are reduced then
 \(\lambda_{\mathfrak X}(x)=1\); moreover \(x\) is the unique preimage of
 \(\xi\), \(|\cH(x)^\times|=|K^\times|\), and
 \(\widetilde{\cH(x)}^1=\widetilde K^1(V)\).
\item\label{itm:ch17-def:def}
 If \(\mathfrak X\) and \(\mathfrak X_s\) are geometrically reduced then
 \(\mathrm{def}_K(x)=1\).
\end{enumerate}
Consequently, if \(\mathfrak X\) and \(\mathfrak X_s\) are geometrically
reduced and \(\mathfrak X\) is proper, then
\(\operatorname{cycle}(\mathfrak X)=[\mathfrak X_s]\).
\end{proposition}

\begin{proof}
\cref{itm:ch17-def:m} is immediate: the generic multiplicity of a reduced
space is one.

\cref{itm:ch17-def:lambda}.  The assertion is local, so assume
\(\mathfrak X=\Spf A\) with \(A\) a topologically finitely presented flat
\(K^\circ\)-algebra such that \(A_{\widetilde K^1}\) is integral, and
\(X=\mathcal M(A_K)\).  Since \(A_K\) and \(A_{\widetilde K^1}\) are
reduced, there is a presentation \(K\{T_1,\dots,T_n\}\to A_K\) whose
quotient norm is the spectral norm of \(A_K\)
\cite[Proposition~1.1]{BoschLutkebohmert1993}.  Hence the spectral norm of
\(A_K\) takes values in \(|K|\), \(A_K^\circ=A\), and
\(A_K^{\circ\circ}=K^{\circ\circ}\cdot A\); so
\(A\otimes_{K^\circ}\widetilde K^1\to\widetilde{A_K}^1\) is an isomorphism.
The structure theory of affinoid reductions
(\Cref{sec:berk-reduction}) then gives that the generic point of \(V\) has a
unique preimage in \(X\), which is therefore \(x\), that \(\{x\}\) is the
Shilov boundary of \(A_K\), that the corresponding seminorm is the spectral
norm---whence \(|\cH(x)^\times|=|K^\times|\)---and that
\(\widetilde{\cH(x)}^1=\widetilde K^1(V)\).  Both factors in
\(\lambda_{\mathfrak X}(x)\) are then \(1\).

\cref{itm:ch17-def:def}.  Reduce as above to \(\mathfrak X=\Spf A\), let
\(K'\) be the completion of an algebraic closure of \(K\), and put
\(\mathfrak Y=\mathfrak X\widehat\otimes_{K^\circ}K'^\circ\) with generic
fibre \(Y=X_{K'}\).  Let \(W_1,\dots,W_r\) be the components of
\(\mathfrak Y_s\) dominating \(V\), with generic points \(\xi_i\).  By
\cref{itm:ch17-def:lambda} applied over \(K'\), each \(\xi_i\) has a unique
preimage \(y_i\in Y\), and the \(y_i\) are exactly the preimages of \(x\).
Since \(\mathfrak X_s\) is geometrically reduced, the algebra
\(\widetilde K^1(V)\otimes_{\widetilde K^1}\widetilde{K'}^1\) is the product
of the \(\widetilde{K'}^1(W_i)\) and has length \(r\); so the graded length
\(\widetilde\ell\) of
\(\widetilde{\cH(x)}\otimes_{\widetilde K}\widetilde{K'}\) equals \(r\).
On the other hand \(Y\) is reduced, so \(\cH(x)\widehat\otimes_KK'\) is the
product of the \(\cH(y_i)\) and has length \(\ell=r\).  Finally
\(|K'^\times|\) is divisible, so the ramification contribution
\(\widetilde{d^e}\) is \(1\).  Since
\(\mathrm{def}_K(x)=\ell\,\widetilde{d^e}/\widetilde\ell\), we get
\(\mathrm{def}_K(x)=1\).

The final assertion is now \Cref{def:ch17-defect}: all three factors equal
\(1\), and the components of \(\mathfrak X_s\) are reduced.  This is
\cite[Proposition~18.5.6 and Corollary~18.5.7]{CLD25}.
\end{proof}

\begin{proposition}
\label{prop:ch17-defect-dvr}
Suppose the valuation of \(K\) is discrete and write \(\widetilde\cdot\) for
the classical reduction.  Let \(\mathfrak X\) be a formal
\(K^\circ\)-scheme, locally topologically of finite presentation, relatively
normal in its generic fibre \(X\), which we assume reduced.  Let \(\xi\) be
a generic point of \(\mathfrak X_s\).  Then:
\begin{enumerate}
\item\label{itm:ch17-dvr:unique}
 \(\xi\) has a unique preimage \(x\in X\);
\item\label{itm:ch17-dvr:mult}
 the multiplicity of the component \(\overline{\{\xi\}}\) of
 \(\mathfrak X_s\), that is, the length of the artinian local ring
 \(\cO_{\mathfrak X_s,\xi}\), equals the index
 \(\bigl(|\cH(x)^\times|:|K^\times|\bigr)\);
\item\label{itm:ch17-dvr:residue}
 \(\widetilde{\cH(x)}^1=\widetilde K^1(\xi)\);
\item\label{itm:ch17-dvr:defect}
 \(m(x)=\mathrm{def}_K(x)=1\) and \(\lambda_{\mathfrak X}(x)\) is the
 multiplicity of \(\overline{\{\xi\}}\) in \(\mathfrak X_s\).
\end{enumerate}
Consequently, if \(\mathfrak X\) is proper then
\(\operatorname{cycle}(\mathfrak X)=[\mathfrak X_s]\).
\end{proposition}

\begin{proof}
The question is local, so let \(\mathfrak X=\Spf A\) with \(A\) flat,
topologically of finite presentation, and integrally closed in \(A_K\); we
may assume \(\mathfrak X\) equidimensional of dimension \(d\).  Then
\(X=\mathcal M(A_K)\) is equidimensional of dimension \(d\) and reduced, so
its non-normal locus \(Z\) is a Zariski-closed subset of dimension \(<d\).
The point \(x\) we shall construct has
\(\operatorname{tr.deg}(\widetilde{\cH(x)}/\widetilde K)=d\), so it does not
lie in \(Z\), and \(\xi\) does not lie in the image
\(\operatorname{red}(Z)\)---a constructible, indeed closed, set of dimension
\(<d\) by Chevalley's theorem and finiteness of the reduction of an affinoid.
Shrinking, we may assume \(X\) integral and normal; then \(A\) is a normal
domain with fraction field \(\mathcal K\).

The scheme \(\mathfrak X_s\) is \(\Spec A_{\widetilde K}\), and \(\xi\)
corresponds to a height-one point of the noetherian normal scheme
\(\Spec A\).  Its local ring \(\cO_{\Spec A,\xi}\) is therefore a discrete
valuation ring of \(\mathcal K\) dominating \(K^\circ\).  Normalize the
corresponding absolute value \(|\cdot|:\mathcal K\to\mathbf R_+\) so that it
restricts to the given one on \(K^\circ\).  Composing with
\(A_K\hookrightarrow\mathcal K\) gives a bounded multiplicative seminorm on
\(A_K\), that is, a point \(x\in X\); by construction the evaluation at
\(x\) is an isometric embedding of \(\mathcal K\) into \(\cH(x)\) with dense
image, so the residue field of \(\cH(x)\) is \(\widetilde K^1(\xi)\), which
is \cref{itm:ch17-dvr:residue}, and its value group is
\(|\mathcal K^\times|=|K^\times|^{1/e}\) where \(e\) is the ramification
index of \(\cO_{\Spec A,\xi}\) over \(K^\circ\).  That index is the length
of \(\cO_{\Spec A,\xi}/(\pi)\), which is the length of
\(\cO_{\mathfrak X_s,\xi}\); this is \cref{itm:ch17-dvr:mult}.

For \cref{itm:ch17-dvr:unique}, let \(y\in X\) reduce to \(\xi\).  Then every
\(a\in A\) with \(a(\xi)=0\) satisfies \(|a(y)|<1\), so the valuation ring
of \(\mathcal K\) induced by \(y\) dominates \(\cO_{\Spec A,\xi}\); the
latter being a valuation ring, the two coincide.  As the restrictions to
\(K\) agree, \(y=x\).

\cref{itm:ch17-dvr:defect}: \(m(x)=1\) because \(X\) is reduced;
\(\mathrm{def}_K(x)=1\) because a discretely valued field is stable;
and
\(\lambda_{\mathfrak X}(x)
 =\bigl(|\cH(x)^\times|:|K^\times|\bigr)\cdot
  [\widetilde{\cH(x)}^1:\widetilde K^1(\xi)]\)
equals the multiplicity by \cref{itm:ch17-dvr:mult} and
\cref{itm:ch17-dvr:residue}.  The last assertion follows from
\Cref{def:ch17-defect}.  This is \cite[Proposition~18.5.8]{CLD25}.
\end{proof}

\begin{remark}[Where the two hypotheses are used]
\label{rem:ch17-dvr-hypotheses}
Discreteness of the valuation entered twice: to make
\(\cO_{\Spec A,\xi}\) a \emph{discrete} valuation ring, so that the
ramification index is a length; and to make \(K\) stable, so that the
geometric defect is one.  Relative normality entered once, but decisively,
in \cref{itm:ch17-dvr:unique}: without it a generic point of the special
fibre may have several preimages, and the atom \(b_Y\delta_{x_Y}\) of
\eqref{eq:cld-chambert-loir-formula} would have to be replaced by a sum over
the fibre.  Both hypotheses are exactly the ones under which the
Chambert--Loir measure of \Cref{sec:metrics-cl-measure} was defined.
\end{remark}

\subsection{A formal metric with a multiple component}
\label{sec:ch17-multiple-fibre}

The factor \(\lambda_{\mathfrak X}(x)\) of \Cref{def:ch17-defect} is
invisible in \Cref{ex:ch17-check-p1} and
\Cref{ex:ch17-check-toric-surface}, where every multiplicity is one.  The
following example makes it visible, computes it in two independent ways, and
checks the resulting special-fibre intersection numbers against a slope
computation.

\begin{example}[Two blowups, a component of multiplicity two, and all four
intersection numbers]
\label{ex:ch17-multiple-fibre}
\index{special fibre!multiple component}\index{defect factor!computed}
Let \(K\) be discretely valued with uniformizer \(\pi\), \(v_K(\pi)=1\), let
\(X=\mathbf P^1_K\) with coordinate \(T\), and let \(L=\cO(1)\), so
\(\deg L=1\).

\emph{The model.}  Let \(\mathcal X'\to\mathbf P^1_{K^\circ}\) be the blowup
of the closed point \(p=\{\pi=T=0\}\) of the special fibre, as in
\Cref{ex:metrics-p1-blowup}, with
\(\mathcal X'_s=\widetilde C'+E_1\) and
\((E_1^2)=(\widetilde C'^2)=-1\), \((\widetilde C'\cdot E_1)=1\).  Let
\(\rho:\mathcal X\to\mathcal X'\) be the blowup of the point
\(p'=\widetilde C'\cap E_1\), with exceptional divisor \(E_2\).  Writing
\(C\) and \(E\) for the strict transforms of \(\widetilde C'\) and \(E_1\),
\[
 \rho^*\widetilde C'=C+E_2,\qquad \rho^*E_1=E+E_2,
\]
so that
\begin{equation}
\label{eq:ch17-multiple-fibre-cycle}
 \mathcal X_s=\rho^*\mathcal X'_s=C+E+2E_2 .
\end{equation}
The component \(E_2\) has multiplicity \(b_{E_2}=2\).  The intersection
numbers on the regular surface \(\mathcal X\) are
\[
 (E_2^2)=-1,\quad (C^2)=-2,\quad (E^2)=-2,\quad
 (C\cdot E_2)=(E\cdot E_2)=1,\quad (C\cdot E)=0,
\]
the values for \(C\) and \(E\) coming from
\((\rho^*D-E_2)^2=(D^2)-1\) and
\((C\cdot E)=(\widetilde C'\cdot E_1)+(E_2^2)=1-1=0\).  As a check,
\eqref{eq:ch17-multiple-fibre-cycle} gives
\((\mathcal X_s\cdot C)=-2+0+2=0\),
\((\mathcal X_s\cdot E)=0-2+2=0\) and
\((\mathcal X_s\cdot E_2)=1+1-2=0\), as it must, \(\mathcal X_s\) being a
principal divisor.

\emph{The three divisorial points.}  In the chart of the first blowup with
coordinates \((T,u)\) and \(\pi=Tu\), the point \(p'\) is \(T=u=0\); the
second blowup has a chart with \(u=Tw\), so \(\pi=T^2w\) and
\(E_2=\{T=0\}\).  Hence
\[
 \ord_{E_2}(\pi)=2,\qquad \ord_{E_2}(T)=1,
\]
and the normalized divisorial point \(x_{E_2}=b_{E_2}^{-1}\ord_{E_2}\)
satisfies
\[
 v_{x_{E_2}}(\pi)=1,\qquad v_{x_{E_2}}(T)=\tfrac12 .
\]
So \(x_{E_2}=\zeta_{0,|\pi|^{1/2}}\), and on the standard apartment
parametrized by \(t=v(T)\) the three divisorial points sit at
\[
 x_C\ (t=0),\qquad x_{E_2}\ (t=\tfrac12),\qquad x_E\ (t=1),
\]
in that order, matching the dual graph of \(\mathcal X_s\): \(C\) and \(E\)
each meet \(E_2\) and do not meet each other.

\emph{The defect factor, computed twice.}  First, from
\Cref{prop:ch17-defect-dvr}\,\cref{itm:ch17-dvr:defect}:
\(\lambda_{\mathfrak X}(x_{E_2})\) is the multiplicity of \(E_2\) in
\(\mathcal X_s\), which is \(2\) by
\eqref{eq:ch17-multiple-fibre-cycle}.  Second, directly from
\Cref{def:ch17-defect}: the value group of \(\cH(x_{E_2})\) contains
\(|T(x_{E_2})|=|\pi|^{1/2}\), so
\(\bigl(|\cH(x_{E_2})^\times|:|K^\times|\bigr)\ge2\); and it is exactly
\(2\), because \(x_{E_2}\) is the monomial point of the disc of radius
\(|\pi|^{1/2}\), whose value group is \(|K^\times|^{1/2}\).  The residue
degree \([\widetilde{\cH(x_{E_2})}^1:\widetilde K(E_2)]\) is \(1\).  So
\(\lambda_{\mathfrak X}(x_{E_2})=2\cdot1=2\), and
\(m=\mathrm{def}_K=1\), giving \(\epsilon_{\mathfrak X}(x_{E_2})=2\).  The
two computations agree, and the second one used no model at all: it read the
factor off the value group of the Berkovich point.

\emph{A model line bundle and its special-fibre degrees.}  Let \(h\) be the
total pullback of \(\cO(1)\) from \(\mathbf P^1_{K^\circ}\) and put
\(\mathcal L=h\otimes\cO_{\mathcal X}(E_2)\), a model of \(L=\cO(1)\).  Since
\(C\to\mathbf P^1_{\widetilde K}\) is an isomorphism and \(E,E_2\) are
contracted, \(\deg(h|_C)=1\) and \(\deg(h|_E)=\deg(h|_{E_2})=0\).  Therefore
\begin{equation}
\label{eq:ch17-multiple-fibre-degrees}
 \deg(\mathcal L|_C)=1+(E_2\cdot C)=2,\quad
 \deg(\mathcal L|_E)=(E_2\cdot E)=1,\quad
 \deg(\mathcal L|_{E_2})=(E_2^2)=-1 .
\end{equation}

\emph{The measure from \Cref{thm:cld-formal-intersection}.}  Multiplying each
degree by its defect factor \(\epsilon_{\mathfrak X}=b_Y\),
\begin{equation}
\label{eq:ch17-multiple-fibre-measure}
 c_1(\overline L)
 =1\cdot2\,\delta_{x_C}+1\cdot1\,\delta_{x_E}+2\cdot(-1)\,\delta_{x_{E_2}}
 =2\,\delta_{x_C}+\delta_{x_E}-2\,\delta_{x_{E_2}} .
\end{equation}
Total mass \(2+1-2=1=\deg\cO(1)\), as \Cref{cor:ch17-cl-mass} requires.

\emph{The same measure from slopes.}  The metric of \(\mathcal L\) differs
from the Weil metric by the model function \(\varphi_{E_2}\), whose values
at the divisorial points are \(a_Y/b_Y\), that is,
\[
 \varphi_{E_2}(x_C)=0,\qquad
 \varphi_{E_2}(x_{E_2})=\tfrac12,\qquad
 \varphi_{E_2}(x_E)=0,
\]
and which is affine on each edge of the skeleton and constant on every
branch off it.  On \([0,\tfrac12]\) its slope is \(+1\) and on
\([\tfrac12,1]\) it is \(-1\).  Summing outgoing slopes
(\Cref{prop:skeletal-curvature-curve}),
\[
 \dprimeddprime[\varphi_{E_2}]
 =\delta_{x_C}-2\,\delta_{x_{E_2}}+\delta_{x_E},
\]
of total mass \(0\).  Adding the Weil curvature \(\delta_{x_C}\) of
\Cref{ex:ch17-weil-curvature}\,(a) and using
\Cref{prop:cld-c1-metric-change} reproduces
\eqref{eq:ch17-multiple-fibre-measure} exactly.

\emph{What has been checked.}  Three separate normalizations had to be right
for the two computations to agree: the multiplicity \(b_{E_2}=2\) in the
Chambert--Loir formula; the normalization
\(x_{E_2}=b_{E_2}^{-1}\ord_{E_2}\) of the divisorial point, which is what
places the atom at \(t=\tfrac12\) rather than at \(t=1\); and the identity
\(\varphi_D(x_Y)=a_Y/b_Y\) of
\Cref{prop:metrics-divisor-model-function}, which is what makes the slope of
\(\varphi_{E_2}\) equal to \(1\) rather than \(2\).  Using \(\ord_{E_2}\)
instead of \(x_{E_2}\), or dropping \(b_{E_2}\), would break the agreement
and also the total mass.  This is the warning after
\eqref{eq:metrics-cl-pairing}, seen in the smallest example in which it
bites.

\emph{Positivity.}  The metric is \emph{not} semipositive: the degree
\(-1\) on \(E_2\) shows that \(\mathcal L\) is not relatively nef, and
correspondingly \eqref{eq:ch17-multiple-fibre-measure} has a negative atom.
By \Cref{thm:ch17-psh-residual-nef} the residual line bundle at
\(x_{E_2}\) is not nef, which one can see directly: its premodel is
\((E_2,\cO(E_2)|_{E_2})=(\mathbf P^1_{\widetilde K},\cO(-1))\).
\end{example}

\section{Recovery of Chambert--Loir measures}
\label{sec:ch17-cl}

We now give the form used most often in global pluripotential theory.  This
is the theorem announced in \Cref{sec:ch17-payoff}: it identifies a
Bedford--Taylor product of curvature currents, constructed by purely local
analytic means, with the atomic measure that \Cref{chap:metrics} defined by
intersection theory on a model.

\begin{theorem}
\label{thm:cld-recovers-cl}
Assume that \(K\) is discretely valued, \(X\) is a proper reduced
\(n\)-dimensional \(K\)-scheme, and \(\mathcal X\) is a proper flat model
that is relatively normal in \(X\).  Write
\[
                         \mathcal X_s=\sum_Y b_Y Y.
\]
Let \(\mathcal L_i\) extend \(L_i^{\otimes m_i}\), with \(m_i>0\), and
give \(L_i\) the corresponding rational formal metrics.  Then the CLD
curvature product is
\begin{equation}
 \bigwedge_{i=1}^n c_1(\overline L_i)
   =\frac{1}{m_1\cdots m_n}
     \sum_Y b_Y
       \bigl(c_1(\mathcal L_1|_Y)\cdots
             c_1(\mathcal L_n|_Y)\bigr)\,\delta_{x_Y},
\label{eq:cld-chambert-loir-formula}
\end{equation}
where \(x_Y\) is the unique divisorial point reducing to the generic point
of \(Y\).  Consequently this measure is exactly the mixed
Chambert--Loir measure constructed from the same model.
\end{theorem}

\begin{proof}
Pass to the formal completion of \(\mathcal X\) along its special fibre.
Relative normality gives one analytic point \(x_Y\) over the generic point of
each component.  In the discretely valued calculation, the product of the
defect factor in \Cref{thm:cld-formal-intersection} and the fundamental
cycle coefficient is the scheme-theoretic multiplicity \(b_Y\).  Applying
that theorem to \(\mathcal L_i\), and using multilinearity under the tensor
powers \(m_i\), gives \eqref{eq:cld-chambert-loir-formula}.  Chambert--Loir's
definition is the right-hand side, so the two measures agree.  This final
identification is precisely \cite[Proposition~18.5.8]{CLD25}; it also agrees
with \cite{CL06}.
\end{proof}

\begin{proof}[Detailed proof]
We give the argument in five steps, so that every hypothesis is visible at
the point where it is used.

\emph{Step 1: passage to the formal model.}  Let
\(\mathfrak X=\widehat{\mathcal X}\) be the completion of \(\mathcal X\)
along its special fibre.  It is a flat formal \(K^\circ\)-scheme, locally
topologically of finite presentation, with
\(\mathfrak X_\eta=X^{\An}\) and \(\mathfrak X_s=\mathcal X_s\); properness
of \(\mathcal X\) makes \(X^{\An}\) proper, hence boundaryless
(\Cref{thm:berk-proper-compact}), and pure of dimension \(n\).  Each
\(\mathcal L_i\) completes to a line bundle \(\mathfrak L_i\) on
\(\mathfrak X\), and the formal metric of \(\mathfrak L_i\) on
\((\mathfrak L_i)_\eta=L_i^{\otimes m_i,\An}\) is the analytification of the
algebraic model metric of \(\mathcal L_i\)---this is the equivalence of
algebraic model metrics and formal metrics on a proper scheme,
\Cref{thm:gm-metric-comparison}\,(2).  The rational formal metric on
\(L_i\) is the one whose \(m_i\)-th tensor power is that metric.  Relative
normality of \(\mathcal X\) in \(X\) passes to \(\mathfrak X\), because it
is a statement about the integral closure of the structure sheaf in the
generic fibre and completion along an ideal preserves it.

\emph{Step 2: the support and the divisorial points.}  Write
\(\mathcal X_s=\sum_Yb_YY\).  By
\Cref{prop:ch17-defect-dvr}\,\cref{itm:ch17-dvr:unique}, each generic point
\(\eta_Y\) of \(\mathcal X_s\) has a unique preimage \(x_Y\in X^{\An}\) under
the reduction map.  By
\Cref{prop:ch17-defect-dvr}\,\cref{itm:ch17-dvr:mult} and
\cref{itm:ch17-dvr:residue}, the point \(x_Y\) has value group of index
\(b_Y\) in \(|K^\times|\) and residue field \(\widetilde K(\eta_Y)\); by
\Cref{prop:component-valuation} the corresponding valuation is
\(b_Y^{-1}\ord_Y\), which is the normalization of the divisorial point used
in \Cref{sec:metrics-cl-measure}.  So the points \(x_Y\) appearing here are
literally the points \(x_Y\) of \Cref{chap:metrics}, with no rescaling.

\emph{Step 3: the measure for the integral metrics.}  Apply
\Cref{thm:cld-formal-intersection} to \(\mathfrak X\) and the
\(\mathfrak L_i\).  Its support is contained in the set of points reducing
to generic points of components of \(\mathcal X_s\), which by Step~2 is
\(\{x_Y\}_Y\), a finite set.  The mass at \(x_Y\) is
\[
 \epsilon_{\mathfrak X}(x_Y)\int_Y c_1(\mathfrak L_{1,s})\cdots
   c_1(\mathfrak L_{n,s})
 =b_Y\bigl(c_1(\mathcal L_1|_Y)\cdots c_1(\mathcal L_n|_Y)\bigr),
\]
because \(\epsilon_{\mathfrak X}(x_Y)=b_Y\) by
\Cref{prop:ch17-defect-dvr}\,\cref{itm:ch17-dvr:defect}, and because the
restriction of \(\mathfrak L_{i,s}\) to \(Y\) is the restriction of
\(\mathcal L_i\) to \(Y\).  Hence
\begin{equation}
\label{eq:ch17-cl-integral-case}
 \bigwedge_{i=1}^nc_1\bigl(L_i^{\otimes m_i},\|\cdot\|_{\mathcal L_i}\bigr)
 =\sum_Yb_Y\bigl(c_1(\mathcal L_1|_Y)\cdots c_1(\mathcal L_n|_Y)\bigr)
  \delta_{x_Y}.
\end{equation}

\emph{Step 4: descent through the tensor powers.}  The curvature current is
additive in tensor powers by \Cref{cor:ch17-c1-functorial}, so
\(c_1(L_i^{\otimes m_i},\|\cdot\|_{\mathcal L_i})
 =m_i\,c_1(L_i,\|\cdot\|_i)\) where \(\|\cdot\|_i\) is the rational formal
metric.  The mixed product is multilinear
(\Cref{prop:ch17-product-welldefined}), so dividing
\eqref{eq:ch17-cl-integral-case} by \(m_1\cdots m_n\) gives
\eqref{eq:cld-chambert-loir-formula}.  Note that this step is where the
extension of the Bedford--Taylor product to \emph{differences} of psh
potentials is needed: the metrics \(\|\cdot\|_i\) need not be psh, and the
multilinearity being invoked is that of \Cref{thm:ch16-dpsh-product}, not a
statement about positive currents.

\emph{Step 5: comparison with the Chambert--Loir measure.}  The right-hand
side of \eqref{eq:cld-chambert-loir-formula} is, verbatim, the mixed model
measure of \Cref{sec:metrics-cl-measure}, with the same normal model, the
same tensor powers, the same multiplicities \(b_Y\), the same divisorial
points \(x_Y\), and the same special-fibre intersection numbers.  Hence the
two measures coincide.  Independence of the chosen model is then available
from either side: from the Chambert--Loir side by
\Cref{thm:cl-model-properties}, and from the CLD side because the left-hand
side of \eqref{eq:cld-chambert-loir-formula} was defined without reference
to a model at all.  It is worth recording that this gives a second, purely
local proof of the model-independence asserted in
\Cref{thm:cl-model-properties}, whose original proof went through the
projection formula on a dominating model.
\end{proof}

\begin{remark}[What the theorem does and does not say]
\label{rem:ch17-cl-scope}
Two delimitations.

\emph{(i) It is a statement about model metrics.}  Both sides of
\eqref{eq:cld-chambert-loir-formula} are defined for arbitrary
\(\mathbf Q\)-model metrics, with no positivity assumption; the measure is
signed in general.  The extension to \emph{continuous semipositive} metrics
is \Cref{thm:ch17-cl-continuous} below and requires a limit argument on both
sides.  Beyond that---for arbitrary bounded or finite-energy psh
metrics---the two theories genuinely diverge, and the delimitation is the
one drawn in \Cref{sec:ch16-local-vs-global}.

\emph{(ii) It is a statement over a discretely valued field.}  Over a
general complete field there is no Chow-theoretic degree on the special
fibre of an admissible formal model of arbitrary rank, which is exactly why
\Cref{sec:metrics-cl-measure} restricted to the discretely valued case.  The
CLD side, by contrast, is defined in complete generality; the correct
general statement is \Cref{thm:cld-formal-intersection}, in which the
Chow-theoretic degree on \(\mathfrak X_s\) is replaced by a graded degree
and the multiplicity \(b_Y\) by the defect factor
\(\epsilon_{\mathfrak X}\).
\end{remark}

\begin{corollary}
\label{cor:ch17-cl-mass}
Under \Cref{thm:cld-recovers-cl}, the total mass is
\[
                    (L_1\cdots L_n)_X.
\]
If each \(\mathcal L_i\) is relatively nef, the measure is positive.  If
\(\mathcal X_s\) is irreducible and reduced, the measure for a single
\(\overline L\) is
\[
                         c_1(\overline L)^n=(L^n)\delta_{x_{\mathcal X_s}}.
\]
\end{corollary}
\begin{proof}
Flatness and the specialization invariance of intersection numbers identify
the sum in \eqref{eq:cld-chambert-loir-formula} with the generic-fibre
degree.  Relative nefness makes every special-fibre mixed intersection
nonnegative.  In the good-reduction case the sum has one term of
multiplicity one.
\end{proof}

\begin{example}[The standard model of the projective line]
Let \(K\) be discretely valued,
\(\mathcal X=\mathbf P^1_{K^\circ}\), and
\(\mathcal L=\mathcal O_{\mathcal X}(1)\).  Give
\(L=\mathcal O_{\mathbf P^1_K}(1)\) its formal metric.  The special fibre is
the single reduced component
\(Y=\mathbf P^1_{\widetilde K}\), so \(b_Y=1\), and

\[
            \deg_Y(\mathcal L|_Y)=1.
\]

Formula \eqref{eq:cld-chambert-loir-formula} therefore gives

\[
                 c_1(L,\|\cdot\|)=\delta_{x_Y},
\]

where \(x_Y\) is the Gauss point of \(\mathbf P^{1,\An}_K\).  Replacing
\(\mathcal L\) by \(\mathcal O_{\mathcal X}(d)\) multiplies the residual
degree by \(d\), hence

\[
          c_1\bigl(\mathcal O(d),\|\cdot\|_{\mathrm{form}}\bigr)
             =d\,\delta_{x_Y}.
\]

This calculation displays both normalization factors in the simplest case:
the component multiplicity is one and the entire mass is the degree of the
generic-fibre line bundle.
\end{example}

\subsection{Continuous semipositive metrics}
\label{sec:ch17-cl-continuous}

\Cref{thm:cld-recovers-cl} is a statement about model metrics.  The
Chambert--Loir measure was extended to continuous semipositive metrics in
\Cref{thm:cl-continuous-extension} by a Cauchy estimate on model
approximations.  The CLD product has its own continuity theorem,
\Cref{thm:cld-bt-continuity}.  The two extensions agree, and the proof is
short once both are available.

\begin{theorem}
\label{thm:ch17-cl-continuous}
Let \(K\) be discretely valued, let \(X\) be a proper reduced
\(n\)-dimensional \(K\)-scheme, and let \(L_1,\dots,L_n\) be line bundles on
\(X\), each with a continuous semipositive metric in the sense of
\Cref{sec:metrics-semipositive}.  Then the CLD mixed curvature product
\(\bigwedge_ic_1(\overline L_i)\), formed by the Bedford--Taylor recursion
of \Cref{chap:cld-bedford-taylor}, is defined and equals the
Chambert--Loir measure of \Cref{thm:cl-continuous-extension}.  Its total
mass is \((L_1\cdots L_n)_X\).
\end{theorem}

\begin{proof}
By definition of semipositivity there are, for each \(i\), semipositive
\(\mathbf Q\)-model metrics \(\|\cdot\|_{i,j}\) converging uniformly to
\(\|\cdot\|_i\) as \(j\to\infty\).  Each \(\|\cdot\|_{i,j}\) is a formal
metric that is psh: its residual bundles are nef by
\Cref{thm:ch17-psh-residual-nef}, applied through the relative nefness of
the corresponding model line bundle, which is what semipositivity of a model
metric means (\Cref{thm:model-semipositive-closed}).  A uniform limit of
psh metrics is psh by \Cref{prop:ch15-uniform-limit}, so
\(\|\cdot\|_i\) is a continuous psh metric and the recursive product of
\Cref{chap:cld-bedford-taylor} applies to it.

The continuity theorem \Cref{thm:cld-bt-continuity} gives
\[
 \bigwedge_ic_1(\overline L_{i,j})\longrightarrow\bigwedge_ic_1(\overline L_i)
\]
weakly, the potentials converging uniformly on the compact space
\(X^{\An}\).  By \Cref{thm:cld-recovers-cl}, the left-hand side is the
Chambert--Loir measure of the model metrics \(\|\cdot\|_{i,j}\); by
\Cref{thm:cl-continuous-extension}, those measures converge weakly to the
Chambert--Loir measure of the limit metrics.  Weak limits of Radon measures
on a compact Hausdorff space being unique, the two limits coincide.

The mass statement follows because each approximant has mass
\((L_1\cdots L_n)_X\) by \Cref{cor:ch17-cl-mass}, and mass is continuous
under weak convergence of positive measures on a compact space.
\end{proof}

\begin{warning}{Semipositivity is not decoration here}
The proof used semipositivity twice, and neither use can be removed by an
easy argument.  It was used to know that the limit metric is psh, so that
the Bedford--Taylor recursion applies at all; and it was used to know that
weak convergence of the approximating measures is convergence of
\emph{positive} measures, so that no mass escapes.  For a continuous metric
that is merely a uniform limit of arbitrary model metrics, the sequence
\(\bigwedge_ic_1(\overline L_{i,j})\) is a sequence of signed measures of
bounded total variation only under an extra hypothesis, and
\Cref{ex:ch17-weil-curvature}\,(b) already shows that the individual
measures need not be positive.  The correct general statement in that
regime is not a theorem of this chapter.
\end{warning}

\section{Checking the recovery against the model computations}
\label{sec:ch17-cl-checks}

\Cref{thm:cld-recovers-cl} asserts an equality between two computations that
were carried out independently: the Chambert--Loir side in
\Cref{sec:metrics-cl-measure}, from intersection numbers on models, and the
CLD side here, from local weights and tropical Monge--Amp\`ere.  A
comparison theorem of this kind is worth nothing unless one checks it on
cases where both sides have been computed to numbers.  We do that now, in
dimensions one and two, and we verify the total mass \((L^n)\) each time.

\subsection{The blown-up projective line}
\label{sec:ch17-check-p1}

\begin{example}[Verification in dimension one]
\label{ex:ch17-check-p1}
\index{Chambert--Loir measure!verified against CLD}
Let \(K\) be discretely valued with uniformizer \(\pi\), \(v_K(\pi)=1\), and
let \(X=\mathbf P^1_K\), \(L=\cO(1)\), so \((L)=\deg L=1\).  Let
\(\rho:\mathcal X'\to\mathbf P^1_{K^\circ}\) be the blowup of the closed
point \(p\) of the special fibre given by \(T\equiv0\), and let
\(\mathcal L'=\rho^*\cO(1)\otimes\cO_{\mathcal X'}(-E)\).

\emph{The Chambert--Loir side.}  \Cref{ex:metrics-cl-p1-blowup} computes
\(\mathcal X'_s=\widetilde C+E\) with \(b_{\widetilde C}=b_E=1\),
\[
 \deg(\mathcal L'|_E)=1,\qquad
 \deg(\mathcal L'|_{\widetilde C})=0,
\]
hence
\(c_1(L,\|\cdot\|_{\mathcal L'})=\delta_{x_E}\) with
\(x_E=\zeta_{0,|\pi|}\), of total mass \(1\).

\emph{The CLD side.}  \Cref{ex:metrics-p1-blowup} identifies the weight of
the metric \(\|\cdot\|_{\mathcal L'}\) in the frame \(Z_0\) as
\(\max\{-1,\log|T|\}\); this is precisely the metric
\(\|\cdot\|''\) of \Cref{ex:ch17-weil-curvature}\,(c), whose curvature was
computed there from slopes alone:
\[
 c_1(\overline L'')=\delta_{\zeta_{0,|\pi|}}.
\]

\emph{Verification.}  The two atoms sit at the same point
\(\zeta_{0,|\pi|}=x_E\), with the same mass \(1\), and
\[
 \int_{X^{\An}}c_1(\overline L'')=1=(L)=\deg\cO(1),
\]
as \Cref{cor:ch17-cl-mass} requires.  Note what has been checked: the CLD
computation used only the piecewise affine weight and the rule ``mass equals
slope jump''; the Chambert--Loir computation used only
\((E^2)=-1\) and \((\widetilde C\cdot E)=1\).  Their agreement is
\Cref{thm:cld-recovers-cl} for \(n=1\), and it also confirms
\eqref{eq:metrics-p1-blowup-laplacian}, the graph-Laplacian form of the same
identity.
\end{example}

\begin{example}[The semistable chain, all three atoms]
\label{ex:ch17-check-chain}
\Cref{ex:metrics-cl-curve-chain} treats the conic
\(\mathcal C^{(2)}=\{XY=\pi^2Z^2\}\) with semistable chain model
\(\widetilde{\mathcal C}\), special fibre the chain \(C_0,E,C_2\) with all
multiplicities one, skeleton \(\Gamma=[0,2]\), and \(L=\cO(1)|_C\) of degree
\(2\).  For a vertical divisor \(D=a_0C_0+a_1E+a_2C_2\) it computes
\begin{equation}
\label{eq:ch17-chain-cl}
 c_1(L,\|\cdot\|_{\mathcal L})
 =(1+a_1-a_0)\delta_{v_0}
 +(a_0+a_2-2a_1)\delta_{v_1}
 +(1+a_1-a_2)\delta_{v_2},
\end{equation}
which is \eqref{eq:metrics-cl-chain-measure}.

The CLD side is a slope computation on \(\Gamma\).  The weight of the metric
in the frame coming from \(\mathcal L_0\) is the piecewise affine function
\(\varphi\) with \(\varphi(i)=a_i\), affine on \([0,1]\) and on \([1,2]\),
added to the weight of the reference model metric \(\mathcal L_0\).  By
\Cref{prop:cld-c1-metric-change} the curvature is the curvature of the
reference metric plus \(\dprimeddprime[\varphi]\).  On a metric graph the
latter is the sum of outgoing slopes at each vertex
(\Cref{prop:skeletal-curvature-curve}), that is,
\[
 (a_1-a_0)\delta_{v_0}
 +\bigl[(a_0-a_1)+(a_2-a_1)\bigr]\delta_{v_1}
 +(a_1-a_2)\delta_{v_2},
\]
while the reference curvature is \(\delta_{v_0}+\delta_{v_2}\) because
\(\mathcal L_0\) has fibre degrees \(1,0,1\).  Adding gives exactly
\eqref{eq:ch17-chain-cl}.

\emph{Total mass.}  Summing the three coefficients,
\[
 (1+a_1-a_0)+(a_0+a_2-2a_1)+(1+a_1-a_2)=2=\deg_CL=(L),
\]
for every \((a_0,a_1,a_2)\), as \Cref{cor:ch17-cl-mass} and
\Cref{thm:ch17-degree} both require: the \(a_i\) cancel because the
Laplacian of a function on a compact graph has total mass zero, which is the
one-dimensional case of
\(\int_X\dprimeddprime[\varphi]=0\) in
\Cref{prop:ch17-mass-independent}.

\emph{Positivity.}  The measure is positive exactly when
\(a_0-a_1\le1\), \(a_0+a_2\ge2a_1\), \(a_2-a_1\le1\); this is the nef cone
computed in \Cref{ex:chain-nef-functions}, and by
\Cref{thm:ch17-psh-residual-nef} it is also exactly the locus where the
three residual line bundles are nef.  The reader may check the middle
condition against \Cref{cor:ch17-atoms-nonneg}: the atom at \(v_1\) is the
degree of \(\mathcal L\) on the contracted component \(E\), and its
nonnegativity is convexity of \(\varphi\) at the interior vertex.
\end{example}

\subsection{A surface: the toric blowup, with total mass \texorpdfstring{$(L^2)$}{(L2)}}
\label{sec:ch17-check-surface}

\begin{example}[Verification in dimension two]
\label{ex:ch17-check-toric-surface}
\index{Chambert--Loir measure!computed on a toric surface}
This is the example the plan for the chapter singles out, so we carry it
through both computations in full.

\emph{The data.}  Let \(K\) be discretely valued with \(v_K(\pi)=1\), let
\(X=\mathbf P^2_K\) with homogeneous coordinates \(Z_0,Z_1,Z_2\) and torus
coordinates \(x=Z_1/Z_0\), \(y=Z_2/Z_0\), and let \(L=\cO(1)\), so that
\((L^2)=1\).  Consider the Fubini--Study metric attached to the three
sections \(s_0=\pi Z_0\), \(s_1=Z_1\), \(s_2=Z_2\) with zero weights.

\emph{The CLD side.}  In the frame \(Z_0\) and the tropical coordinates
\(u=(v(x),v(y))\), the weight is
\[
 F(u)=\max\{-1,\,-u_1,\,-u_2\},
\]
by \eqref{eq:metrics-toric-weight}.  This is a convex piecewise affine
function of \(u\in\mathbf R^2\), so
\Cref{prop:ch16-tropically-convex} makes the metric tropically convex, hence
psh, and \Cref{thm:ch17-atomic-support} makes \(c_1(\overline L)^2\) a
locally finite atomic measure.  Its atoms sit at the points of the skeleton
over which \(F\) is not affine near which the subgradient image is
two-dimensional: that is, at the vertices of the corner locus of \(F\).  The
function \(\max\{-1,-u_1,-u_2\}\) has exactly one vertex, at
\(u=(1,1)\), where all three affine pieces meet; the two rays
\(\{u_1=u_2\le1\}\), \(\{u_1=1\le u_2\}\), \(\{u_2=1\le u_1\}\) emanating
from it are edges, along which the subgradient image is one-dimensional and
contributes nothing to the top product.  The mass at the vertex is the
lattice volume of the subgradient image, which is the convex hull of the
three gradients \((0,0)\), \((-1,0)\), \((0,-1)\); by the normalization
\eqref{eq:appJ-alexandrov} recorded in \Cref{app:normalizations} and used in
\Cref{ex:ch16-fan-p2}, that mass is
\[
 2!\cdot\vol\bigl(\operatorname{conv}\{(0,0),(-1,0),(0,-1)\}\bigr)
 =2\cdot\tfrac12=1 .
\]
So
\begin{equation}
\label{eq:ch17-toric-cld}
 c_1(\overline L)^2=\delta_{u=(1,1)},
\end{equation}
a single atom of mass \(1\), at the monomial point with
\(v(x)=v(y)=1\).

\emph{The Chambert--Loir side.}  \Cref{ex:metrics-cl-toric-surface} computes
the same measure from a model.  The metric is induced by the ideal
\(\mathfrak a=(\pi,x,y)\); let \(\rho:\mathcal X\to\mathbf P^2_{K^\circ}\)
be the blowup of the corresponding closed point \(p\) of the special fibre.
Then \(\mathcal X_s=\widetilde{X_s}+E\) with both multiplicities one, where
\(E\simeq\mathbf P^2_{\widetilde K}\) and \(\widetilde{X_s}\) is
\(\mathbf P^2_{\widetilde K}\) blown up at \(p\).  With
\(\mathcal L=\rho^*\cO(1)\otimes\cO_{\mathcal X}(-E)\),
\[
 \bigl(\mathcal L|_E\bigr)^2=1,
 \qquad
 \bigl(\mathcal L|_{\widetilde{X_s}}\bigr)^2=(h-e)^2=1-1=0,
\]
so \eqref{eq:cld-chambert-loir-formula} gives
\[
 c_1\bigl(\cO(1),\|\cdot\|_{\mathcal L}\bigr)^2
 =1\cdot1\cdot\delta_{x_E}+1\cdot0\cdot\delta_{x_{\widetilde{X_s}}}
 =\delta_{x_E},
\]
which is \eqref{eq:metrics-toric-measure}.

\emph{Verification.}  The divisorial point \(x_E\) is the monomial
valuation with \(v(\pi)=v(x)=v(y)=1\)---the exceptional divisor of the
blowup of \((\pi,x,y)\) is the exceptional divisor of the monomial ideal, so
\(\ord_E(\pi)=\ord_E(x)=\ord_E(y)=1\) and \(b_E=1\).  In the tropical
coordinates \(u\) that is the point \(u=(1,1)\).  Hence
\eqref{eq:ch17-toric-cld} and \eqref{eq:metrics-toric-measure} are the same
measure: one atom, at the same point, of the same mass \(1\).  Finally
\[
 \int_{X^{\An}}c_1(\overline L)^2=1=(L^2)=\bigl(\cO_{\mathbf P^2}(1)^2\bigr),
\]
which is the total-mass assertion of \Cref{cor:ch17-cl-mass}, of
\Cref{thm:cl-model-properties}, and of \Cref{thm:ch17-degree}, all three in
agreement.

\emph{What was independent about the two computations.}  The CLD side used a
convex piecewise affine function on \(\mathbf R^2\) and the Alexandrov
normalization \eqref{eq:appJ-alexandrov}; it never mentioned a model, a
special fibre, or an intersection number.  The Chambert--Loir side used
\(\mathcal O(-E)|_E\simeq\cO_{\mathbf P^2}(1)\), the self-intersection
\((h-e)^2=0\) on a Hirzebruch surface, and the multiplicities
\(b_{\widetilde{X_s}}=b_E=1\); it never mentioned a superform.  That the two
produce the same atom, at the same point, with the same mass, is the content
of \Cref{thm:cld-recovers-cl}, and it is also a check on the two
normalization conventions \eqref{eq:appJ-cl-measure} and
\eqref{eq:appJ-alexandrov} simultaneously: a discrepancy of a factor
\(n!=2\), or of the multiplicity \(b_E\), would have shown up here.

\emph{Relation to the cross-check already performed in
\Cref{chap:metrics}.}  \Cref{ex:metrics-cl-toric-surface} also compares its
intersection-theoretic answer with a convex computation, namely with
\(\MA_{\mathbf R}(F)=\tfrac12\delta_{(1,1)}\) and the factor \(n!=2\) of
\Cref{lem:toric-real-total-mass}, using the toric comparison
\Cref{thm:toric-ma-comparison}.  What is added here is the identification of
the object on the left: in \Cref{chap:metrics} the symbol
\(c_1(\cO(1),\|\cdot\|_{\mathcal L})^2\) denoted the model measure by
definition, and \Cref{thm:toric-ma-comparison} compared \emph{that} with a
real Monge--Amp\`ere measure.  Here the same symbol denotes an honest
Bedford--Taylor product of curvature currents on
\(\mathbf P^{2,\An}_K\), and \Cref{thm:cld-recovers-cl} is what makes the two
readings agree.  The chain of three computations---Bedford--Taylor product,
real Monge--Amp\`ere on \(N_{\mathbf R}\), intersection numbers on
\(\mathcal X_s\)---is therefore closed, with the normalization factor
\(n!\) appearing exactly once, in the middle link, as
\Cref{ex:ch16-toric-normalization} records.
\end{example}

\begin{remark}[A third check: the Fubini--Study description]
\label{rem:ch17-fs-check}
The metric of \Cref{ex:ch17-check-toric-surface} is the Fubini--Study metric
of \(\pi Z_0,Z_1,Z_2\), hence a semipositive model metric by
\Cref{prop:metrics-fs-model}; \Cref{thm:ch17-psh-residual-nef} then predicts
that its curvature current is positive and that its residual line bundles
are nef, both of which are visible in the computation:
\(F\) is convex, and the two special-fibre degrees \(1\) and \(0\) are
nonnegative.  The vanishing degree on \(\widetilde{X_s}\) is the
two-dimensional analogue of the vanishing degree on the contracted component
in \Cref{ex:ch17-check-chain}, and of the vanishing degree on
\(\widetilde C\) in \Cref{ex:ch17-check-p1}.  In all three examples the
entire mass sits on the exceptional component, which is the component on
which the model line bundle is relatively ample.
\end{remark}

\section{Bochner--Martinelli and Levine identities}
\label{sec:ch17-bm-levine}

Poincar\'e--Lelong is the codimension-one member of a family of identities.
We record two higher-codimensional forms because they provide useful tests
for signs and extension across polar sets.

The two identities have the same shape.  Each exhibits the integration
current \(\delta_Y\) of a subvariety \(Y\) of codimension \(d\) as the
\(\dprimeddprime\) of an explicit current, plus a correction which is a
curvature product.  In codimension one the correction is the curvature of
the metrized line bundle \(\cO(Y)\) and the identity is
\Cref{thm:ch17-pl-metrized}; in higher codimension the correction is
\(\alpha^p\) for the Levine formula and vanishes for
Bochner--Martinelli.  The proofs of both go through
\Cref{cor:ch17-truncated-max-fibre}, which is why they belong to this
chapter rather than to \Cref{chap:cld-poincare-lelong}: the identification of
a Monge--Amp\`ere current with a fibre is a statement about parlinear
metrics.

\subsection{Setting up the extension}
\label{sec:ch17-bm-setup}

Let \(X\) be boundaryless of pure dimension \(n\), let
\(f_1,\dots,f_d\) be a regular sequence of analytic functions on \(X\), and
let \(Y=V(f_1,\dots,f_d)\), of pure codimension \(d\).  Put
\(U=X\setminus Y\) and \(u=\log\max\{|f_1|,\dots,|f_d|\}\) on \(U\), a
continuous psh function there.  For every \(r\ge0\) the Bedford--Taylor
theory of \Cref{chap:cld-bedford-taylor} defines
\((\dprimeddprime u)^r\) and \(u(\dprimeddprime u)^r\) on \(U\), of bidegree
\((r,r)\).

\begin{lemma}
\label{lem:ch17-bm-extension}
For \(r<d\), the currents \((\dprimeddprime u)^r\) and
\(u(\dprimeddprime u)^r\) extend uniquely by zero to currents on \(X\); and
if \(r+1<d\) the extensions satisfy
\(\dprimeddprime\bigl(u(\dprimeddprime u)^r\bigr)=(\dprimeddprime u)^{r+1}\).
\end{lemma}

\begin{proof}
Every point of \(Y\) has \(d_K(x)\le n-d\), because \(Y\) has pure
codimension \(d\).  Thus \(Y\) is a closed subset all of whose points have
Abhyankar dimension \(\le n-d<n-r\), and the extension-by-zero criterion of
\Cref{lem:ch14-extension-by-zero} applies to a current of bidegree
\((r,r)\); uniqueness is \Cref{cor:ch14-rigid-invisible}.  The displayed
identity holds on \(U\) by \eqref{eq:ch16-dpsh-outer}, and both sides extend
by zero when \(r+1<d\), so it holds on \(X\).
\end{proof}

\begin{theorem}
\label{thm:cld-bochner-martinelli}
Let \(X\) be a pure \(n\)-dimensional boundaryless analytic space, and let
\(f_1,\ldots,f_d\) be a regular sequence of analytic functions.  Put
\(Y=V(f_1,\ldots,f_d)\) and, on \(X\setminus Y\),
\[
                  u=\log\max\{|f_1|,\ldots,|f_d|\}.
\]
The currents \(u(\dprimeddprime u)^{d-1}\) have a unique extension across \(Y\) in
the sense of \cite[\S18.7]{CLD25}, and
\[
             \dprimeddprime\bigl(u(\dprimeddprime u)^{d-1}\bigr)=\delta_Y.
\]
\end{theorem}
\begin{proof}
This is exactly \Cref{thm:cld-local-bochner-martinelli}, applied to the
regular sequence $f_1,\ldots,f_d$.  Its proof uses the truncations
\(
 u_r=\log\max\{r,|f_1|,\ldots,|f_d|\}
\), whose $d$-fold calibrated Monge--Amp\`ere currents are the fibres over
the Gauss point of radius $r$.  As $r\downarrow0$, calibrated conservation
of degree identifies the limiting fibre with the fundamental cycle of $Y$,
including its complete-intersection multiplicities.  Integration by parts
then gives the displayed identity.  Regularity supplies pure codimension
$d$ and the dimension bound required for the canonical extension; without
it, the cited theorem does not apply.
\end{proof}

\begin{proof}[Second proof, via the truncated maximum]
For \(r>0\), set
\[
 u_r=\log\max\{r,|f_1|,\dots,|f_d|\}
     =\log\max\{r,\rme^{u}\}.
\]
It is a parlinear function on all of \(X\), and by
\Cref{cor:ch17-truncated-max-fibre} with \(t=r\) and all \(r_i=1\),
\begin{equation}
\label{eq:ch17-bm-truncated}
 (\dprimeddprime u_r)^d=\delta_{f^{-1}(\eta_{r})},
\end{equation}
where \(\eta_r\) is the Gauss point of the polydisc of polyradius
\((r,\dots,r)\).

Let \(\omega\) be a smooth \((n-d,n-d)\)-form with compact support on \(X\).
By definition of the differential of a current and
\Cref{lem:ch17-bm-extension},
\[
 \bigl\langle\dprimeddprime\bigl(u(\dprimeddprime u)^{d-1}\bigr),\omega\bigr\rangle
 =\bigl\langle u(\dprimeddprime u)^{d-1},\dprimeddprime\omega\bigr\rangle
 =\bigl\langle u(\dprimeddprime u)^{d-1},\dprimeddprime\omega|_U\bigr\rangle,
\]
the form \(\dprimeddprime\omega|_U\) having compact support in \(U\) because
the extension is by zero.  On that compact support the continuous function
\(\max_i|f_i|=\rme^{u}\) does not vanish, hence is bounded below by some
\(r_0>0\); so \(u_r=u\) there for every \(0<r\le r_0\).  For such \(r\),
\[
 \begin{aligned}
  &\bigl\langle\dprimeddprime
    \bigl(u(\dprimeddprime u)^{d-1}\bigr),\omega\bigr\rangle\\
  &\qquad=\bigl\langle u_r(\dprimeddprime u_r)^{d-1},
    \dprimeddprime\omega|_U\bigr\rangle\\
  &\qquad=\bigl\langle\dprimeddprime
    \bigl(u_r(\dprimeddprime u_r)^{d-1}\bigr),\omega\bigr\rangle,
 \end{aligned}
\]
and the outer identity \eqref{eq:ch16-dpsh-outer} for the continuous psh
function \(u_r\) turns the right-hand side into
\(\langle(\dprimeddprime u_r)^d,\omega\rangle\).  By
\eqref{eq:ch17-bm-truncated} this is \(\langle\delta_{f^{-1}(\eta_r)},\omega\rangle\).

It remains to let \(r\downarrow0\) in the fibre.  That is
\Cref{prop:ch15-bm-fibre-stabilizes}: for \(t\) close enough to the origin
in \(\mathbf A^{d,\An}\), one has
\(\langle\delta_{f^{-1}(t)},\omega\rangle=\langle\delta_Y,\omega\rangle\),
because the calibrated tropicalizations of the fibres are eventually
constant.  Hence the left-hand side equals \(\langle\delta_Y,\omega\rangle\)
for all small \(r\), and being independent of \(r\) it equals it outright.
\end{proof}

\begin{remark}[What the second proof adds]
\label{rem:ch17-bm-second-proof}
The proof just given is not a different argument from
\Cref{thm:cld-local-bochner-martinelli}; it is the same argument, written so
that its one nonformal ingredient is isolated.  That ingredient is
\eqref{eq:ch17-bm-truncated}, the identification of
\((\dprimeddprime u_r)^d\) with a fibre, and it is a statement about
\emph{parlinear} metrics: \(u_r\) is the weight of the integral parlinear
metric of \Cref{cor:ch17-truncated-max-fibre} on a trivial bundle over
\(\mathbf A^{d,\An}\).  Everything else---the extension by zero, the
identity \eqref{eq:ch16-dpsh-outer}, the stabilization of the fibre---is
formal or was proved earlier.  This is the sense in which
Bochner--Martinelli is a corollary of the theory of metrized bundles rather
than of the theory of psh functions.
\end{remark}

\begin{example}[The two extreme codimensions]
\label{ex:ch17-bm-extremes}
For \(d=1\) the identity of \Cref{thm:cld-bochner-martinelli} reads
\(\dprimeddprime[\log|f_1|]=\delta_{V(f_1)}\), which is
\Cref{thm:cld-poincare-lelong}; the truncated proof specializes to the
truncated Poincar\'e--Lelong computation of
\Cref{ex:ch14-line-model-function}, and no extension lemma is needed.

For \(d=n\), so that \(Y\) is a finite set of rigid points, the identity says
that the \(n\)-fold Monge--Amp\`ere current of
\(\log\max\{|f_1|,\dots,|f_n|\}\)---understood through the extension of
\Cref{lem:ch17-bm-extension}---is the cycle \(\sum_Pm_P\delta_P\), where
\(m_P\) is the intersection multiplicity of the complete intersection at
\(P\).  Taking \(X=\mathbf A^{n,\An}_K\) and \(f_i=T_i\), the fibre over
\(\eta_r\) is a single point, of mass \(1\) by
\Cref{cor:ch17-truncated-max-fibre}, and the limit is \(\delta_0\); so
\[
 \dprimeddprime\Bigl(u\,(\dprimeddprime u)^{n-1}\Bigr)=\delta_0,
 \qquad u=\log\max\{|T_1|,\dots,|T_n|\}.
\]
This is the computation that \Cref{ex:ch16-polar-failure} contrasts with:
the naive product \((\dprimeddprime u)^n\) formed on
\(\mathbf A^{n,\An}\setminus\{0\}\) and extended by zero is \emph{zero}, and
the mass \(1\) is exactly the mass that the extension by zero discards.  The
two statements are not in conflict; they compute different objects, and the
Bochner--Martinelli identity is the one that sees the cycle.
\end{example}

\subsection{The Levine formula}
\label{sec:ch17-levine}

For the linear subspace
\(Y=\{x_0=\cdots=x_{p-1}=0\}\subset\mathbf P^n\), let \(\alpha\) be the
curvature current of the standard formal metric on \(\mathcal O(1)\), and
let \(\beta\) be the pullback of the corresponding current on
\(\mathbf P^{p-1}\) over \(\mathbf P^n\setminus Y\).  With
\[
 h=\log\frac{\max_i|x_i|}{\max_{i<p}|x_i|},\qquad
 \omega=\sum_{j=0}^{p-1}\alpha^{p-1-j}\beta^j,
\]
the current \(h\omega\) has a canonical extension \(\Lambda\) and the Levine
identity is
\[
                         \dprimeddprime\Lambda+\delta_Y=\alpha^p.
\]
The construction and extension theorem are proved in
\cite[\S18.7]{CLD25}; the displayed formula should not be used for an
arbitrary singular metric without the stated extension mechanism.

We indicate the mechanism, both because it is a model application of the
present chapter's tools and because it explains why the sum \(\omega\)
appears.  The full argument is \Cref{thm:cld-levine}.

\begin{proof}[Sketch of the Levine identity]
Let \(g:\widetilde X\to\mathbf P^n\) be the blowup along \(Y\), with
exceptional divisor \(E\), and let \(f':\widetilde X\to\mathbf P^{p-1}\)
extend the rational projection.  Put \(L_a=g^*\cO(1)\) and
\(L_b=(f')^*\cO(1)\) with their formal metrics; both are parlinear and psh
by \Cref{thm:ch17-weil-globally-psh}, and their curvature currents
\(\widetilde\alpha,\widetilde\beta\) restrict over
\(U'=\widetilde X\setminus E\) to \(\alpha,\beta\).  The canonical
isomorphism of \(L_a\) and \(L_b\) over \(U'\) extends to a morphism
\(\rho:L_a\to L_b\) with zero locus exactly \(E\), and
\(\log\|\rho\|^{-1}=g^*h\).  \Cref{cor:ch15-pl-section} gives
\begin{equation}
\label{eq:ch17-levine-blowup}
 \dprimeddprime[g^*h]+\delta_E=\widetilde\alpha-\widetilde\beta .
\end{equation}

Let \(\widetilde\omega=\sum_{j=0}^{p-1}\widetilde\alpha^{p-1-j}
\widetilde\beta^{j}\), a parlinear current of bidimension
\((n-p+1,n-p+1)\); its support is contained in a compact parlinear subset of
\(\widetilde X\) by
\Cref{prop:ch17-product-welldefined}\,\cref{itm:ch17-prod:skel}, hence meets
no strict Zariski-closed subset, hence lies in \(U'\).  So
\(\Lambda'=g^*h\cdot\widetilde\omega\) is defined, and
\(g_*\Lambda'=\Lambda\).  Multiplying
\eqref{eq:ch17-levine-blowup} by \(\widetilde\omega\), which is legitimate
because both \(\widetilde\alpha\) and \(\widetilde\beta\) are curvature
currents of parlinear psh metrics, gives
\[
 \dprimeddprime\Lambda'
 =(\widetilde\alpha-\widetilde\beta)\wedge\widetilde\omega
  -\widetilde\omega\wedge\delta_E
 =\widetilde\alpha^{\,p}-\widetilde\beta^{\,p}
  -\widetilde\omega\wedge\delta_E
 =\widetilde\alpha^{\,p}-\widetilde\omega\wedge\delta_E,
\]
where the middle equality is the telescoping identity
\((\widetilde\alpha-\widetilde\beta)\wedge\widetilde\omega
 =\widetilde\alpha^{\,p}-\widetilde\beta^{\,p}\)
and the last one uses \(\widetilde\beta^{\,p}=0\), the current
\(\widetilde\beta\) being pulled back from the \((p-1)\)-dimensional
\(\mathbf P^{p-1}\).
Pushing forward by \(g\), which is birational, and using
\Cref{prop:ch17-product-pushforward}\,\cref{itm:ch17-pf:proper}, it remains
to identify \(g_*(\widetilde\omega\wedge\delta_E)\) with \(\delta_Y\).
Restricting to \(E\), a projective bundle of relative dimension \(p-1\) over
\(Y\) on which \(L_a\) is trivial and \(L_b\) is the relative \(\cO(1)\),
the fibre integral of \(\widetilde\omega|_E\) over a fibre \(E_x\) is
\[
 \sum_{j=0}^{p-1}\bigl(c_1(L_a)^{p-1-j}c_1(L_b)^{j}\cap[E_x]\bigr)
 =\bigl(c_1(\cO(1))^{p-1}\cap[\mathbf P^{p-1}]\bigr)=1
\]
for every \(x\in Y\), only the term \(j=p-1\) surviving.  Fibre integration
(\Cref{thm:cld-fibre-projection}) then gives
\(g_*(\widetilde\omega\wedge\delta_E)=\delta_Y\), and the identity follows.
\end{proof}

\begin{example}[Levine on \(\mathbf P^n\), the two extreme cases]
\label{ex:ch17-levine-extremes}
For \(p=1\) the subspace \(Y=\{x_0=0\}\) is a hyperplane,
\(\omega=\alpha^0=1\), \(h=\log(\max_i|x_i|/|x_0|)\), and the identity
becomes \(\dprimeddprime[h]+\delta_Y=\alpha\), which is
\Cref{thm:ch17-pl-metrized} for the section \(x_0\) of
\(\overline{\cO(1)}\).  This recovers \Cref{ex:ch15-levine-p1}.

For \(p=n\) the subspace \(Y\) is the point \([0:\cdots:0:1]\), and
\(\alpha^n=\delta_{\eta_1}\) by
\Cref{ex:ch17-projective-good-reduction}.  The identity then reads
\(\dprimeddprime\Lambda=\delta_{\eta_1}-\delta_{Y}\): the mass \(1\) sitting
at the Gauss point differs from the mass \(1\) at the rigid point \(Y\) by
an exact current.  Both sides have total mass \(0\), as they must on a
compact boundaryless space; this is the consistency check performed in
\Cref{ex:ch15-levine-pn}.
\end{example}

\section{Functoriality and what positivity detects}
\label{sec:ch17-functoriality}

For a proper morphism \(f:Y\to X\), smooth forms satisfy the projection
formula
\[
       f_*\bigl(f^*\alpha\wedge T\bigr)=\alpha\wedge f_*T
\]
whenever all terms are defined.  The same statement holds for the formal PL
products above.  In particular, for a generically finite morphism of degree
\(e\) between proper \(n\)-folds,
\[
 f_*\!\left(\bigwedge_{i=1}^n c_1(f^*\overline L_i)\right)
      =e\bigwedge_{i=1}^n c_1(\overline L_i).
\]
On the model locus this follows equally from the algebraic projection
formula on special fibres.

Let us record the precise statements, since the projection formula is used
in \Cref{chap:compare-ma} and in \Cref{part:bfj}.

\begin{proposition}
\label{prop:ch17-projection}
Let \(f:Y\to X\) be a morphism of proper \(K\)-analytic spaces of pure
dimension \(n\) and let \(\overline L_1,\dots,\overline L_n\) be dpsh
metrized line bundles on \(X\).
\begin{enumerate}
\item\label{itm:ch17-proj:gf}
 If \(f\) is generically finite of calibrated degree \(e\), then
 \(f_*\bigl(\bigwedge_ic_1(f^*\overline L_i)\bigr)
   =e\bigwedge_ic_1(\overline L_i)\), and in particular
 \(\int_Y\bigwedge_ic_1(f^*\overline L_i)=e\int_X\bigwedge_ic_1(\overline L_i)\).
\item\label{itm:ch17-proj:bir}
 If \(f\) is birational then \(e=1\), and the total masses agree.
\item\label{itm:ch17-proj:model}
 If \(K\) is discretely valued, \(X\) and \(Y\) are proper varieties, and
 the metrics are model metrics, then \cref{itm:ch17-proj:gf} coincides with
 the algebraic projection formula \Cref{thm:model-projection-formula} applied
 to the special fibres of a common model.
\end{enumerate}
\end{proposition}

\begin{proof}
\cref{itm:ch17-proj:gf} is
\Cref{prop:ch17-product-pushforward}\,\cref{itm:ch17-pf:finite} in the finite
case; in general, factor \(f\) through the image of its finite locus and use
that the calibrated degree of the pushforward of a current is constant on
\(X^{\mathrm{gen}}\) by \Cref{prop:current-pushforward-formal}, the
exceptional locus being of smaller dimension and hence negligible for a top
product (\Cref{thm:ch14-support-dimension}).  Integrating and using
properness gives the mass statement.  \cref{itm:ch17-proj:bir} is the special
case \(e=1\).  For \cref{itm:ch17-proj:model}, both sides may be computed by
\Cref{thm:cld-recovers-cl} on a common model dominating \(X\) and \(Y\), and
the resulting identity of atomic measures is the special-fibre projection
formula, atom by atom; the divisorial points match because a component of
the special fibre of the model of \(Y\) dominating a component \(Y'\) of the
special fibre of the model of \(X\) defines the same normalized valuation
(\Cref{prop:component-valuation}).
\end{proof}

\section{Detection of positivity on a subvariety}
\label{sec:ch17-positivity-detection}

For an integral PL metric, nefness of its residual line bundle at \(x\)
implies positivity of the curvature current near \(x\).  Conversely, CLD's
subvariety-detection theorem shows, under its Abhyankar and PL hypotheses,
that a nonzero positive residual intersection can be detected after a
complete valued extension on a suitable analytic subvariety; see
\cite[\S18.9]{CLD25}.  This is a strict-positivity detection statement.  It
is not a blanket theorem that every positive current comes from a nef model.

\subsection{Positivity and vanishing are residual numerical conditions}
\label{sec:ch17-formal-positivity}

\begin{theorem}
\label{thm:ch17-formal-positivity}
Let \(X\) be boundaryless and topologically separated, let
\(\overline L_1,\dots,\overline L_d\) be line bundles on \(X\) with integral
parlinear metrics, and let \(x\in X\).  The current
\(c_1(\overline L_1)\wedge\cdots\wedge c_1(\overline L_d)\) is positive
(resp.\ zero) on a neighbourhood of \(x\) if and only if the product of
graded Chern classes
\(c_1(\widetilde{L_{1,x}})\cdots c_1(\widetilde{L_{d,x}})\) on
\(\widetilde{(X,x)}\) is numerically effective (resp.\ numerically trivial).
\end{theorem}

\begin{proof}[Proof, in outline]
A current, like a numerical cycle class, vanishes exactly when it is both
positive and negative, so it suffices to treat positivity.  Let \(m\) be the
analytic dimension of definition at \(x\) of the family
\((\overline L_i)\), after replacing the \(L_i\) by a common power so that
\(m\) is attained.  Choose a flat morphism of germs \(f:(X,x)\to(Y,y)\) with
\(\dim Y=m\) and metrized bundles \(\overline M_i\) on \((Y,y)\) with
\(\overline L_i\simeq f^*\overline M_i\); write \(T\) for the current in
question.

If \(m<d\), then \(T=0\) near \(x\) by the tropical-dimension argument in
the proof of \Cref{thm:ch17-atomic-support}, and the graded class vanishes
too, because a premodel of dimension \(<d\) carries no nonzero product of
\(d\) Chern classes.

If \(m=d\), \Cref{thm:cld-relative-ma} gives \(T=\lambda\delta_{f^{-1}(y)}\)
near \(x\), where \(\lambda\) is a strictly positive multiple of the graded
intersection number
\(\bigl(c_1(\widetilde{M_{1,y}})\cdots c_1(\widetilde{M_{d,y}})
 \cap[\widetilde{(Y,y)}]\bigr)\).
Since \(\delta_{f^{-1}(y)}\) is positive and nonzero near \(x\), positivity
of \(T\) is equivalent to \(\lambda\ge0\), which is equivalent to nefness of
the graded product; and any premodel of the \(\widetilde{M_{i,y}}\) yields a
premodel of the \(\widetilde{L_{i,x}}\), giving the asserted equivalence.

If \(m>d\), one localizes: choose finitely many
\(\varphi_i\in\widetilde{\cH(x)}^\times\) so that on each of the
\(2^{|I|}\) pieces \(U_\Phi^\varepsilon\) of the induced affine covering of
\(\widetilde{(X,x)}\) each \(\widetilde{L_{j,x}}\) has an invertible
homogeneous section, lift the \(\varphi_i\) and the sections to analytic
data, and take an affinoid \(V\ni x\) on which they are defined.  The
reductions \(\widetilde{V_\varepsilon}\) of the resulting affinoid domains
glue to a graded scheme \(\mathscr S\) of finite type, the reduction maps
glue to an anticontinuous \(\pi:V\to\mathscr S\), and the sections glue to
graded line bundles \(\mathscr L_j\) on \(\mathscr S\) whose restriction to
\(\overline{\{\pi(y)\}}\) is a premodel of \(\widetilde{L_{j,y}}\) for every
\(y\).  Nefness of \(c_1(\widetilde{L_{1,x}})\cdots c_1(\widetilde{L_{d,x}})\)
is then equivalent to
\(c_1(\mathscr L_1|_W)\cdots c_1(\mathscr L_d|_W)\ge0\) for every
irreducible \(W\subseteq\overline{\{\pi(x)\}}\) of dimension \(d\), and the
case \(m\le d\) already proved, applied at points \(y\) with
\(\dim\overline{\{\pi(y)\}}\le d\), propagates positivity to a
neighbourhood; the exceptional locus is weakly analytic of dimension
\(<n-d\) by \Cref{prop:ch17-dimdef-discrete}, hence tropically
\(H^{n-d}\)-negligible and invisible to the current
(\Cref{thm:cld-polar-negligibility}).  The converse implication runs the
same comparison backwards.  This is \cite[Theorem~18.8.4]{CLD25}.
\end{proof}

\begin{corollary}
\label{cor:ch17-positivity-stability}
Let \(\overline L_1,\dots,\overline L_d\) be parlinear metrized line bundles
on a boundaryless \(X\).
\begin{enumerate}
\item\label{itm:ch17-posst:prod}
 If each \(c_1(\overline L_i)\) is positive, so is
 \(c_1(\overline L_1)\wedge\cdots\wedge c_1(\overline L_d)\).
\item\label{itm:ch17-posst:pull}
 If \(c_1(\overline L_1)\wedge\cdots\wedge c_1(\overline L_d)\) is positive
 and \(f:Y\to X\) is a \(K\)-morphism from a boundaryless space over a
 complete extension, then
 \(c_1(f^*\overline L_1)\wedge\cdots\wedge c_1(f^*\overline L_d)\) is
 positive.
\item\label{itm:ch17-posst:harm}
 If \(u\) is a parlinear pluriharmonic function on \(X\) and \(f:Y\to X\) is
 as above, then \(u\circ f\) is pluriharmonic.
\end{enumerate}
\end{corollary}

\begin{proof}
Replacing the \(L_i\) by suitable powers we may assume the metrics integral
parlinear.  For \cref{itm:ch17-posst:prod}, positivity of each
\(c_1(\overline L_i)\) makes each residual bundle nef by
\Cref{thm:ch17-psh-residual-nef}, hence their product nef by Kleiman's
theorem, hence the product current positive by
\Cref{thm:ch17-formal-positivity}.  For \cref{itm:ch17-posst:pull}, nefness
of the graded product at \(f(y)\) implies nefness of its pullback at \(y\)
by \Cref{prop:ch17-residual-functorial}\,\cref{itm:ch17-res:pull}, and one
concludes by \Cref{thm:ch17-formal-positivity} again.
\cref{itm:ch17-posst:harm} is \cref{itm:ch17-posst:pull} applied to the
trivial bundle with the metric \(\|1\|=\rme^{-u}\) and to its inverse:
pluriharmonicity of \(u\) means that both curvature currents vanish.
\end{proof}

\begin{example}[Support of \(c_1(\cO(1))^d\) on projective space]
\label{ex:ch17-support-projective}
Let \(\cO(1)\) on \(\mathbf P^{n,\An}_K\) carry the Weil metric and let
\(1\le d\le n\).  By \Cref{lem:ch17-residual-formal}, the residual bundle at
\(x\) has premodel \((V_x,\cO(1)|_{V_x})\) with
\(V_x=\overline{\{\operatorname{red}(x)\}}\subseteq\mathbf P^n_{\widetilde K}\),
and \(\cO(1)\) is ample on \(V_x\).  Therefore the numerical class
\(c_1(\cO(1)|_{V_x})^d\cap[V_x]\) is nonzero exactly when
\(\dim V_x\ge d\).  By \Cref{thm:ch17-formal-positivity},
\[
 \Supp\bigl(c_1(\overline{\cO(1)})^d\bigr)
 =\bigl\{x\in\mathbf P^{n,\An}_K:
   \dim\overline{\{\operatorname{red}(x)\}}\ge d\bigr\}.
\]
For \(d=n\) this is the single Gauss point of
\Cref{ex:ch17-projective-good-reduction}; for \(d=1\) it is the set of
points whose reduction is not a closed point, that is, the complement of the
union of the open sets \(\operatorname{red}^{-1}(\text{closed point})\)---by
anticontinuity of the reduction map
(\Cref{prop:berk-reduction-anticontinuous}) that complement is closed, as it
must be.  For \(n=1\), \(d=1\) it is \(\{\zeta_{0,1}\}\), recovering
\Cref{ex:ch17-dimdef-p1}.
\end{example}

\subsection{Detecting strict positivity on a curve}
\label{sec:ch17-detection}

\Cref{thm:ch17-formal-positivity} converts an analytic question into a
numerical one on a graded variety.  The remaining question is whether a
nonzero positive product can be \emph{seen} by restricting to a small
subvariety, as in the complex case, where pshness of a metric is detected on
curves.  The answer is yes, at Abhyankar points, and the two graded lemmas
below are the geometric input; we quote them from
\cite[\S18.9]{CLD25} and use them only here.

\refstepcounter{lemma}\label{lem:ch17-detect-fibration}
\begin{quote}
\textbf{Foundational input \thelemma\ (positive degree on a generic fibre;
source boundary).}
Let \(\mathcal K\) be a graded field, let \(q:V\to Z\) be a surjective
morphism of proper integral graded \(\mathcal K\)-schemes of generic
relative dimension \(d\), and let \(L_1,\dots,L_d\) be line bundles on
\(V\) with \(P=c_1(L_1)\cdots c_1(L_d)\) numerically effective.  Suppose
there are a point \(z\in Z\) with graded residue field finite over
\(\mathcal K\) and an irreducible component \(H\) of \(q^{-1}(z)\) with
\(\int_HP>0\).  Then the degree of \(P\) on the generic fibre of \(q\) is
strictly positive.
This statement is imported in full from \cite[Lemma~18.9.1]{CLD25}; no
proof is claimed here.
\end{quote}

\refstepcounter{lemma}\label{lem:ch17-detect-model}
\begin{quote}
\textbf{Foundational input \thelemma\ (fibring a nef nontrivial class;
source boundary).}
Let \(\mathcal K\) be a graded field, let \(V\) be a proper integral graded
\(\mathcal K\)-scheme of dimension \(n\), and let \(L_1,\dots,L_d\) be line
bundles with \(P=c_1(L_1)\cdots c_1(L_d)\) nef and numerically nontrivial.
Then there are a proper integral graded \(\mathcal K\)-scheme \(W\), a
birational \(p:W\to V\), a polyradius \(r\) of length \(n-d\), and a
dominant \(q:W\to\mathbf P^{1,r}_{\mathcal K}\) such that the degree of
\(c_1(p^*L_1)\cdots c_1(p^*L_d)\) on the generic fibre of \(q\) is strictly
positive.
This statement is imported in full from \cite[Proposition~18.9.2]{CLD25};
no proof is claimed here.
\end{quote}

\begin{theorem}
\label{thm:ch17-detect-subvariety}
\index{positivity!detected on a subvariety}
Let \((X,x)\) be a boundaryless germ of \(K\)-analytic space, and let
\(\overline L_1,\dots,\overline L_d\) be line bundles on \(X\) with
parlinear metrics.  Assume that
\(c_1(\overline L_1)\wedge\cdots\wedge c_1(\overline L_d)\) is positive and
nonzero, and that \((X,x)\) admits a Zariski-closed subset of dimension
\(d_K(x)\).  Then there is an immersion \(\varphi:(Y,x)\hookrightarrow(X,x)\)
of a boundaryless germ of dimension \(d\) such that
\(c_1(\varphi^*\overline L_1)\wedge\cdots\wedge c_1(\varphi^*\overline L_d)\)
is nonzero.
\end{theorem}

\begin{proof}
Set \(n=d_K(x)\) and \(m=n-d\).  Replacing \((X,x)\) by the postulated
Zariski-closed subset of dimension \(n\), we may assume
\(\dim(X,x)=n\); and replacing the \(L_i\) by powers, that the metrics are
integral parlinear.

By \Cref{thm:ch17-formal-positivity} the graded class
\(P=c_1(\widetilde{L_{1,x}})\cdots c_1(\widetilde{L_{d,x}})\) on
\(\widetilde{(X,x)}\) is nef and numerically nontrivial.  Choose a model
\[
 (V,M_1,\dots,M_d)
 \quad\text{of}\quad
 \bigl(\widetilde{(X,x)},\widetilde{L_{1,x}},\dots,
       \widetilde{L_{d,x}}\bigr)
\]
such that the class
\[
 c_1(M_1)\cdots c_1(M_d)
\]
is nef and numerically nontrivial on \(V\).  By
\Cref{lem:ch17-detect-model}, after
modifying \(V\) we may assume there are a polyradius
\(r=(r_1,\dots,r_m)\) and a dominant
\(q:V\to\mathbf P^{1,r}_{\mathcal K}\) whose generic fibre \(W\) carries
\(\int_Wc_1(M_1|_W)\cdots c_1(M_d|_W)>0\); \Cref{lem:ch17-detect-fibration}
is what makes that possible.

Choose homogeneous affine coordinates \(T_1,\dots,T_m\) on the graded affine
space inside \(\mathbf P^{1,r}\), with \(\deg T_i=r_i\), and let
\(f_i\in\widetilde K(V)=\widetilde{\cH(x)}\) be their images under \(q\).
Since \((X,x)\) is boundaryless, hence almost good, there are invertible
analytic functions \(g_1,\dots,g_m\) on the germ \((X,x)\) with
\(\widetilde{g_i(x)}=f_i\).  Let \(g:(X,x)\to(\mathbf A^{m,\An}_K,z)\) be
the induced morphism of germs, and set \((Y,x)=(g^{-1}(\eta_r),x)\).

We claim \(\dim(Y,x)=d\) and \(d_{\cH(\eta_r)}(x)=d\).  Indeed, for every
\(y\in Y\),
\(n\ge d_K(y)=d_{\cH(\eta_r)}(y)+d_K(\eta_r)=d_{\cH(\eta_r)}(y)+m\), so
\(d_{\cH(\eta_r)}(y)\le d\); this gives \(\dim(Y,x)\le d\).  Applying the
same computation at \(x\), where \(d_K(x)=n\), gives equality there, hence
\(\dim(Y,x)=d\).

Finally, put \(\Lambda_i=L_i|_Y\).  By
\Cref{thm:ch17-formal-positivity} applied on \(Y\), it suffices to show that
\[
 \int_{P_{\widetilde{\cH(x)}/\widetilde{\cH(z)}}}
 c_1(\widetilde{\Lambda_{1,x}})\cdots
 c_1(\widetilde{\Lambda_{d,x}})\ne0.
\]
By functoriality of the residual bundle
(\Cref{prop:ch17-residual-functorial}\,\cref{itm:ch17-res:pull}) the tuple
\[
 (W,M_1|_W,\dots,M_d|_W)
\]
is a model of
\[
 \bigl(P_{\widetilde{\cH(x)}/\widetilde{\cH(z)}},
       \widetilde{\Lambda_{1,x}},\dots,
       \widetilde{\Lambda_{d,x}}\bigr).
\]
The integral above therefore equals
\[
 \int_Wc_1(M_1|_W)\cdots c_1(M_d|_W)>0.
\]
\end{proof}

\begin{corollary}
\label{cor:ch17-curve-test}
Let \((X,x)\) be a boundaryless germ and let \(L\) be a line bundle with a
parlinear psh metric such that \(c_1(\overline L)\ne0\).  If \((X,x)\)
admits a Zariski-closed subset of dimension \(d_K(x)\), then there is a germ
of curve \((Y,x)\) in \((X,x)\) with \(c_1(\overline L|_Y)\ne0\).
\end{corollary}

\begin{proof}
Apply \Cref{thm:ch17-detect-subvariety} with \(d=1\).
\end{proof}

\begin{example}[Detection on a subvariety, computed]
\label{ex:ch17-detection-computed}
Let \(K\) be discretely valued, \(X=\mathbf P^{2,\An}_K\), and let
\(\overline L=\overline{\cO(1)}\) carry the Weil metric.  Take \(x\) to be
the Gauss point \(\eta_1\), so \(d_K(x)=2\) and
\(\widetilde{\cH(x)}=\widetilde K(s,t)\) with \(s,t\) the reductions of the
torus coordinates \(x_1/x_0\), \(x_2/x_0\).

\emph{The curvature is nonzero.}  By
\Cref{ex:ch17-projective-good-reduction}, \(c_1(\overline L)^2
=\delta_{\eta_1}\ne0\); and \(c_1(\overline L)\ne0\) a fortiori.  The
residual bundle at \(\eta_1\) has premodel
\((\mathbf P^2_{\widetilde K},\cO(1))\), so its class is nef and
numerically nontrivial, as \Cref{thm:ch17-formal-positivity} predicts.

\emph{The subvariety produced by the theorem.}  Apply
\Cref{thm:ch17-detect-subvariety} with \(d=1\), so \(m=1\).  A fibration of
\(V=\mathbf P^2_{\widetilde K}\) with positive degree on the generic fibre
is the linear projection \(q:V\dashrightarrow\mathbf P^1_{\widetilde K}\)
from a point, resolved on the blowup \(W\to V\); its generic fibre is a
line, and \(\int_{\text{line}}c_1(\cO(1))=1>0\).  Concretely, take
\(q\) to be given by \(f=s\in\widetilde K(V)\), a coordinate of degree
\(1\).  Lift \(f\) to the analytic unit \(g=x_1/x_0\), take
\(r=1\), and set
\[
 Y=g^{-1}(\eta_{1})
 =\{y\in\mathbf P^{2,\An}_K:\ |x_1/x_0|(y)=1
    \text{ and }x_1/x_0\text{ is a Gauss point over }y\},
\]
which is the fibre of \(x_1/x_0:\mathbf P^{2,\An}_K\dashrightarrow
\mathbf A^{1,\An}_K\) over the Gauss point of the unit disc: a
one-dimensional \(\cH(\eta_1)\)-analytic space, namely a projective line
over the complete residue field \(\cH(\eta_1)\).

\emph{Verification.}  On that fibre, \(\cO(1)|_Y\) is
\(\cO_{\mathbf P^1}(1)\) over \(\cH(\eta_1)\) with its Weil metric, so by
\Cref{ex:ch17-weil-curvature}\,(a),
\[
 c_1(\overline L|_Y)=\delta_{\zeta},\qquad
 \int_Yc_1(\overline L|_Y)=1=\deg\bigl(\cO(1)|_Y\bigr)\ne0,
\]
where \(\zeta\) is the Gauss point of \(Y\).  So the restriction to a
one-dimensional subvariety detects the positivity, with the same numerical
value \(1\) that \Cref{lem:ch17-detect-fibration} produced on the graded
side.  The two \(1\)'s are the same number: the degree of \(\cO(1)\) on a
line in \(\mathbf P^2_{\widetilde K}\), and the degree of the restricted
metrized bundle on \(Y\).

\emph{What would go wrong without the Abhyankar hypothesis.}  Had we taken
for \(x\) a point with \(d_K(x)=0\)---a rigid point of
\(\mathbf P^{2,\An}_K\)---the residual bundle would be defined over a
zero-dimensional graded variety, the curvature current would vanish near
\(x\) by \Cref{thm:ch17-atomic-support}, and there would be nothing to
detect.  The hypothesis that \((X,x)\) contains a Zariski-closed subset of
dimension \(d_K(x)\) is what allows the germ to be replaced by an
\(n\)-dimensional one; whether it can be dropped is not known
\cite[Remark~18.9.6]{CLD25}.
\end{example}

\begin{warning}{Detection is not a converse to the model construction}
\Cref{thm:ch17-detect-subvariety} says that a nonzero positive parlinear
product is visible on some \(d\)-dimensional subvariety, possibly after a
complete valued extension of the base.  It does \emph{not} say that a
positive current arises from a nef model, nor that pshness may be
\emph{defined} by restriction to curves.  In the complex case those two
further statements are true and are how pshness is usually characterized;
here the first is the open question of the following warning, and the second
fails for the class of curve-test psh functions of \cite{GR25Classical},
which is compared with the CLD class in
\Cref{thm:cld-to-classical-psh,thm:gr-pl-comparison} and is not identified
with it outside the parlinear case.
\end{warning}

\begin{warning}{The remaining local-to-global gap}
The first question in \cite[Remark~17.8.10]{CLD25} asks for global
\emph{smooth psh} approximation of continuous locally psh metrics on
projective varieties.  It is not the same question as approximation on a
fixed line bundle by nef rational model metrics.  The latter implication
holds in positive degree on smooth proper geometrically connected curves
over every complete nontrivially valued field by
\cref{thm:p081-corrected-curve-comparison}, and fails without the
model-existence condition in degree zero by
\cref{thm:p081-irrational-tate-period-obstruction}.  The second paragraph of that
remark is different: it asks whether the published implication
``parlinear implies globally dpsh'' extends to a proper strict analytic
space admitting a formal model with projective special fibre.  That second
question is answered negatively by
\Cref{prop:p082-formal-twistor-counterexample-candidate}; the first remains
open in its generality and is recorded precisely in
\cref{prob:frontier-cld-global-metric}.  The positive-degree comparison
does not settle higher dimension or singular curves.  Its arbitrary-field
scope does not establish a higher-dimensional envelope theorem.
\end{warning}

\section{Chapter dictionary}
\label{sec:ch17-dictionary}

\subsection{From the objects of this chapter to their use}
\label{sec:ch17-dict-forward}

\begin{center}
\small
\begin{tabular}{p{.27\textwidth}|p{.35\textwidth}|p{.30\textwidth}}
\textbf{Object}&\textbf{Established property}&\textbf{First use}\\ \hline
metric weight \(\phi_e=-\log\|e\|\)
 &frame ambiguity is \(\log|g|\)
  (\eqref{eq:cld-weight-change-frame})
 &\Cref{prop:ch17-c1-construction}\\
orthogonal frame
 &unique norms, stable under perturbation
  (\Cref{lem:ch17-orthogonal-frames})
 &\Cref{sec:ch17-residual}\\
smooth metric
 &exists in rank one, never in rank \(\ge2\)
  (\Cref{prop:ch17-smooth-existence},
   \Cref{prop:ch17-no-smooth-higher-rank})
 &\Cref{thm:ch17-degree}\\
integral parlinear metric
 &formal, in the strict case
  (\Cref{thm:cld-formal-pl-comparison})
 &\Cref{sec:ch17-formal}\\
residual bundle \(\widetilde L_x\)
 &graded bundle on \(\widetilde{(X,x)}\); premodel from a formal model
  (\Cref{lem:ch17-residual-formal})
 &\Cref{thm:ch17-atomic-masses}\\
dimension of definition
 &high-dimension locus is discrete
  (\Cref{prop:ch17-dimdef-discrete})
 &\Cref{thm:ch17-atomic-support}\\
first Chern current \(c_1(\overline L)\)
 &closed; positive iff residual nef
  (\Cref{thm:ch17-psh-residual-nef})
 &\Cref{chap:delta-forms}, \Cref{chap:compare-ma}\\
Poincar\'e--Lelong for sections
 &\(\dprimeddprime[\log\|s\|^{-1}]+\delta_{\divis(s)}=c_1(\overline L)\)
  (\Cref{thm:ch17-pl-metrized})
 &\Cref{cor:ch17-curve-total-mass},
  \Cref{thm:ch17-degree}\\
mixed curvature product
 &symmetric, multilinear, functorial
  (\Cref{prop:ch17-product-welldefined})
 &\Cref{thm:cld-recovers-cl}\\
degree theorem
 &total mass \(=(L_1\cdots L_n)\)
  (\Cref{thm:ch17-degree})
 &\Cref{cor:ch17-cl-mass}, \Cref{chap:compare-ma}\\
atomic top product
 &\(\sum_x\lambda_x\delta_x\), \(\lambda_x\in\mathbf Z\)
  (\Cref{thm:ch17-atomic-support,thm:ch17-atomic-masses})
 &\Cref{thm:cld-formal-intersection}\\
relative Monge--Amp\`ere
 &atoms become fibres
  (\Cref{thm:cld-relative-ma})
 &\Cref{cor:ch17-truncated-max-fibre}\\
defect factor \(\epsilon_{\mathfrak X}\)
 &\(=b_Y\) in the discretely valued normal case
  (\Cref{prop:ch17-defect-dvr})
 &\Cref{thm:cld-recovers-cl}\\
recovery of Chambert--Loir
 &\eqref{eq:cld-chambert-loir-formula}
 &\Cref{chap:compare-ma},
  \Cref{chap:master-dictionary}\\
Bochner--Martinelli
 &\(\dprimeddprime(u(\dprimeddprime u)^{d-1})=\delta_Y\)
  (\Cref{thm:cld-bochner-martinelli})
 &\Cref{chap:delta-forms}\\
detection on a subvariety
 &nonzero positive product restricts nonzero
  (\Cref{thm:ch17-detect-subvariety})
 &\Cref{chap:cld-examples}\\
locally CLD-psh metric on a smooth proper geometrically connected algebraic curve
 &FS/model approximable in positive degree; nef rational-model
  approximable in degree zero iff one such metric exists
  (\cref{thm:p081-corrected-curve-comparison})
 &\Cref{prob:frontier-cld-global-metric}\\
smooth harmonic cubist metric on a degree-zero Tate bundle
 &need not lie in the fixed-bundle nef rational-model closure
  (\cref{thm:p081-irrational-tate-period-obstruction})
 &\cref{prob:frontier-cld-global-metric}\\
parlinear metric on a proper strict analytic space with projective-special-
fibre formal model
 &need not be globally dpsh
  (\Cref{prop:p082-formal-twistor-counterexample-candidate})
 &\Cref{prob:frontier-cld-global-metric}\\
\end{tabular}
\end{center}

\subsection{From the later theory back to this chapter}
\label{sec:ch17-dict-backward}

\begin{center}
\small
\begin{tabular}{p{.30\textwidth}|p{.32\textwidth}|p{.30\textwidth}}
\textbf{Later object}&\textbf{Smallest instance}&\textbf{Computed in}\\
\hline
Chambert--Loir measure of a model metric
 &\(\delta_{\zeta_{0,|\pi|}}\), mass \(1\)
 &\Cref{ex:ch17-check-p1}\\
mixed model measure on a curve
 &three atoms summing to \(2\)
 &\Cref{ex:ch17-check-chain}\\
Chambert--Loir measure on a surface
 &one atom at \(u=(1,1)\), mass \(1=(L^2)\)
 &\Cref{ex:ch17-check-toric-surface}\\
defect factor \(\ne1\)
 &\(b_{E_2}=2\), atoms \(2,1,-2\)
 &\Cref{ex:ch17-multiple-fibre}\\
curvature of a rank-two metric
 &one atom at \(|x_2/x_1|=a_1/a_2\)
 &\Cref{ex:ch17-rank-two}\\
support of \(c_1(\cO(1))^d\)
 &reduction of dimension \(\ge d\)
 &\Cref{ex:ch17-support-projective}\\
positivity detected on a curve
 &a line in \(\mathbf P^2\), degree \(1\)
 &\Cref{ex:ch17-detection-computed}\\
\(\mathbf R\)-parlinear, non-model metric
 &atoms \(1+\theta\) and \(-\theta\)
 &\Cref{ex:ch17-weil-curvature}\,(d)\\
\end{tabular}
\end{center}

\subsection{Hypothesis ledger}
\label{sec:ch17-ledger}

\begin{center}
\small
\begin{tabular}{p{.30\textwidth}|p{.30\textwidth}|p{.32\textwidth}}
\textbf{Statement}&\textbf{Hypotheses that cannot be dropped}
 &\textbf{Where it breaks}\\ \hline
\Cref{prop:ch17-c1-construction}
 &boundaryless; locally integrable weight
 &Green's formula acquires a flux term\\
\Cref{thm:ch17-degree}
 &\(X\) proper and algebraic; dpsh metrics
 &nonalgebraic case needs
  \Cref{rem:ch17-integrality-analytic}\\
\Cref{thm:ch17-atomic-support}
 &integral parlinear; boundaryless
 &an \(\mathbf R\)-parlinear metric still works after scaling, a
  merely continuous one does not\\
\Cref{thm:cld-relative-ma}
 &generic flatness; \(\dim X=m\)
 &fibre dimension jumps and \(u_\omega\) is not
  \(\langle\delta_{Y_x},\omega\rangle\)\\
\Cref{thm:cld-formal-intersection}
 &\(\mathfrak X\) flat, pure relative dimension \(n\); \(X\)
  boundaryless
 &components of \(\mathfrak X_s\) fail to be proper\\
\Cref{thm:cld-recovers-cl}
 &\(K\) discretely valued; \(\mathcal X\) relatively normal
 &several points over a generic point; no Chow degree\\
\Cref{thm:ch17-cl-continuous}
 &semipositivity
 &no psh limit, and mass may escape\\
\Cref{thm:cld-bochner-martinelli}
 &regular sequence
 &codimension and the extension-by-zero bound fail\\
\Cref{thm:ch17-detect-subvariety}
 &Abhyankar: a Zariski-closed subset of dimension \(d_K(x)\)
 &not known to be removable
  \cite[Remark~18.9.6]{CLD25}\\
\cref{thm:p081-corrected-curve-comparison}
 &\(K\) complete and nontrivially valued; \(C\) smooth, proper,
  geometrically connected, and algebraic
 &positive degree is essential to the unconditional approximation;
  degree zero requires one actual nef rational model metric\\
\cref{thm:p081-irrational-tate-period-obstruction}
 &\(K\) algebraically closed; a Tate period ratio
  \(a/\ell\notin\mathbf Q\); fixed line bundle and rational model metrics
 &does not refute smooth global pshness, ample-class envelopes, or a
  definition allowing an ample perturbation\\
\end{tabular}
\end{center}

\section{Source note and interface}
\label{sec:ch17-source-note}

\paragraph{The construction.}
The metric, residual-bundle, and psh material in this chapter is reorganized
from \cite[Chapters~17--18]{CLD25}.  The formal intersection formula and the
comparison with Chambert--Loir are cited at their exact locations,
\cite[Theorem~18.5.5 and Proposition~18.5.8]{CLD25}; the original measure is
\cite{CL06}.  The distinction between algebraic, formal, and PL metrics uses
\cite{GM19}.  \Cref{chap:delta-forms} explains how Gubler--K\"unnemann
enlarge the calculus so that the curvature of formal metrics can be wedged
as a geometric object rather than only as a current or an atomic measure.

\paragraph{Item by item.}
\Cref{sec:ch17-metrics} follows \cite[\S17.1]{CLD25}; the nonexistence of
smooth metrics in higher rank is \cite[Remark~17.1.9]{CLD25}, and the
orthogonalization it rests on is Temkin's \cite[\S2.3]{Tem16}.
\Cref{sec:ch17-parlinear} is \cite[\S\S17.2 and 17.7]{CLD25}; the
equivalence of formal and strictly integral parlinear metrics is
\cite[Proposition~17.7.5]{CLD25}, which for line bundles was proved by
Gubler and, without hypotheses, in \cite[Proposition~2.10]{GM19}---the
statement quoted in this book as \Cref{thm:gm-metric-comparison}.
\Cref{sec:ch17-residual,sec:ch17-dimdef} are
\cite[\S\S17.3--17.6]{CLD25}; the graded reduction they use is Temkin's
\cite{Tem16,Temkin2004}, summarized in \Cref{sec:appG-graded-reduction}.
\Cref{sec:ch17-chern} is \cite[\S18.1]{CLD25}; the total-mass statement on a
curve is \cite[Corollary~18.1.7]{CLD25}.  \Cref{sec:ch17-psh} is
\cite[\S17.8]{CLD25}, and the global psh statement for the Weil metric uses
Demailly's regularized maximum, quoted from \cite[Lemma~5.18]{Dem12} in the
form recalled in the proof of \Cref{thm:ch17-weil-globally-psh}.
\Cref{sec:ch17-products} is \cite[\S\S18.3--18.4]{CLD25}, with the degree
theorem \cite[Theorem~18.3.11]{CLD25}; \Cref{sec:ch17-relative-ma} is
\cite[\S18.6]{CLD25}; \Cref{sec:ch17-formal} is \cite[\S18.5]{CLD25};
\Cref{sec:ch17-bm-levine} is \cite[\S18.7]{CLD25}, the Bochner--Martinelli
identity being due to Mihatsch and reproved there;
\Cref{sec:ch17-positivity-detection} is \cite[\S\S18.8--18.9]{CLD25}.

The corrected curve comparison
\cref{thm:p081-corrected-curve-comparison} uses Thuillier's global
approximation and potential theory \cite{Thu05}, the comparison of
continuous CLD-psh and classical subharmonic functions \cite{Wanner19},
and the formal/algebraic metric and finite-field descent results
\cite{GM19,GJKM19}.  It applies over every complete nontrivially valued
ground field in the displayed smooth proper geometrically connected
algebraic curve scope.  Positive degree supplies a strict curvature margin
for rational approximation.  In degree zero, existence of one nef rational
model metric is an additional hypothesis, not a consequence of generic
nefness.

The two-edge calculation in
\cref{thm:p081-irrational-tate-period-obstruction} uses the Tate
uniformization \cite{Tate95}, semistable models \cite{BL85}, and the
slope formula \cite{BPR14}.  It disproves the unrestricted degree-zero
model-existence assertion used in an earlier version of this chapter.
The old proof tried to rationalize a globally defined real-PL metric while
preserving nefness and its fixed line-bundle cocycle.  That step is not
valid without a further argument; local real-PL presentations do not
supply it.  The earlier theorem and its dependent transfer proofs have
therefore been withdrawn from the live text.  Their source files and
verification records remain archived as superseded evidence.

Both corrected statements and their complete proofs passed independent
source-aware review and strict production Rethlas verification using
\texttt{gpt-5.6-sol} at \texttt{xhigh}.  The exact-source provenance is
recorded in \cref{sec:appL-curve-correction-20260831}.  These are
model-assisted checks, not formal proof-assistant certificates or claims
of novelty.  In particular, they neither settle the higher-dimensional
local-to-global question nor disprove continuity of ample-class envelopes.

The formal-twistor counterexample
\Cref{prop:p082-formal-twistor-counterexample-candidate} answers the second
question of \cite[Remark~17.8.10]{CLD25} negatively.  Its independent
twelve-obligation exact-byte audit passed, and the dependency-closed
production request using \texttt{gpt-5.6-sol} at \texttt{xhigh} returned
\texttt{correct}, with zero critical errors and zero gaps.  The frozen
statement/proof and archived request/response SHA-256 values are
\begin{center}
\small
\begin{tabular}{@{}c@{}}
\nolinkurl{1e9b7c31a70c55f407de07d76c67d6912e61e1673ef81c66c60789afd2202d2f}\\
\nolinkurl{78fab726abcdfb3086988150359568ccd8a9b29d63c3a94e0c513f14b9fc148c}\\
\nolinkurl{7916d3750482946cd933591a3f836d6e3b308dcbddb8a24a26c836de41d73e1b}\\
\nolinkurl{f5e780184515b39bbb7d4fd099366da98b96a74896fa1d693b43f66cce19d57e}.
\end{tabular}
\end{center}
The conclusion retires P08.2 only.  P08.1 remains open outside the positive
curve scopes stated above.

\paragraph{What we have organized differently.}
Four choices depart from the source's arrangement, and the reader comparing
the two should be aware of them.  First, we have separated the
\emph{construction} of the curvature current
(\Cref{prop:ch17-c1-construction}) from its \emph{computation}
(\Cref{sec:ch17-formal}), so that the Chambert--Loir comparison can be
stated as a single theorem with a single hypothesis list rather than as the
last of a chain.  Second, we have made the three defect factors into a named
definition, \Cref{def:ch17-defect}, and proved their vanishing in the two
cases the book uses (\Cref{prop:ch17-defect-reduced,prop:ch17-defect-dvr})
before the Chambert--Loir theorem rather than after it.  Third, we have
given \Cref{thm:cld-relative-ma} a second proof that uses only fibre
integration, because that is the form in which
\Cref{cor:ch17-truncated-max-fibre}---and hence Bochner--Martinelli---needs
it.  Fourth, all worked examples are carried to explicit numbers and checked
against \Cref{chap:metrics}: this is \Cref{sec:ch17-cl-checks}, which has no
counterpart in the source, and it is the only place in the book where the
normalizations \eqref{eq:appJ-cl-measure}, \eqref{eq:appJ-cld-c1} and
\eqref{eq:appJ-alexandrov} are tested simultaneously against one another.

\paragraph{The complex prototype, and where it fails.}
In complex geometry the curvature of a smooth Hermitian metric is a smooth
form, positivity of that form is equivalent to pshness of the local weights,
and pshness may be tested on curves.  All three statements have partial
analogues here, and each degrades in a different way.  Smooth metrics exist
only in rank one (\Cref{prop:ch17-no-smooth-higher-rank}), and the metrics
that matter---the formal ones---are not smooth even then.  Positivity of the
curvature current is equivalent to pshness only for parlinear metrics, where
both are equivalent to a numerical condition on a residual bundle
(\Cref{thm:ch17-psh-residual-nef}); for a general continuous metric the
equivalence is not available, and the transfinite hierarchy of
\Cref{def:ch15-hierarchy} is the reason.  The curve test survives as
\Cref{cor:ch17-curve-test}, but only as a detection statement at an
Abhyankar point and only for parlinear metrics; the curve-test psh functions
of \cite{GR25Classical,GR25PL} form a genuinely different class, compared
with the CLD one in \Cref{thm:cld-to-classical-psh} and
\Cref{thm:gr-pl-comparison}.

\paragraph{The connection this chapter establishes.}
\Cref{chap:metrics} defined the Chambert--Loir measure of a model
metric and left the notation \(c_1(\overline L_1)\wedge\cdots\wedge
c_1(\overline L_n)\) unjustified; \Cref{thm:cld-recovers-cl} justifies it.
This is the first arrow of the master diagram of
\Cref{chap:master-dictionary} to be established in full, and it is the
foundation of the comparison of Monge--Amp\`ere operators in
\Cref{chap:compare-ma}: on model metrics the BFJ operator, the
Chambert--Loir measure, and the CLD Bedford--Taylor product are one object.
Beyond the model locus the three theories separate, and the delimitation
table of \Cref{sec:ch16-delimitation-table} says where.

\paragraph{What is new here, and what is not.}
No theorem of this chapter is new.  The numerical verifications of
\Cref{sec:ch17-cl-checks} and \Cref{ex:ch17-multiple-fibre} are elementary,
and are included because a comparison theorem between two independently
normalized constructions is worth only as much as the smallest example in
which both sides have been computed.  \Cref{ex:ch17-multiple-fibre} is the
smallest such example in which the defect factor is not \(1\); we have not
found it written out in the literature, and we make no claim beyond that.
